\chapter{The super generalization}\label{ch:L17}
\begingroup\setlength{\emergencystretch}{8em}

\providecommand{\lseventeenvac}{\mathbf{1}}
\providecommand{\lseventeenV}{\mathcal{V}}
\providecommand{\lseventeeneps}{\varepsilon}
\providecommand{\lseventeenpar}[1]{|#1|}
\providecommand{\lseventeenGt}{\Gamma_\theta}
\providecommand{\lseventeenstr}{\operatorname{str}}

This chapter documents the release of the LeanVOA library named in the front matter; it was first written against an earlier state of the library, and checks that were run
there and not re-run for this edition are marked as such. Every Lean quotation below is taken from the files of the release (docstrings
and comments inside a quotation are omitted; the remaining text is machine-checked by
\texttt{scripts/textbook/sync.py} to occur verbatim, up to whitespace, in the source; a line
\texttt{...} marks an elision). The chapter covers the files assigned to it in
\texttt{Notes/Textbook/assignment.txt}: in \leanpath{VOA/\allowbreak Core/} the files
\leanpath{Super*.lean} and \leanpath{DongSuper*.lean}; in \leanpath{VOA/\allowbreak Examples/} the files
\leanpath{FreeFermion.lean}, \leanpath{GhostSystem.lean}, \leanpath{BCGhosts.lean},
\leanpath{BetaGamma*.lean}, \leanpath{SymplecticFermion*.lean}, \leanpath{SympFerm*.lean} and
\leanpath{LatticeSuperAmbient*.lean}; in \leanpath{VOA/\allowbreak Modules/} the files
\leanpath{Super*.lean}; and in \leanpath{VOA/\allowbreak Modularity/} the files \leanpath{Super*.lean},
\leanpath{TwistedSuperSymmetricFunctionals.lean}, \leanpath{ThetaHalfAutFactor.lean},
\leanpath{HalfIntegralGridFoundations.lean}, \leanpath{QueerOddTrace.lean} and
\leanpath{ThetaGroup.lean}. Together these are about one hundred and ten files and forty-five
thousand lines. They are not a single theory with a single headline; they are a foundation (super
vertex algebras, their Borcherds identity, their modules), four model examples (free fermion, $bc$
ghosts, $\beta\gamma$ symplectic bosons, symplectic fermions), a super Zhu theory, and a large,
explicitly incomplete programme for super and half-integral modularity. The chapter presents all of
these, at the level of the statements on which the headlines rest. (The release audit of 2026-10-27 re-elaborated the
headline statements at a later pre-release commit on branch \texttt{deligne-lean}, and re-verified the quotations below
against the source at that commit; see Section~\ref{l17:sec:status}.)

Some results of this chapter were checked in Lean by short checks written against the library;
these checks are not part of the library or of this release.

%======================================================================
\section{What this chapter is about, and what it found}\label{l17:sec:intro}
%======================================================================

Ordinary vertex operator algebras (Chapters~\ref{ch:01}, \ref{ch:L03}, \ref{ch:L12}) describe the
chiral algebra of a two-dimensional conformal field theory whose fields all \emph{commute} at
space-like separation and whose states all have \emph{integer} conformal weight. Physics needs two
generalisations at once. Fermions anticommute, so the locality axiom must carry a sign; and the
energy-momentum tensor of a theory with fermions typically has half-integer eigenvalues on some
states. The library generalises the axioms along two \emph{independent} axes:

\begin{itemize}
\item \textbf{Axis A (weights).} The conformal-weight grading is indexed by $\C$ instead of $\Z$, so
  that weights $\tfrac12$, $\tfrac32$, $\dots$ (and, for the lattice examples, arbitrary rational
  weights) are allowed. The mode index is \emph{not} changed: a field is still
  $Y(u,z)=\sum_{n\in\Z}u_{(n)}z^{-n-1}$.
\item \textbf{Axis B (parity).} A second grading by $\Z/2$ (even/odd, bosonic/fermionic) enters the
  locality axiom through the Koszul sign $(-1)^{\lseventeenpar u\lseventeenpar v}$.
\end{itemize}

The relation ``integer weight $\Leftrightarrow$ even'' (spin--statistics) is deliberately \emph{not}
an axiom; it is a separate predicate. The four model examples are chosen to separate the axes: the
$\beta\gamma$ system is half-integrally graded and purely even; the $bc$ system is integrally graded
and odd; symplectic fermions are integrally graded and odd; the free fermion is both half-integral
and odd.

\textbf{What is solidly proved.} The Koszul-signed formal calculus, Dong's lemma and Goddard
uniqueness (§\ref{l17:sec:goddard}); the full super Borcherds (Jacobi) identity in all three indices
for a super near-VOA, and from it super skew-symmetry (§\ref{l17:sec:borcherds}); the canonical
parity automorphism and the even part as an ordinary VOA (§\ref{l17:sec:aut}); the adjoint super
module; two super VOA instances: symplectic fermions ($c=-2$), the \emph{only} instance with genuinely odd states (so the
only one in which the Koszul sign rule is exercised), and $\beta\gamma$ ($c=-1$), which is half-integrally graded but
\emph{purely even} (its parity is \lean{ParityGrading.even}, so it never exercises the sign rule); and a super vertex algebra
for the $bc$ system in the library's weak sense, without translation covariance (§\ref{l17:sec:examples}); super
$C_2$-cofiniteness and a \emph{non-semisimple} super Zhu algebra of dimension \emph{at most} $4$ (and at least $1$; exact
dimension not proved) for symplectic fermions; failure of super
$C_2$-cofiniteness together with super Etingof-finiteness for $\beta\gamma$; the honest associative
super Zhu ring of every integrally graded super VOA; finite-dimensionality of twisted super Zhu
quotients under $C_2$ (§\ref{l17:sec:zhu}); the theta group $\Gamma_\theta=\langle S,T^2\rangle$ of
index $3$, a half-integral automorphy factor on it, the theta-characteristic $S$-law, the $(+,+)$
spin-sector $S$-cells, and a super Arike--Nagatomo upper bound (§\ref{l17:sec:modularity}).

\textbf{What is not proved.} The project organises super modularity as a ``four-by-three grid''
(defined precisely in §\ref{l17:sec:grid}): four \emph{columns}, one for each combination of
integral or half-integral weights with purely even or genuinely super parity, and three \emph{rows},
one for each known route to modular invariance (Zhu's rational route, the Arike--Nagatomo
$C_2$-cofinite route, and the Arakawa--Kawasetsu quasi-lisse route). A cell is \emph{complete} when
the modular-invariance theorem of that route is proved for that class and instantiated by a model.
The first column (integral, purely even) is the bosonic theory of earlier chapters; the other three
are ``new''. No modular-invariance theorem exists in any of the three new columns of the grid; the free fermion exists only as a pair of
anticommuting operator families; the $bc$ Virasoro relation is stated and unproved; ``symplectic
fermions are not super-rational'' is proved only \emph{from a hypothesis that is itself the missing
module}. Neither the free fermion nor ``symplectic fermions are not super-rational'' may be announced as a result.

\textbf{Findings of this chapter} (each is detailed in a \texttt{warning} box where it arises):
\begin{enumerate}[label=(F\arabic*)]
\item \textbf{\lean{SuperTwistedModule} is uninhabitable for $N\ge 2$} (confirmed by a Lean
  check, Theorem~\ref{l17:thm:twistvacuous}). Its \lean{jacobi} field has the same
  over-quantification as the pre-repair \lean{TwistedModule.jacobi} of Chapter~\ref{ch:06}. Hence
  every parity-twisted (Ramond) carrier, every denominator-two half-integral trace carrier, the
  sector-rationality realization and the relative Etingof layer for twisted modules are statements
  about the zero vector space. \textbf{When this chapter was first written, and still in the released source (the \lean{jacobi} field is
  unchanged), the only \lean{SuperTwistedModule} with $N\ge2$ is the zero module $M=0$} (checked in Lean against the library;
  the check is not part of this release).
  Throughout this chapter, no statement about a \lean{SuperTwistedModule} with $N\ge2$ is presented
  as having mathematical content; wherever such a statement is mentioned, it is marked ``(F1)'' and
  called vacuous.
\item \textbf{\lean{SuperVertexAlgebra} has no translation-covariance axiom}, exactly like the bosonic
  \lean{PreVertexAlgebra} of Chapter~\ref{ch:L03}: a Lean check (not part of this release) exhibits a super vertex algebra in
  the library's sense violating the super Jacobi identity (Theorem~\ref{l17:thm:nocov}). The library
  is consistent with this: \lean{superJacobi} is stated for \lean{SuperNearVOA}, which has the axiom.
\item \textbf{Orientation defect} in \lean{FreeFermion.psi}, \lean{BCGhosts.bfield},
  \lean{BCGhosts.cfield}, \lean{BetaGamma.beta}, \lean{BetaGamma.gamma}: these are indexed by mode
  number while \lean{FormalDist} is indexed by $z$-exponent, so they are not fields (not
  lower-truncated). Certificate for $\psi$ (Theorem~\ref{l17:thm:psinotfield}). Their locality
  statements remain true. The later files \leanpath{GhostSystem.lean} and
  \leanpath{BetaGammaField.lean} redo $bc$ and $\beta\gamma$ correctly; nothing redoes the free
  fermion.
\item \textbf{\lean{SympFerm.sympFermSVOA_not_isSuperRational} is conditional on
  \lean{NonSplitSuperModule sympFermSVOA}, which is never constructed} and is essentially the
  conclusion. The status docstring of \leanpath{SympFermIrrational.lean} is stale (it says the
  library has no super Jacobi identity).
\item \textbf{Super skew-symmetry exists twice}: gated (\lean{skew_modeε}, conditional on
  \lean{WeakAssocLocε}) in \leanpath{Core/\allowbreak SuperSkew.lean}, and unconditional
  (\lean{superSmul_eq_skew_Lpow}) in \leanpath{Modularity/\allowbreak SuperZhuCommutator.lean}. The
  ``gated'' remarks in \leanpath{Core/\allowbreak SuperSkew.lean},
  \leanpath{Core/\allowbreak ScreeningKernel.lean}, \leanpath{Modules/\allowbreak SuperAdjointModule.lean}
  are stale.
\item \textbf{$bc$ system}: \lean{IsGhostVirasoro} is a \lean{Prop}, not a theorem; the string value
  $c=-26$ is the arithmetic identity \lean{ghostCentralCharge_two}; the grading and the super vertex
  algebra exist only at $\lambda=1$.
\item \textbf{Twist conventions differ} between the bosonic \lean{TwistedModule} (support
  $k\equiv -j$) and \lean{SuperTwistedModule} (support $k\equiv j$) for the same eigenvalue
  $e^{-2\pi i j/N}$; harmless but undocumented.
\item \lean{IsSuperRational} quantifies over \lean{SympFermSuperOrdVModule}, a second copy of
  \lean{SuperOrdVModule} that \emph{lacks} the mode-weight axiom \lean{weight_mode}; the library
  itself flags the name as non-standard.
\end{enumerate}

\begin{warning}[title={Read first: \texttt{SuperTwistedModule} can only be the zero module for $N\ge2$}]\label{l17:warn:twistzero}
\textbf{The library structure \lean{SuperTwistedModule W g N M}, for every super VOA $W$, every even automorphism $g$ and every
twist order $N\ge2$, is inhabited only by the zero space $M=0$.} The cause is the over-quantified \lean{jacobi} field
(finding (F1)); the proof takes the \lean{jacobi} field at $u=w=\lseventeenvac$, $\ell=-1$, $k=1$, $m=-N-1$ and is checked
by Lean (Theorem~\ref{l17:thm:twistvacuous}; the check is not part of this release). The release audit of 2026-10-27 re-ran the check at a later pre-release commit and confirmed that \textbf{the
defect was still present there}; at that commit the docstrings of \leanpath{VOA/\allowbreak Modules/\allowbreak
SuperTwistedModule.lean} carried no warning. In the released source the \lean{jacobi} field is still over-quantified, and the
docstrings of the file and of the structure now flag the defect. Twelve library files use
the structure. Consequently \textbf{every claim about twisted or Ramond super modules is vacuous}: every parity-twisted
(Ramond) carrier \lean{ParityTwistedModule}, every half-integral (denominator-two) trace or supertrace carrier, the
sector-rationality realization, irreducible twisted modules (\lean{SuperTwistedModule.IsIrreducible}), the twisted zero-mode
bridge to the twisted Zhu quotient, and the relative Etingof statements for twisted super modules are statements about the zero
vector space. Throughout this chapter such statements are marked ``(F1)'' and \emph{vacuous}; none of them may be announced.
Only the identity-twist case $N=1$ carries content. Twisted \emph{Zhu quotients} $A_g(W)$ (an algebra-level construction,
Theorem~\ref{l17:thm:twistedzhufinite}) do not involve the structure and are unaffected; only their intended link to twisted
modules is.
\end{warning}

\textbf{How to read the chapter.} Section~\ref{l17:sec:conv} recalls conventions.
Sections~\ref{l17:sec:sva}--\ref{l17:sec:aut} are the foundations. Section~\ref{l17:sec:modules}
treats modules, including the defect (F1). Section~\ref{l17:sec:examples} treats the examples,
Section~\ref{l17:sec:zhu} Zhu algebras, Section~\ref{l17:sec:modularity} the modular toolkit, and
Section~\ref{l17:sec:grid} states precisely what each cell of the grid contains.

%======================================================================
\section{Conventions recalled}\label{l17:sec:conv}
%======================================================================

\textbf{Modes and $z$-exponents.} As everywhere in the library (Chapter~\ref{ch:01},
Example~\ref{ch01:ex:mode}; Chapter~\ref{ch:L01}), a formal distribution $f$ with values in $W$ is a
function $\Z\to W$, $f(k)$ being the coefficient of $z^k$; the mode is
\lean{mode f n = f (-n - 1)}. For a field $Y(u,z)=\sum_nu_{(n)}z^{-n-1}$ one has
\lean{mode (Y u) n} $=u_{(n)}$, while \lean{Y u k} $=u_{(-k-1)}$. The $n$-th product is
\lean{V.smul u n v} $=u_{(n)}v$. A distribution is a \emph{field}
(\lean{IsField}, \leanpath{Core/\allowbreak Dong.lean}) when for every vector $x$ the coefficients
$f(k)x$ vanish for $k\ll0$; that is $f(z)x\in W((z))$. We shall see (F3) that getting the
direction of this indexing wrong is the commonest mistake in the early free-field files.

\textbf{Physics modes.} A field of conformal weight $h$ is written by physicists as
$\phi(z)=\sum_r\phi_rz^{-r-h}$, so $\phi_r=\phi_{(r+h-1)}$. For $h=\tfrac12$ the physics modes are
$r\in\Z+\tfrac12$ while the vertex-algebra modes $n=r-\tfrac12$ are integers. This is why Axis A
does not touch the mode index.

\textbf{Parity, the Koszul sign and the parity involution.} The file
\leanpath{Foundations/\allowbreak Parity.lean} (Chapter~\ref{ch:L01}) defines the following.

\begin{definition}[Parity sign, Koszul sign]\label{l17:def:koszul}
For $p\in\Z/2$ let $\operatorname{parSign}(p)=(-1)^p$, and for $p,q\in\Z/2$ let
$\operatorname{koszul}(p,q)=(-1)^{pq}$, which is $-1$ exactly when $p=q=1$.
\end{definition}
\noindent File \leanpath{VOA/\allowbreak Foundations/\allowbreak Parity.lean}.
\begin{leancode}
def parSign (p : ZMod 2) : ℂ := if p = 0 then 1 else -1
...
def koszul (p q : ZMod 2) : ℂ := parSign (p * q)
\end{leancode}
\begin{dictionary}
\lean{ZMod 2} is $\Z/2$, with elements \lean{0} (even) and \lean{1} (odd); the product $p*q$ is
computed in $\Z/2$, so $\operatorname{koszul}(1,1)=\operatorname{parSign}(1)=-1$ and all other values
are $1$ (\lean{koszul_one_one}, \lean{koszul_zero_left}). The sign is a complex number, so it can
multiply operators. This is the standard Koszul rule of Kac, \emph{Vertex algebras for beginners}
§1.3, $p(a,b)=(-1)^{p(a)p(b)}$.
\end{dictionary}
\begin{physics}
Exchanging two fermionic operators costs a minus sign; exchanging a boson with anything costs
nothing. The Koszul sign is that rule, and it is the only place statistics enters the axioms.
\end{physics}
\noindent\textbf{Where the Koszul sign is used.} The Koszul sign enters every super axiom and
identity of this chapter in one way only: whenever two homogeneous objects of parities $p,q$ are
exchanged, a factor $\operatorname{koszul}(p,q)$ appears. Concretely: in the locality axiom VA4
(Definition~\ref{l17:def:sva}) it is the scalar $\varepsilon$ of $\varepsilon$-locality
(Definition~\ref{l17:def:localeps} below); in the super Jacobi identity
(Theorem~\ref{l17:thm:superjacobi}) it multiplies the term in which $v$ acts before $u$; in
skew-symmetry (Theorem~\ref{l17:thm:skew}) it is the overall sign; and in the supertrace it is the
sign of cyclicity, $\operatorname{str}(fg)=\operatorname{koszul}(p,q)\operatorname{str}(gf)$.
\begin{status}
Definitions; standard axioms only (they appear in every \texttt{\#print axioms} below).
\end{status}

A \lean{ParityGrading} on a complex vector space $M$ is a pair of subspaces $M_{\bar0},M_{\bar1}$
forming an internal direct sum (Mathlib's \lean{DirectSum.IsInternal}: every vector is uniquely a
sum of an even and an odd vector). From it the library builds, without choosing a basis, the
\emph{parity involution} $(-1)^F$ (\lean{parityInvol}, $+1$ on $M_{\bar0}$ and $-1$ on $M_{\bar1}$,
defined as $2\pi_0-\mathrm{id}$ with $\pi_0$ the projection onto $M_{\bar0}$ along $M_{\bar1}$). In a
formula: for a homogeneous vector $u\in M_{\bar p}$,
\begin{equation*}
(-1)^F u=(-1)^{p}\,u=\operatorname{parSign}(p)\,u ,
\end{equation*}
and a general vector $u=u_{\bar0}+u_{\bar1}$ is sent to $u_{\bar0}-u_{\bar1}$. The library also builds
the Koszul operator \lean{koszulOp P p} $=((-1)^F)^p$. The purely even grading
\lean{ParityGrading.even M} has $M_{\bar0}=M$ and $M_{\bar1}=0$.

\textbf{Super-locality.} Chapter~\ref{ch:L01}, Definition~\ref{l01:def:bracketeps}, defines for a
scalar $\varepsilon$ the twisted commutator and $\varepsilon$-locality.

\begin{definition}[$\varepsilon$-twisted commutator, $\varepsilon$-locality]\label{l17:def:localeps}
For operator-valued distributions $a,b$ and $\varepsilon\in\C$,
$[a,b]_\varepsilon(z_1,z_2)=a(z_1)b(z_2)-\varepsilon\,b(z_2)a(z_1)$ (coefficientwise), and $a,b$ are
$\varepsilon$-local if $(z_1-z_2)^N[a,b]_\varepsilon=0$ for some $N\in\N$.
\end{definition}
\noindent File \leanpath{VOA/\allowbreak Foundations/\allowbreak Locality.lean}.
\begin{leancode}
def bracketDistε (ε : ℂ) (a b : FormalDist (Module.End ℂ V)) : FormalDist₂ (Module.End ℂ V) :=
  fun m n => a m * b n - ε • (b n * a m)
...
def Localε (ε : ℂ) (a b : FormalDist (Module.End ℂ V)) : Prop :=
  ∃ N : ℕ, mulBySubPow N (bracketDistε ε a b) = 0
\end{leancode}
\begin{dictionary}
\lean{a m * b n} is composition of the $z_1^m$-coefficient of $a$ with the $z_2^n$-coefficient of $b$
(exponent indexing). \lean{mulBySubPow N} is multiplication by $(z_1-z_2)^N$. The sign $\varepsilon$ is
a bare scalar: $\varepsilon=1$ is ordinary locality (\lean{localε_one_iff}) and the super case uses
$\varepsilon=\operatorname{koszul}(p,q)$. Chapter~\ref{ch:L01}, Theorem~\ref{l01:thm:car}, proves the
free-field criterion: if $[a_m,b_n]_\varepsilon=\delta_{m+n+1,0}S$ then $a,b$ are $\varepsilon$-local
with $N=1$ (\lean{localε_of_car}, \lean{localε_neg_one_of_anticomm}). Note that this criterion is
symmetric under reindexing $m\mapsto-m-1$ in both slots, which is why locality cannot detect an
orientation error (F3).
\end{dictionary}
\begin{physics}
$(z-w)^N[a(z),b(w)]_\varepsilon=0$ says that the (anti)commutator is supported at $z=w$: fields
(anti)commute at space-like separation, and their operator product has a pole of order at most $N$.
For two fermions $\psi(z)\psi(w)\sim 1/(z-w)$ the anticommutator is a delta function and $N=1$.
\end{physics}
\begin{status}
Definitions (Chapter~\ref{ch:L01}).
\end{status}

%======================================================================
\section{Super vertex algebras}\label{l17:sec:sva}
%======================================================================

\subsection{The axioms}

The library's definition of a super vertex algebra was quoted verbatim, field by field, in
Chapter~\ref{ch:01}, §\ref{ch01:sec:super}, from \leanpath{VOA/\allowbreak Core/\allowbreak
SuperVertexAlgebra.lean}. We restate it mathematically and compare it with the literature.

\begin{definition}[Super vertex algebra, library sense]\label{l17:def:sva}
A \emph{super vertex algebra} (\lean{SuperVertexAlgebra 𝒱}) on a complex vector space $\lseventeenV$
consists of
\begin{itemize}
\item a grading $\lseventeenV=\bigoplus_{h\in\C}\lseventeenV_h$ (internal direct sum) with every
  $\lseventeenV_h$ finite-dimensional (\lean{grading}, \lean{isInternal}, \lean{finrank});
\item a second grading $\lseventeenV=\lseventeenV_{\bar0}\oplus\lseventeenV_{\bar1}$ (\lean{parity},
  \lean{parity_isInternal});
\item a linear map $Y:\lseventeenV\to\End(\lseventeenV)[[z^{\pm1}]]$ and a vector $\lseventeenvac$;
\end{itemize}
such that: $\lseventeenvac\in\lseventeenV_0\cap\lseventeenV_{\bar0}$; (VA1) $u\in\lseventeenV_h$,
$v\in\lseventeenV_k$ imply $u_{(n)}v\in\lseventeenV_{h+k-n-1}$; (VA1$'$) $u\in\lseventeenV_{\bar p}$,
$v\in\lseventeenV_{\bar q}$ imply $u_{(n)}v\in\lseventeenV_{\overline{p+q}}$; (VA2)
$Y(\lseventeenvac,z)=\mathrm{id}$; (VA3) $u_{(-1)}\lseventeenvac=u$ and $u_{(n)}\lseventeenvac=0$ for
$n\ge0$; (VA4) for parity-homogeneous $u,v$ of parities $p,q$, $Y(u,z)$ and $Y(v,z)$ are
$\operatorname{koszul}(p,q)$-local; (VA5) for all $u,v$ there is $N$ with $u_{(n)}v=0$ for $n\ge N$.
There is \emph{no} translation operator and \emph{no} translation-covariance axiom; because of this the
super Borcherds identity can fail for such an object (finding (F2), Theorem~\ref{l17:thm:nocov}), so
this is weaker than (and incomparable with) a vertex superalgebra in Kac's sense. Translation
covariance is added in Definition~\ref{l17:def:snvoa}.
\end{definition}
\begin{dictionary}
\textbf{Decoding.} Each clause of the definition is one field of the structure quoted in
§\ref{ch01:sec:super}; \lean{finrank} is a proposition (finite-dimensionality), not a number. Locality
is demanded only for homogeneous pairs, which is the standard formulation (it cannot be demanded for
inhomogeneous pairs, since the sign would not be defined).

\textbf{Comparison with Kac.} Kac (\emph{Vertex algebras for beginners}, Definition~1.3 and §4.1)
defines a vertex superalgebra as $(V,\lseventeenvac,T,Y)$ with $V=V_{\bar0}\oplus V_{\bar1}$, $Y(a,z)$ a
field of parity $p(a)$, $T\lseventeenvac=0$, $[T,Y(a,z)]=\partial_zY(a,z)$ (translation covariance),
$Y(\lseventeenvac,z)=\mathrm{id}$, $Y(a,z)\lseventeenvac|_{z=0}=a$, and super-locality. The library's
definition differs in three ways. (i) It \emph{adds} a $\C$-grading with finite-dimensional pieces
compatible with the modes (VA1); Kac has no grading in a vertex superalgebra. (ii) It \emph{adds}
VA1$'$ (``$Y(a,z)$ has parity $p(a)$''), which in Kac is part of the data. (iii) It \emph{omits} the
translation operator and translation covariance. Exactly as for the bosonic \lean{PreVertexAlgebra}
(Chapter~\ref{ch:L03}, §\ref{l03:sec:counter}), omission (iii) is not harmless: the Borcherds identity
does not follow (Theorem~\ref{l17:thm:nocov}). So the library's \lean{SuperVertexAlgebra} is
\emph{incomparable} with Kac's vertex superalgebra. With the conformal layer added
(\lean{SuperNearVOA}, Definition~\ref{l17:def:snvoa}), translation covariance comes back through
VOA3, and the notion becomes Kac's vertex superalgebra with a compatible $\C$-grading and conformal
vector.

\textbf{No spin--statistics.} Nothing relates the two gradings. This is intended
(Definition~\ref{l17:def:spinstat}).

\textbf{How it could have been wrong.} An unsatisfiable axiom system is excluded by the instances of
§\ref{l17:sec:examples} (symplectic fermions, $\beta\gamma$, $bc$ are genuine inhabitants with
nonzero odd or half-integral parts) and by \lean{PreVertexAlgebra.toSuper} below.
\end{dictionary}
\begin{physics}
States $\leftrightarrow$ fields: $Y(u,z)$ is the local field creating the state $u$ from the vacuum.
Odd states create fermionic fields; the super-locality axiom says their operator products are
radially ordered with the fermionic sign. The $\C$-grading is the spectrum of the dilatation operator
$L_0$ (once a conformal vector is present); allowing $h\in\frac12+\Z$ is needed for any theory with
free fermions in the Neveu--Schwarz sector.
\end{physics}
\begin{status}
A definition. \texttt{\#print axioms VOA.SuperVertexAlgebra}: \lean{[propext, Classical.choice,
Quot.sound]} (run for this chapter).
\end{status}

\begin{theorem}[No translation covariance]\label{l17:thm:nocov}
There is a three-dimensional super vertex algebra in the library's sense (purely even, weights $0$
and $2$) on which the super Jacobi identity fails at $(u,v,x,\ell,m,n)=(x_1,\lseventeenvac,\lseventeenvac,-2,0,-2)$
and on which $Y(Tu,z)\neq\partial_zY(u,z)$.
\end{theorem}
\begin{dictionary}
The example is the counterexample of Chapter~\ref{ch:L03}, Theorem~\ref{l03:thm:counter}:
$V=\operatorname{span}\{\lseventeenvac,x_1,y\}$, $Y(u,z)$ = multiplication by
$\varphi_z(u)$ with $\varphi_z(x_1)=x_1+z^2y$, read as a super vertex algebra
(\lean{PreVertexAlgebra.toSuper}, Proposition~\ref{l17:prop:tosuper}). The failing
identity is \emph{exactly} the conclusion of \lean{superJacobi} (Theorem~\ref{l17:thm:superjacobi})
at parities $p=q=0$.
\end{dictionary}
\begin{physics}
A ``vertex algebra'' without a momentum operator is not a chiral algebra: nothing ties the
$z$-dependence of a field to the action of translations on states, and operator product
associativity fails. Physically every chiral algebra has $L_{-1}$; the library adds it in
\lean{SuperNearVOA}.
\end{physics}
\begin{proof}[Proof on paper]
Since $\varphi_z(x_1)\lseventeenvac$ has no $z^1$ term, $Tx_1=x_{1(-2)}\lseventeenvac=0$, but
$\partial_zY(x_1,z)=2zy\neq0$. In the Jacobi identity at the stated indices the right side is
$(x_{1(-2)}\lseventeenvac)_{(-2)}\lseventeenvac=0$ while the left side is
$2\,x_{1(-3)}\lseventeenvac=2y\neq0$ (Chapter~\ref{ch:L03}).
\end{proof}
\begin{status}
Checked in Lean against the library (the check is not part of this release); not a library theorem.
\texttt{\#print axioms} of both statements: \lean{[propext, Classical.choice, Quot.sound]}.
\end{status}

\begin{warning}[title={Finding (F2): the bare super vertex algebra has the L03 defect}]
\lean{SuperVertexAlgebra} omits translation covariance, so a theorem about a bare
\lean{SuperVertexAlgebra} cannot use the Borcherds identity. The library is consistent: Goddard
uniqueness and Dong's lemma (which do not need it) are stated for \lean{SuperVertexAlgebra}; the
Borcherds/Jacobi identity, skew-symmetry and everything downstream are stated for \lean{SuperNearVOA}
or \lean{SuperVertexOperatorAlgebra}. The danger is only in the name. The bosonic structure with the same omission was
renamed \lean{VOA.PreVertexAlgebra} for exactly this reason, the name \lean{VertexAlgebra} now denoting the
translation-covariant structure (Chapter~\ref{ch:L03}); the super structure still carries the name
\lean{SuperVertexAlgebra} in the released source. Severity: moderate (the remedy is the one applied on the bosonic side).
\end{warning}

\begin{definition}[Spin--statistics]\label{l17:def:spinstat}
A super vertex algebra satisfies \emph{spin--statistics} if every nonzero homogeneous vector $u$ of
weight $h$ and parity $p$ satisfies: $h\in\Z$ if and only if $p=\bar0$.
\end{definition}
\noindent File \leanpath{VOA/\allowbreak Core/\allowbreak SuperVertexAlgebra.lean}.
\begin{leancode}
def SpinStatistics {𝒱 : Type u} [AddCommGroup 𝒱] [Module ℂ 𝒱]
    (W : SuperVertexAlgebra 𝒱) : Prop :=
  ∀ (h : ℂ) (p : ZMod 2) (u : 𝒱), u ≠ 0 → u ∈ W.grading h → u ∈ W.parity p →
    ((∃ m : ℤ, (m : ℂ) = h) ↔ p = 0)
\end{leancode}
\begin{dictionary}
``$h\in\Z$'' is written $\exists m:\Z,\ (m:\C)=h$. The quantification is over homogeneous nonzero $u$
(an inhomogeneous vector has no single weight). It is a predicate, never a field; theorems needing
it must say so. Kac's §4.1 notion ``$V_{\bar0}=\bigoplus_{n\in\Z}V_n$, $V_{\bar1}=\bigoplus_{n\in\frac12+\Z}V_n$''
is this predicate together with half-integrality of all weights.
\end{dictionary}
\begin{physics}
The spin--statistics theorem of relativistic quantum field theory: integer spin $\Leftrightarrow$
Bose statistics. In a chiral algebra the spin of a field is its conformal weight. Ghost systems
($bc$, $\beta\gamma$) and symplectic fermions violate it, because they are not unitary theories; the
library must admit them, hence the predicate.
\end{physics}
\begin{status}
A definition; standard axioms. It is satisfiable (Theorem~\ref{l17:thm:tosuperspin}) and fails for
$\beta\gamma$ (\lean{bgSVA_not_spinStatistics}, §\ref{l17:sec:bg}). For symplectic fermions its
failure is noted in a docstring but not formalised.
\end{status}

\begin{proposition}[Every vertex algebra is a super vertex algebra]\label{l17:prop:tosuper}
Every ordinary vertex algebra $V$ (Chapter~\ref{ch:L03}, Definition~\ref{l03:def:VA}) defines a super
vertex algebra with the same $Y$ and vacuum, the $\Z$-grading pushed forward along $\Z\hookrightarrow\C$
(zero off the integers), and purely even parity.
\end{proposition}
\noindent File \leanpath{VOA/\allowbreak Core/\allowbreak SuperVertexAlgebra.lean}.
\begin{leancode}
noncomputable def toSuper (V : PreVertexAlgebra 𝒱) : SuperVertexAlgebra 𝒱 where
  grading := V.gradingC
...
  parity := (ParityGrading.even 𝒱).parity
  parity_isInternal := (ParityGrading.even 𝒱).isInternal
  Y := V.Y
  vacuum := V.vacuum
\end{leancode}
\begin{dictionary}
\lean{V.gradingC} is \lean{extendGrading (fun n : ℤ => (n : ℂ)) V.grading}: equal to $V_n$ at
$h=n\in\Z$ and $0$ elsewhere. The remaining fields are proved: locality at Koszul sign
$\operatorname{koszul}(0,0)=1$ is ordinary locality; pairs with an odd member are zero. The
analogous \lean{VertexOperatorAlgebra.toSuper} exists in \leanpath{Core/\allowbreak
SuperVertexOperatorAlgebra.lean}. So the super theory is a generalisation, not a parallel hierarchy.
\end{dictionary}
\begin{physics}
A bosonic chiral algebra is a fermionic one with no fermions. This embedding is what lets the
super modularity statements specialise to the ordinary ones (the ``integral non-super'' column of
§\ref{l17:sec:grid}).
\end{physics}
\begin{proof}[Proof on paper]
Internality of the pushed-forward grading is \lean{isInternal_extend} (an internal decomposition
reindexed along an injection, with zero pieces added, is still internal). Every other axiom is the
bosonic one, the Koszul signs being $1$.
\end{proof}
\begin{status}
\texttt{\#print axioms}: standard (checked for the dependent theorem below).
\end{status}

\begin{theorem}[Spin--statistics is satisfiable]\label{l17:thm:tosuperspin}
For every ordinary VOA $W$, the super VOA \lean{W.toSuper} satisfies spin--statistics.
\end{theorem}
\noindent File \leanpath{VOA/\allowbreak Core/\allowbreak SuperVertexOperatorAlgebra.lean}.
\begin{leancode}
theorem toSuper_spinStatistics (W : VertexOperatorAlgebra 𝒱) :
    SpinStatistics W.toSuper.toSuperVertexAlgebra := by
\end{leancode}
\begin{dictionary}
Unconditional. Its role is a non-vacuity check: the predicate is not malformed.
\end{dictionary}
\begin{physics}
A purely bosonic theory with integer weights trivially obeys spin--statistics.
\end{physics}
\begin{proof}[Proof on paper]
Every nonzero vector is even and has integer weight, so both sides of the equivalence hold.
\end{proof}
\begin{status}
\texttt{\#print axioms VOA.VertexOperatorAlgebra.toSuper\_spinStatistics}: \lean{[propext,
Classical.choice, Quot.sound]}.
\end{status}

\subsection{The conformal layer}

\begin{definition}[Super near-VOA]\label{l17:def:snvoa}
A \emph{super near-VOA} is a super vertex algebra with an even vector $\omega\in\lseventeenV_2$ and
$c\in\C$ such that $L_n:=\omega_{(n+1)}$ satisfy the Virasoro relations with central charge $c$
(VOA1), $L_0v=hv$ for $v\in\lseventeenV_h$ (VOA2), and $Y(L_{-1}v,z)=\partial_zY(v,z)$ for all $v$
(VOA3).
\end{definition}
\noindent File \leanpath{VOA/\allowbreak Core/\allowbreak SuperNearVOA.lean}.
\begin{leancode}
structure SuperNearVOA (𝒱 : Type u) [AddCommGroup 𝒱] [Module ℂ 𝒱]
    extends SuperVertexAlgebra 𝒱 where
  conformalVector : 𝒱
  conformal_mem : conformalVector ∈ grading 2
  conformal_parity : conformalVector ∈ parity 0
  centralCharge : ℂ
  virasoro : ∀ m n : ℤ,
      ⁅mode (Y conformalVector) (m + 1), mode (Y conformalVector) (n + 1)⁆
        = ((m : ℂ) - n) • mode (Y conformalVector) (m + n + 1)
          + (if m + n = 0 then (((m : ℂ) ^ 3 - m) / 12 * centralCharge) • (1 : End ℂ 𝒱)
             else 0)
  L0_grading : ∀ {h : ℂ} (v : 𝒱), v ∈ grading h → mode (Y conformalVector) 1 v = h • v
  translation : ∀ v : 𝒱, Y (mode (Y conformalVector) 0 v) = deriv (Y v)
\end{leancode}
\begin{dictionary}
\lean{mode (Y conformalVector) (m + 1)} is $\omega_{(m+1)}=L_m$; \lean{⁅·,·⁆} is the commutator of
endomorphisms; VOA1 reads $[L_m,L_n]=(m-n)L_{m+n}+\delta_{m+n,0}\frac{m^3-m}{12}c$. VOA2 has no cast
because the grading index is the eigenvalue itself. VOA3 uses \lean{mode (Y ω) 0} $=L_{-1}$, and
\lean{deriv} is coefficientwise $\partial_z$. The evenness of $\omega$ is an axiom (it does not follow,
since spin--statistics is not assumed). Missing, by design: $\lseventeenV_0=\C\lseventeenvac$ and
``no weights of negative real part'' (the CFT-type conditions), because ambient lattice algebras with
background charge (§\ref{l17:sec:lattice}) violate positivity. The weaker \lean{BoundedBelow}
($\exists B\in\mathbb{R}$, $\lseventeenV_h=0$ whenever $\operatorname{Re}h<B$) is a separate predicate. This is
the super analogue of the bosonic \lean{NearVOA} of Chapter~\ref{ch:L03}.
\end{dictionary}
\begin{physics}
$\omega$ is the stress tensor $T(z)=\sum L_nz^{-n-2}$; $c$ the central charge (conformal anomaly);
$L_0$ the energy; $L_{-1}$ the momentum $\partial_z$. The stress tensor is a bosonic field even in a
fermionic theory, hence even.
\end{physics}
\begin{status}
A definition; standard axioms.
\end{status}

\begin{definition}[Super VOA]\label{l17:def:svoa}
A \emph{super vertex operator algebra} (\lean{SuperVertexOperatorAlgebra}) has the same data and
axioms as a super near-VOA, plus (VOA4) $\lseventeenV_0=\C\lseventeenvac$ and $\lseventeenV_h=0$ whenever
$\operatorname{Re}h<0$.
\end{definition}
\noindent File \leanpath{VOA/\allowbreak Core/\allowbreak SuperVertexOperatorAlgebra.lean}.
\begin{leancode}
structure SuperVertexOperatorAlgebra (𝒱 : Type u) [AddCommGroup 𝒱] [Module ℂ 𝒱]
    extends SuperVertexAlgebra 𝒱 where
...
  vacuum_grade : grading 0 = ℂ ∙ vacuum
  positivity : ∀ h : ℂ, h.re < 0 → grading h = ⊥
\end{leancode}
\begin{dictionary}
The elided fields are verbatim those of \lean{SuperNearVOA} (\lean{conformalVector} through
\lean{translation}). \lean{ℂ ∙ vacuum} is the line spanned by $\lseventeenvac$; \lean{⊥} the zero
subspace. Conversions: \lean{SuperVertexOperatorAlgebra.toSuperNearVOA} (forget VOA4) and
\lean{SuperNearVOA.toSVOA W hvac hpos} (supply VOA4). Compared with Kac §4.1 and with Dong--Zhao's
``vertex operator superalgebra'', the library allows complex weights (they require
$\frac12\Z$) and does not require $\dim\lseventeenV_h<\infty$ to come from a $\frac12\Z_{\ge0}$-grading;
these are generalisations, and every theorem below states the integrality it needs:
\lean{IntegralWeights} says $\lseventeenV_h=0$ unless $h\in\Z$, and \lean{IsHalfIntegerGraded} says
$\lseventeenV_h=0$ unless $h\in\frac12\Z$ (both are predicates on the weight grading, defined in
\leanpath{Core/\allowbreak SuperEvenPart.lean} and \leanpath{Modularity/\allowbreak
GammaThetaStatement.lean} respectively, and recalled where they are used).
\end{dictionary}
\begin{physics}
CFT type: a unique vacuum and energy bounded below by zero, as in a unitary theory. Symplectic
fermions and $\beta\gamma$ at $\lambda=\frac12$ are CFT type though non-unitary.
\end{physics}
\begin{status}
A definition; standard axioms.
\end{status}

\begin{theorem}[Dong's lemma for super vertex algebras]\label{l17:thm:dong}
If $u,v,w$ are homogeneous of parities $p,q,r$, then for every $k\ge0$ the field
$Y(u,z)_{(k)}Y(v,z)$ (the $k$-th OPE coefficient of the $\operatorname{koszul}(p,q)$-twisted
commutator) is $\operatorname{koszul}(p+q,r)$-local with $Y(w,z)$.
\end{theorem}
\noindent File \leanpath{VOA/\allowbreak Core/\allowbreak SuperVertexAlgebra.lean}.
\begin{leancode}
theorem dong_fields (V : SuperVertexAlgebra 𝒱) {p q r : ZMod 2} {u v w : 𝒱}
    (hu : u ∈ V.parity p) (hv : v ∈ V.parity q) (hw : w ∈ V.parity r) (k : ℕ) :
    Localε (koszul (p + q) r)
      (cCoeff k (bracketDistε (koszul p q) (V.Y u) (V.Y v))) (V.Y w)
\end{leancode}
\begin{dictionary}
\lean{cCoeff k F} is $\Res_{z_1}(z_1-z_2)^kF(z_1,z_2)$, the $k$-th OPE coefficient (Chapter~\ref{ch:L03},
§\ref{l03:sec:conventions}); for $k\ge0$ it is the field $Y(u_{(k)}v,z)$ once Borcherds is known.
The sign of the product field is $\operatorname{koszul}(p+q,r)$, exactly the Koszul sign of a state
of parity $p+q$, as VA1$'$ predicts. The general $\varepsilon$-generic Dong lemma
(\lean{dong_nonnegε}, \leanpath{Core/\allowbreak DongSuper.lean}) needs only
$\varepsilon_{ab}^2=1$, which for Koszul signs is \lean{koszul_mul_self}; negative products are in
\leanpath{Core/\allowbreak DongSuperNeg.lean} (\lean{dong_negε}). Kac, Lemma 3.2 (super version).
\end{dictionary}
\begin{physics}
The OPE of two local fields produces local fields, with the statistics of the product: two
fermions make a boson. This is what makes the operator algebra close.
\end{physics}
\begin{proof}[Proof on paper]
The bosonic proof (Chapter~\ref{ch:L03}, Theorem~\ref{l03:thm:dong}) multiplies the three-variable
Jacobi identity of commutators by a large power of $(z_1-z_3)$ and $(z_2-z_3)$. With signs, the
twisted Jacobi identity $[[a,b]_{ab},c]_{ac\,bc}=[a,[b,c]_{bc}]_{ab\,ac}-\varepsilon_{ab}[b,[a,c]_{ac}]$
holds provided $\varepsilon_{ab}^2=1$ (the residual otherwise is $(1-\varepsilon_{ab}^2)$ times a
nonzero expression), and the sign of the composite is
$\varepsilon_{(ab)c}=\varepsilon_{ac}\varepsilon_{bc}$, i.e.\ \lean{koszul_add_left}.
\end{proof}
\begin{status}
\texttt{\#print axioms VOA.SuperVertexAlgebra.dong\_fields}: \lean{[propext, Classical.choice,
Quot.sound]}. Unconditional.
\end{status}

%======================================================================
\section{Goddard uniqueness with signs}\label{l17:sec:goddard}
%======================================================================

\begin{theorem}[Goddard uniqueness, $\varepsilon$-generic core]\label{l17:thm:goddardpt}
Let $A$ be an operator-valued distribution, $a,b\in\lseventeenV$ and $\varepsilon\in\C$ arbitrary. If
$A$ and $Y(a,z)$ are both $\varepsilon$-local with $Y(b,z)$, and $A(z)\lseventeenvac=Y(a,z)\lseventeenvac$
coefficientwise, then $A(z)b=Y(a,z)b$.
\end{theorem}
\noindent File \leanpath{VOA/\allowbreak Core/\allowbreak SuperGoddard.lean}.
\begin{leancode}
theorem goddard_pointwiseε (V : SuperVertexAlgebra 𝒱) (A : FormalDist (End ℂ 𝒱)) (a : 𝒱)
    (ε : ℂ) (b : 𝒱) (hAloc : Localε ε A (V.Y b)) (hVloc : Localε ε (V.Y a) (V.Y b))
    (hAvac : ∀ m : ℤ, A m V.vacuum = (V.Y a) m V.vacuum) (m : ℤ) :
    A m b = (V.Y a) m b
\end{leancode}
\begin{dictionary}
\lean{A m} is the $z^m$-coefficient. The conclusion is for every $m$, i.e.\ $A(z)b=Y(a,z)b$. No
parity hypothesis: $\varepsilon$ is any scalar, the \emph{same} for both locality hypotheses. No
translation operator is used, so the statement is correctly placed at the bare level.
\end{dictionary}
\begin{physics}
A local field is determined by the state it creates: there is one field per state. This is the
uniqueness half of the state--field correspondence, and it does not care about statistics.
\end{physics}
\begin{proof}[Proof on paper]
Let $D(z,w)=A(z)Y(b,w)\lseventeenvac-Y(a,z)Y(b,w)\lseventeenvac$. Using the creation hypothesis,
$D$ is the difference of the two twisted brackets applied to $\lseventeenvac$ (the $\varepsilon$-terms
$\varepsilon Y(b,w)A(z)\lseventeenvac$ and $\varepsilon Y(b,w)Y(a,z)\lseventeenvac$ coincide and
cancel). Hence $(z-w)^ND=0$ for large $N$. But $D$ has no negative powers of $w$ (VA3), and in
$\lseventeenV((z))[[w]]$ the element $z-w$ is not a zero divisor, so $D=0$; setting $w=0$ gives
$A(z)b=Y(a,z)b$.
\end{proof}
\begin{status}
\texttt{\#print axioms VOA.SuperVertexAlgebra.goddard\_pointwiseε}: \lean{[propext, Classical.choice,
Quot.sound]}. Unconditional.
\end{status}

\begin{theorem}[Goddard uniqueness for super vertex algebras]\label{l17:thm:goddard}
Let $a\in\lseventeenV_{\bar p}$ and let $A$ be a distribution that is $\operatorname{koszul}(p,q)$-local
with $Y(b,z)$ for every $q$ and every $b\in\lseventeenV_{\bar q}$, and creates the same vacuum tower as
$Y(a,z)$. Then $A=Y(a,z)$.
\end{theorem}
\noindent File \leanpath{VOA/\allowbreak Core/\allowbreak SuperGoddard.lean}.
\begin{leancode}
theorem goddard_uniquenessε (V : SuperVertexAlgebra 𝒱) (A : FormalDist (End ℂ 𝒱)) (a : 𝒱)
    {p : ZMod 2} (_ha : a ∈ V.parity p)
    (hAloc : ∀ (q : ZMod 2) (b : 𝒱), b ∈ V.parity q → Localε (koszul p q) A (V.Y b))
    (hAvac : ∀ m : ℤ, A m V.vacuum = (V.Y a) m V.vacuum) :
    A = V.Y a
\end{leancode}
\begin{dictionary}
The hypothesis \lean{_ha} is used (the underscore is only a naming style): it supplies
\lean{V.locality a _ha b₀ hb₀}. The conclusion is equality of distributions. Inhomogeneous $b$ are
handled by splitting $b=b_0+b_1$; the conclusion is linear in $b$ although locality is not. Matches
Kac, Theorem 4.4 (super case) and Frenkel--Ben-Zvi §3.1.
\end{dictionary}
\begin{physics}
Same as above: uniqueness of the field attached to a state, for bosons and fermions alike.
\end{physics}
\begin{proof}[Proof on paper]
Apply Theorem~\ref{l17:thm:goddardpt} to the even and odd parts of $b$, with
$\varepsilon=\operatorname{koszul}(p,0)$ and $\operatorname{koszul}(p,1)$ respectively, using VA4 for
$Y(a,z)$.
\end{proof}
\begin{status}
\texttt{\#print axioms VOA.SuperVertexAlgebra.goddard\_uniquenessε}: \lean{[propext,
Classical.choice, Quot.sound]}. Unconditional.
\end{status}

The super reconstruction theorem \lean{reconstructε} (\leanpath{Core/\allowbreak
SuperReconstruction.lean}) is the existence counterpart: from a spanning family of states with
$T$-covariant, mutually Koszul-local generating fields and the vacuum conditions, it produces a
linear $Y$ with the creation, locality and covariance properties. It is the super analogue of the
operator-first reconstruction of Chapter~\ref{ch:L03}, Theorem~\ref{l03:thm:reconstruct}, and it is
how every example of §\ref{l17:sec:examples} obtains its vertex operator. Its proof uses
Theorem~\ref{l17:thm:goddardpt} and the super Dong lemma; translation covariance of normal-ordered
products is \lean{DCovariant.nprodε} in \leanpath{Core/\allowbreak SuperCovariant.lean}, whose key
observation is that the twisted normal-ordered product differs from the bosonic one by a scalar
depending only on the first mode index.

%======================================================================
\section{The super Borcherds identity and skew-symmetry}\label{l17:sec:borcherds}
%======================================================================

\begin{definition}[The two sides of the super Borcherds identity]\label{l17:def:BLBR}
For a super near-VOA $V$, vectors $u,v,x$, a scalar $\varepsilon$ and $\ell,m,n\in\Z$ put
\begin{align*}
B_R(u,v,x;\ell,m,n)&=\sum_{i\ge0}\binom{m}{i}\,(u_{(\ell+i)}v)_{(m+n-i)}x,\\
B_L^\varepsilon(u,v,x;\ell,m,n)&=\sum_{i\ge0}(-1)^i\binom{\ell}{i}\Bigl(u_{(\ell+m-i)}v_{(n+i)}x
  -\varepsilon(-1)^\ell\,v_{(\ell+n-i)}u_{(m+i)}x\Bigr).
\end{align*}
\end{definition}
\noindent File \leanpath{VOA/\allowbreak Core/\allowbreak SuperBorcherds.lean}.
\begin{leancode}
noncomputable def BRs (V : SuperNearVOA 𝒱) (u v x : 𝒱) (ℓ m n : ℤ) : 𝒱 :=
  ∑ᶠ i : ℤ, cbz m i • V.smul (V.smul u (ℓ + i) v) (m + n - i) x
noncomputable def Puvs (V : SuperNearVOA 𝒱) (u v x : 𝒱) (ℓ m n : ℤ) : 𝒱 :=
  ∑ᶠ i : ℤ, (negOnePow i * cbz ℓ i) • V.smul u (ℓ + m - i) (V.smul v (n + i) x)
noncomputable def Pvus (V : SuperNearVOA 𝒱) (u v x : 𝒱) (ℓ m n : ℤ) : 𝒱 :=
  ∑ᶠ i : ℤ, (negOnePow i * cbz ℓ i) • V.smul v (ℓ + n - i) (V.smul u (m + i) x)
noncomputable def BLε (V : SuperNearVOA 𝒱) (ε : ℂ) (u v x : 𝒱) (ℓ m n : ℤ) : 𝒱 :=
  Puvs V u v x ℓ m n - (ε * negOnePow ℓ) • Pvus V u v x ℓ m n
\end{leancode}
\begin{dictionary}
\lean{cbz m i} is the generalised binomial coefficient $\binom mi$ for $m\in\Z$, $i\in\Z$, equal to
$0$ for $i<0$ (Chapter~\ref{ch:L03}, Definition~\ref{l03:def:BLBR}); so the $\Z$-indexed
\lean{finsum} $\sum^{\mathrm f}$ is really over $i\ge0$, and it has finite support by VA5 even when
$m<0$. \lean{negOnePow i} is $(-1)^i$. The exchanged term carries the single scalar
$\varepsilon(-1)^\ell$; all sign-free parts are literally the bosonic ones of Chapter~\ref{ch:L03}.
\end{dictionary}
\begin{physics}
$B_L$ is the contour-deformation side of the OPE (the difference of the two radial orderings of
$u(z)v(w)$, with the fermionic sign on the reversed ordering); $B_R$ is the OPE side (the fields
appearing in $u\times v$). Their equality is operator-product associativity.
\end{physics}
\begin{status}
Definitions; standard axioms.
\end{status}

\begin{theorem}[Super Borcherds / Jacobi identity]\label{l17:thm:superjacobi}
Let $V$ be a super near-VOA, $u\in V_{\bar p}$, $v\in V_{\bar q}$, $x\in V$. For all $\ell,m,n\in\Z$,
\begin{equation*}
\sum_{i\ge0}(-1)^i\binom{\ell}{i}\Bigl(u_{(\ell+m-i)}v_{(n+i)}x
-(-1)^{pq}(-1)^\ell\,v_{(\ell+n-i)}u_{(m+i)}x\Bigr)
=\sum_{i\ge0}\binom mi\,(u_{(\ell+i)}v)_{(m+n-i)}x .
\end{equation*}
\end{theorem}
\noindent File \leanpath{VOA/\allowbreak Core/\allowbreak SuperBorcherds.lean}.
\begin{leancode}
theorem superBorcherds (V : SuperNearVOA 𝒱) {p q : ZMod 2} {u v : 𝒱} (hu : u ∈ V.parity p)
    (hv : v ∈ V.parity q) (x : 𝒱) (ℓ m n : ℤ) :
    BLε V (koszul p q) u v x ℓ m n = BRs V u v x ℓ m n := by
...
theorem superJacobi (V : SuperNearVOA 𝒱) {p q : ZMod 2} {u v : 𝒱} (hu : u ∈ V.parity p)
    (hv : v ∈ V.parity q) (x : 𝒱) (ℓ m n : ℤ) :
    (∑ᶠ i : ℤ, (negOnePow i * cbz ℓ i) •
        (V.smul u (ℓ + m - i) (V.smul v (n + i) x)
          - (koszul p q * negOnePow ℓ) • V.smul v (ℓ + n - i) (V.smul u (m + i) x)))
      = ∑ᶠ i : ℤ, cbz m i • V.smul (V.smul u (ℓ + i) v) (m + n - i) x := by
\end{leancode}
\begin{dictionary}
\textbf{Index dictionary with Kac.} Kac's (4.8.3) (super case, \emph{Vertex algebras for
beginners}) reads
$\sum_j\binom mj(a_{(n+j)}b)_{(m+k-j)}c=\sum_j(-1)^j\binom nj\bigl(a_{(m+n-j)}b_{(k+j)}c-p(a,b)(-1)^nb_{(n+k-j)}a_{(m+j)}c\bigr)$.
Renaming Kac's $(n,m,k)$ as the library's $(\ell,m,n)$ gives the library statement character for
character, including the position of the sign. The quantifiers are ``for all homogeneous $u,v$, all
$x$, all $\ell,m,n\in\Z$'': the full identity, not a slice. Inhomogeneous $u,v$ are not covered,
since the sign is undefined; by linearity nothing is lost.

\textbf{Hypotheses.} Only the structure \lean{SuperNearVOA} (no CFT type, no spin--statistics, no
integrality of weights) and parity homogeneity. The second form is the first unfolded
(\lean{BL_as_finsumε}).

\textbf{Why it is not vacuous.} It is consumed definitionally as the \lean{jacobi} field of the adjoint
super module (Proposition~\ref{l17:prop:adjoint}), and it is the input to the super Zhu ring
(Theorem~\ref{l17:thm:superzhuring}); at $p=0$ it reduces to the bosonic Jacobi identity
(\lean{superJacobi_even_left}).
\end{dictionary}
\begin{physics}
This is the full OPE associativity of a fermionic chiral algebra: for every choice of three contour
exponents, the reordered product of $u(z)$ and $v(w)$, with the minus sign for two fermions, equals
the field of the OPE coefficients. All the usual consequences (the commutator formula
$[u_{(m)},v_{(n)}]_\pm=\sum\binom mi(u_{(i)}v)_{(m+n-i)}$, the associativity/normal-ordering formula)
are special cases ($\ell=0$, resp.\ $m=0$).
\end{physics}
\begin{proof}[Proof on paper]
Exactly as in the bosonic case (Chapter~\ref{ch:L03}, Theorem~\ref{l03:thm:borcherds}). Both sides
satisfy the Pascal recursion $B(\ell,m+1,n)=B(\ell,m,n+1)+B(\ell+1,m,n)$; the sign
$\varepsilon(-1)^\ell$ is constant in $m,n$ and changes by $-1$ under $\ell\mapsto\ell+1$, exactly as
$(-1)^\ell$ does in the bosonic recursion. So it suffices to know two boundary lines. At $m=0$ the
identity is the iterate (associativity) formula
$(u_{(\ell)}v)_{(n)}=\sum_i(-1)^i\binom\ell i(u_{(\ell-i)}v_{(n+i)}-\varepsilon(-1)^\ell v_{(\ell+n-i)}u_{(i)})$
(\lean{mode_iterateε}), proved for $\ell\ge0$ from the $n$-th product theorem (Dong's lemma plus
Goddard uniqueness plus translation covariance) and for $\ell<0$ from the twisted normal-ordered
product, where the Koszul factor sits precisely on the reordered branch. At $\ell=0$ it is the
commutator formula $[u_{(m)},v_{(n)}]_\varepsilon=\sum_i\binom mi(u_{(i)}v)_{(m+n-i)}$
(\lean{commBoundε}), which follows from the $\varepsilon$-generic OPE decomposition
(Chapter~\ref{ch:L01}, Theorem~\ref{l01:thm:epsdeltasum}). Induction on $m$ in both directions, with
the anti-diagonal pin of Chapter~\ref{ch:L03}, closes the argument. Translation covariance enters
through the $n$-th product theorem, which is why the statement needs \lean{SuperNearVOA}.
\end{proof}
\begin{status}
\texttt{\#print axioms VOA.SuperNearVOA.superBorcherds} and \lean{VOA.SuperNearVOA.superJacobi}:
\lean{[propext, Classical.choice, Quot.sound]}. Unconditional.
\end{status}

\subsection{Skew-symmetry}

Skew-symmetry $Y(u,z)v=(-1)^{\lseventeenpar u\lseventeenpar v}e^{zT}Y(v,-z)u$ exists twice in the
library.

\begin{theorem}[Super skew-symmetry, unconditional]\label{l17:thm:skew}
Let $V$ be a super near-VOA, $a\in V_{\bar p}$, $b\in V_{\bar q}$. For every $t\in\Z$,
\begin{equation*}
b_{(t)}a=(-1)^{pq}\sum_{j\ge0}\frac{(-1)^{t+j+1}}{j!}\,L_{-1}^{\,j}\bigl(a_{(t+j)}b\bigr).
\end{equation*}
\end{theorem}
\noindent File \leanpath{VOA/\allowbreak Modularity/\allowbreak SuperZhuCommutator.lean}.
\begin{leancode}
theorem superSmul_eq_skew_Lpow {p q : ZMod 2} {a b : 𝒱}
    (ha : a ∈ V.parity p) (hb : b ∈ V.parity q) (t : ℤ) :
    V.smul b t a
      = koszul p q • ∑ᶠ j : ℕ, ((j.factorial : ℂ)⁻¹) •
          (sgn (-t - 1 - (j : ℤ))) •
            ((V.L (-1) ^ j) (V.smul a (t + (j : ℤ)) b))
\end{leancode}
\begin{dictionary}
\lean{V} is a \lean{SuperNearVOA} (file variable). \lean{sgn k} $=(-1)^k$ (\leanpath{Core/\allowbreak
Skew.lean}), so \lean{sgn (-t-1-j)} $=(-1)^{t+j+1}$. \lean{V.L (-1) ^ j} is $L_{-1}^j$; the sum is
finite by VA5. This is Kac's (4.2.1)/(4.8.6) in modes: $b_{(n)}a=p(a,b)\sum_j(-1)^{n+j+1}T^{(j)}(a_{(n+j)}b)$
with $T^{(j)}=T^j/j!$ and $T=L_{-1}$; it is the coefficient form of
$Y(b,z)a=p(a,b)e^{zT}Y(a,-z)b$.
\end{dictionary}
\begin{physics}
Exchanging the two operators in an OPE: $b(z)a(0)=\pm e^{zL_{-1}}a(-z)b(0)$, the translation
moving the insertion point from $0$ to $z$. For two fermions the sign is $-1$; e.g.\ for a free
fermion $\psi_{(0)}\psi=\lseventeenvac$ is symmetric under the exchange only because of the sign.
\end{physics}
\begin{proof}[Proof on paper]
Take the Jacobi identity with $x=\lseventeenvac$, $u=a$, $v=b$, $\ell=t$, $m=-1$, $n=0$. On the left
only the $i=0$ term of the first sum survives and gives $a_{(t-1)}b_{(0)}\lseventeenvac=0$, while the
second sum gives $-\varepsilon(-1)^t\sum_i(-1)^i\binom ti b_{(t-i)}a_{(-1+i)}\lseventeenvac$, in which
only $i=0$ survives: $-\varepsilon(-1)^tb_{(t)}a$. The right side is
$\sum_i\binom{-1}{i}(a_{(t+i)}b)_{(-1-i)}\lseventeenvac=\sum_i(-1)^iL_{-1}^i(a_{(t+i)}b)/i!$ by the
creation formula $w_{(-1-i)}\lseventeenvac=L_{-1}^iw/i!$. Solve for $b_{(t)}a$.
\end{proof}
\begin{status}
\texttt{\#print axioms VOA.SuperNearVOA.superSmul\_eq\_skew\_Lpow}: \lean{[propext,
Classical.choice, Quot.sound]}. Unconditional.
\end{status}

\begin{proposition}[Super skew-symmetry, gated version]\label{l17:prop:skewgated}
In a super VOA $W$, for $u\in W_{\bar p}$, $v\in W_{\bar q}$, if the skew candidate field
$v\mapsto(-1)^{pq}e^{zT}Y(v,-z)u$ is Koszul-local with every $Y(b,z)$ (the ``gate''
\lean{WeakAssocLocε W u p}), then the mode form of skew-symmetry holds.
\end{proposition}
\noindent File \leanpath{VOA/\allowbreak Core/\allowbreak SuperSkew.lean}.
\begin{leancode}
def WeakAssocLocε (W : SuperVertexOperatorAlgebra 𝒱) (u : 𝒱) (p : ZMod 2) : Prop :=
  ∀ (q : ZMod 2) (b : 𝒱), b ∈ W.parity q → Localε (koszul p q) (skewDistε W u p) (W.Y b)
...
theorem skew_modeε (W : SuperVertexOperatorAlgebra 𝒱) (u v : 𝒱) {p q : ZMod 2}
    (hu : u ∈ W.parity p) (hv : v ∈ W.parity q) (hloc : WeakAssocLocε W u p) (n : ℤ) :
    W.smul u n v
      = koszul p q • ∑ᶠ j : ℕ, ((j.factorial : ℂ)⁻¹) • (sgn (-n - 1 - j)) •
          ((W.L (-1) ^ j) ((W.Y v) (-n - 1 - j) u))
\end{leancode}
\begin{dictionary}
\lean{(W.Y v) (-n - 1 - j) u} is the $z^{-n-1-j}$-coefficient, i.e.\ $v_{(n+j)}u$; so this is
Theorem~\ref{l17:thm:skew} with the roles of the two vectors named differently. The gate
\lean{WeakAssocLocε} is never discharged in \leanpath{Core/}: the file reduces it to the sign-free
\lean{SkewVacLocal₀} (\leanpath{Core/\allowbreak SuperWeakAssoc.lean},
\lean{skew_modeε_of_signFree}) and stops. The conclusion is nevertheless true for every super VOA,
by Theorem~\ref{l17:thm:skew} applied to \lean{W.toSuperNearVOA}.
\end{dictionary}
\begin{physics}
Same content as Theorem~\ref{l17:thm:skew}; the gated route follows the bosonic history in which
skew-symmetry was proved from Goddard uniqueness before Borcherds was available.
\end{physics}
\begin{status}
\texttt{\#print axioms VOA.SuperVertexOperatorAlgebra.skew\_modeε}: \lean{[propext,
Classical.choice, Quot.sound]}; conditional on \lean{hloc : WeakAssocLocε W u p}, which is never
proved.
\end{status}

\begin{warning}[title={Finding (F5): a superseded gate and stale documentation}]
The conditional \lean{skew_modeε} is superseded by the unconditional \lean{superSmul_eq_skew_Lpow},
which is derived from \lean{superJacobi}. The remarks ``super skew-symmetry is still gated on
\lean{WeakAssocLocε}'' in \leanpath{Core/\allowbreak ScreeningKernel.lean},
\leanpath{Examples/\allowbreak DSReductionSl2.lean} and \leanpath{Modules/\allowbreak
SuperAdjointModule.lean} (``the \lean{conformal} hypothesis \dots still waits on the super skew gate'')
are stale. Severity: documentation only. Remedy: derive \lean{WeakAssocLocε} (or delete the gated route) from
\lean{superSmul_eq_skew_Lpow}.
\end{warning}

%======================================================================
\section{Automorphisms and the even part}\label{l17:sec:aut}
%======================================================================

\begin{definition}[Even automorphism]\label{l17:def:superaut}
An (even) automorphism of a super VOA $W$ is a linear bijection $g$ with
$g(u_{(n)}v)=(gu)_{(n)}(gv)$ for all $u,v,n$, $g\lseventeenvac=\lseventeenvac$, $g\omega=\omega$, and
$gW_{\bar p}=W_{\bar p}$.
\end{definition}
\noindent File \leanpath{VOA/\allowbreak Core/\allowbreak SuperAutomorphism.lean}.
\begin{leancode}
structure SuperAut (W : SuperVertexOperatorAlgebra 𝒱) where
  toLinearEquiv : 𝒱 ≃ₗ[ℂ] 𝒱
  map_smul' : ∀ (u : 𝒱) (n : ℤ) (v : 𝒱),
    toLinearEquiv (W.smul u n v) = W.smul (toLinearEquiv u) n (toLinearEquiv v)
  map_vacuum : toLinearEquiv W.vacuum = W.vacuum
  map_omega : toLinearEquiv W.conformalVector = W.conformalVector
  map_parity' : ∀ (p : ZMod 2) (u : 𝒱),
    toLinearEquiv u ∈ W.parity p ↔ u ∈ W.parity p
\end{leancode}
\begin{dictionary}
\lean{𝒱 ≃ₗ[ℂ] 𝒱} is a $\C$-linear bijection. Preservation of the conformal weights is a theorem
(\lean{map_mem_grading}: $g$ commutes with $L_0=\omega_{(1)}$). Only parity-preserving automorphisms
are considered; this is the standard notion for twisted modules of super VOAs (Dong--Zhao). The
automorphisms form a group (\lean{instance : Group (SuperAut W)}).
\end{dictionary}
\begin{physics}
Global internal symmetries of the chiral algebra that commute with the stress tensor and with
fermion number. Twisting by such a symmetry gives the twisted sectors (orbifold theory). (In the
library the twisted-sector structure \lean{SuperTwistedModule} is, for every twist of order $N\ge2$,
inhabited only by the zero space; see the warning before Definition~\ref{l17:def:supertwisted}.)
\end{physics}
\begin{status}
A definition; standard axioms.
\end{status}

\begin{theorem}[The parity automorphism]\label{l17:thm:parityaut}
For every super VOA $W$ the fermion-number involution $\sigma=(-1)^F$ is an even automorphism,
$\sigma u=(-1)^pu$ on $W_{\bar p}$, $\sigma^2=1$, and $\sigma$ has order exactly $2$ when
$W_{\bar1}\neq0$ (and is the identity when $W_{\bar1}=0$).
\end{theorem}
\noindent File \leanpath{VOA/\allowbreak Core/\allowbreak SuperAutomorphism.lean}.
\begin{leancode}
noncomputable def parityAut (W : SuperVertexOperatorAlgebra 𝒱) : SuperAut W where
  toLinearEquiv := W.parityGrading.parityLinearEquiv
...
theorem parityAut_homog (W : SuperVertexOperatorAlgebra 𝒱) {p : ZMod 2} {v : 𝒱}
    (hv : v ∈ W.parity p) : W.parityAut v = parSign p • v :=
...
theorem orderOf_parityAut_eq_two (W : SuperVertexOperatorAlgebra 𝒱)
    (hodd : W.parity 1 ≠ ⊥) : orderOf W.parityAut = 2 :=
\end{leancode}
\begin{dictionary}
The hypothesis \lean{hodd} is necessary (without odd vectors $\sigma=1$, \lean{parityAut_eq_one_of_odd_eq_bot}).
The mode law $\sigma(u_{(n)}v)=(\sigma u)_{(n)}(\sigma v)$ is VA1$'$ plus
$\operatorname{parSign}(p+q)=\operatorname{parSign}(p)\operatorname{parSign}(q)$.
\end{dictionary}
\begin{physics}
$(-1)^F$ is a symmetry of every fermionic theory. Twisting by it changes the boundary condition of
fermions around the spatial circle from Neveu--Schwarz to Ramond; inserting it in a trace changes
the temporal boundary condition (the supertrace). The four combinations are the four spin
structures on the torus (§\ref{l17:sec:modularity}). Caution: the automorphism $\sigma$ itself is
genuine, but the library's Ramond-sector carrier \lean{W.ParityTwistedModule M}, i.e.\
\lean{SuperTwistedModule W W.parityAut 2 M}, is inhabited only by $M=0$ (F1,
Theorem~\ref{l17:thm:twistvacuous}); so only the temporal (supertrace) half of this picture has
content in Lean.
\end{physics}
\begin{status}
\texttt{\#print axioms VOA.SuperVertexOperatorAlgebra.orderOf\_parityAut\_eq\_two}: \lean{[propext,
Classical.choice, Quot.sound]}. Unconditional.
\end{status}

\begin{theorem}[The even part is an ordinary VOA]\label{l17:thm:evenpart}
Let $W$ be a super VOA whose even part has only integer weights. Then $W_{\bar0}$, with the restricted
vertex operators, vacuum and conformal vector, is a \lean{VertexOperatorAlgebra} of the same central
charge, with weight spaces $(W_{\bar0})_n=W_n\cap W_{\bar0}$.
\end{theorem}
\noindent File \leanpath{VOA/\allowbreak Core/\allowbreak SuperEvenPart.lean}.
\begin{leancode}
noncomputable def evenPartVOA (W : SuperVertexOperatorAlgebra 𝒱)
    (hint : ∀ h : ℂ, (¬ ∃ n : ℤ, (n : ℂ) = h) → ∀ x ∈ W.parity 0, x ∈ W.grading h → x = 0) :
    VertexOperatorAlgebra (W.parity 0) :=
\end{leancode}
\begin{dictionary}
The carrier is the subtype \lean{W.parity 0}. \lean{hint} is the integrality of the even part, a
genuine hypothesis (it fails for ambient algebras with weight-$\frac52$ states). The identifications
of $Y$, $\lseventeenvac$, $\omega$, $c$ and the grading with the ambient ones are \lean{rfl}-lemmas
(\lean{evenPartVOA_Y_coe}, \lean{evenPartVOA_grading}, \dots), so this is the even part and not some
other VOA structure on the same space.
\end{dictionary}
\begin{physics}
The bosonic subalgebra of a fermionic theory is the $\Z/2$-orbifold by $(-1)^F$; e.g.\ the even
part of symplectic fermions is the triplet algebra $W(2)$ at $c=-2$ (Chapter~\ref{ch:L18}).
\end{physics}
\begin{status}
\texttt{\#print axioms VOA.SuperVertexOperatorAlgebra.evenPartVOA}: \lean{[propext,
Classical.choice, Quot.sound]}. Conditional only on \lean{hint}.
\end{status}

%======================================================================
\section{Super modules, twisted super modules, and the defect}\label{l17:sec:modules}
%======================================================================

\begin{definition}[Super weak module]\label{l17:def:superweak}
A super weak module over a super vertex algebra $W$ is a vector space $M=M_{\bar0}\oplus M_{\bar1}$
with a linear $Y_M:W\to\End(M)[[z^{\pm1}]]$ such that modes add parities, $Y_M(\lseventeenvac,z)=\mathrm{id}$,
$Y_M(u,z)x\in M((z))$, and the super Jacobi identity of Theorem~\ref{l17:thm:superjacobi} holds with
$Y_M$ in place of $Y$, for all homogeneous $u,v$, all $x\in M$ and all $\ell,m,n\in\Z$.
\end{definition}
\noindent File \leanpath{VOA/\allowbreak Modules/\allowbreak SuperWeakModule.lean}.
\begin{leancode}
structure SuperWeakModule (W : SuperVertexAlgebra 𝒱) (M : Type v)
    [AddCommGroup M] [Module ℂ M] where
  YM : 𝒱 →ₗ[ℂ] (End ℂ M)⟦z^±⟧
  parity : ZMod 2 → Submodule ℂ M
  parity_isInternal : DirectSum.IsInternal parity
  parity_mode : ∀ {p q : ZMod 2} (u : 𝒱), u ∈ W.parity p → ∀ (x : M), x ∈ parity q →
    ∀ (n : ℤ), mode (YM u) n x ∈ parity (p + q)
  vacuum_YM : YM W.vacuum = 1
  truncation : ∀ (u : 𝒱) (x : M), ∃ N : ℤ, ∀ n : ℤ, N ≤ n → mode (YM u) n x = 0
  jacobi : ∀ {p q : ZMod 2} (u : 𝒱), u ∈ W.parity p → ∀ (v : 𝒱), v ∈ W.parity q →
    ∀ (x : M) (ℓ m n : ℤ),
    (∑ᶠ i : ℤ, (negOnePow i * cbz ℓ i) •
        (mode (YM u) (ℓ + m - i) (mode (YM v) (n + i) x)
          - (koszul p q * negOnePow ℓ) • mode (YM v) (ℓ + n - i) (mode (YM u) (m + i) x)))
      = ∑ᶠ i : ℤ, cbz m i • mode (YM (W.smul u (ℓ + i) v)) (m + n - i) x
\end{leancode}
\begin{dictionary}
This is the super version of Chapter~\ref{ch:L12}, Definition~\ref{l12:def:weakmodule}, and of Kac's
definition of a module over a vertex superalgebra (Kac §4.1, with the Borcherds identity as the
axiom, as in Li and in Dong--Zhao). Untwisted modules have integer modes, so quantification over all
$\ell,m,n$ is correct here (contrast Definition~\ref{l17:def:supertwisted}). No grading is part of a
weak module.
\end{dictionary}
\begin{physics}
A representation of the fermionic chiral algebra on the state space of a sector; the parity of
$M$ is fermion number in that sector.
\end{physics}
\begin{status}
A definition; standard axioms.
\end{status}

\begin{proposition}[The adjoint super module]\label{l17:prop:adjoint}
Every super near-VOA is a super weak module over itself, with $Y_M=Y$ and its own parity.
\end{proposition}
\noindent File \leanpath{VOA/\allowbreak Modules/\allowbreak SuperAdjointModule.lean}.
\begin{leancode}
def SuperNearVOA.adjointSuperWeakModule (V : SuperNearVOA 𝒱) :
    SuperWeakModule V.toSuperVertexAlgebra 𝒱 where
  YM := V.Y
\end{leancode}
\begin{dictionary}
The \lean{jacobi} field is discharged by \lean{superJacobi} with no rewriting, which certifies that
Theorem~\ref{l17:thm:superjacobi} is exactly the module axiom. First inhabitant with genuinely odd
states.
\end{dictionary}
\begin{physics}
The vacuum sector: the chiral algebra acting on itself.
\end{physics}
\begin{proof}[Proof on paper]
Each field is the corresponding axiom or theorem of $V$.
\end{proof}
\begin{status}
\texttt{\#print axioms VOA.SuperNearVOA.adjointSuperWeakModule}: \lean{[propext, Classical.choice,
Quot.sound]}.
\end{status}

\begin{definition}[Graded super module]\label{l17:def:superord}
A \lean{SuperOrdVModule} over a super VOA $W$ is a super weak module with an internal grading
$M=\bigoplus_{\alpha\in\C}M_\alpha$ by finite-dimensional $L_0$-eigenspaces, real parts of weights
bounded below, and $u_{(n)}M_\alpha\subseteq M_{\alpha+h-n-1}$ for $u\in W_h$.
\end{definition}
\noindent File \leanpath{VOA/\allowbreak Modules/\allowbreak SuperOrdVModule.lean}.
\begin{leancode}
structure SuperOrdVModule (W : SuperVertexOperatorAlgebra 𝒱) (M : Type v)
    [AddCommGroup M] [Module ℂ M]
    extends SuperWeakModule W.toSuperVertexAlgebra M where
  weight : ℂ → Submodule ℂ M
  weight_isInternal : DirectSum.IsInternal weight
  L0_weight : ∀ (α : ℂ) (x : M), x ∈ weight α →
    mode (YM W.conformalVector) 1 x = α • x
  weight_finite : ∀ α, FiniteDimensional ℂ (weight α)
  boundedBelow : ∃ B : ℝ, ∀ α : ℂ, weight α ≠ ⊥ → B ≤ α.re
  weight_mode : ∀ {h α : ℂ} (u : 𝒱), u ∈ W.grading h →
    ∀ (x : M), x ∈ weight α → ∀ (n : ℤ),
      mode (YM u) n x ∈ weight (α + h - n - 1)
\end{leancode}
\begin{dictionary}
The super analogue of \lean{OrdVModule} (Chapter~\ref{ch:L12}, Definition~\ref{l12:def:ordmodule}).
Since $\omega$ is even, $L_0$ preserves parity and each $M_\alpha$ inherits a parity grading
(\lean{weightParityGrading}), on which the basis-free supertrace
$\operatorname{str}f=\tr((-1)^Ff)$ (\lean{supertrace}, \leanpath{Modularity/\allowbreak SuperTrace.lean})
is defined. The only instances constructed in the library are \lean{OrdVModule.toSuper} of bosonic
modules; no \lean{SuperOrdVModule} with a nonzero odd part is built.
\end{dictionary}
\begin{physics}
A sector with finite degeneracies and energy bounded below: the kind of module whose (super)character
$\operatorname{str}_Mq^{L_0-c/24}$ makes sense.
\end{physics}
\begin{status}
A definition; standard axioms.
\end{status}

\subsection{Twisted super modules}

Recall from Chapter~\ref{ch:06} the twisted Jacobi identity in mode form: for $g$ of order dividing
$N$ and $u$ in the eigenspace $V^j$, the field $Y_M(u,z)$ has modes $u_{(r)}$ with $r$ in one coset of
$\Z$ in $\frac1N\Z$, and the identity is required \emph{only} for $r$ on that coset.

\begin{warning}[title={Read this before the definition: for $N\ge2$ the only twisted super module is $M=0$}]
\textbf{At the earlier state of the library against which this chapter was first written, and still in the
released source (its \lean{jacobi} field is unchanged), the structure \lean{SuperTwistedModule W g N M}
defined below is uninhabited for every twist order $N\ge2$ except by the zero space $M=0$.} This
holds for every super VOA $W$ and every even automorphism $g$ (including $g=1$ and the parity
automorphism $g=(-1)^F$, i.e.\ the Ramond sector). It is a theorem, checked by Lean
(Theorem~\ref{l17:thm:twistvacuous} below; the check is not part of this release). The cause is that the \lean{jacobi} field is imposed for
\emph{all} integer labels $k$, not only for those with $k\equiv j\pmod N$ (finding (F1)).
Consequently every library theorem whose hypotheses contain a \lean{SuperTwistedModule} with
$N\ge2$ is a statement about the zero vector space, and \emph{this book does not present any such
theorem as having mathematical content}. Only the case $N=1$ (the untwisted sector, reached through
\lean{SuperOrdVModule.toIdentityTwist}) carries content.
\end{warning}

\begin{definition}[Twisted super module, library sense]\label{l17:def:supertwisted}
For an even automorphism $g$ of a super VOA $W$ and $N>0$ with $g^N=1$, a \lean{SuperTwistedModule W g N M}
is: a linear map $u\mapsto(k\mapsto Y_M(u)_k)$, the integer label $k$ standing for the fractional
mode $u_{(k/N)}$; a parity grading of $M$ with modes adding parities; $\lseventeenvac_{(k/N)}=\delta_{k,-N}$;
lower truncation; support $u_{(k/N)}=0$ unless $k\equiv j\pmod N$ for $u\in W^j$; the Koszul-signed
twisted Jacobi identity \emph{for all integers $\ell,k,m$}; translation covariance; and a
$\C$-grading by finite-dimensional $L_0$-eigenspaces, bounded below, with the fractional weight shift.
\textbf{For $N\ge2$ these axioms force $M=0$} (Theorem~\ref{l17:thm:twistvacuous}); so, at this
commit, this is \emph{not} a definition of a twisted module in the sense of Dong--Li--Mason or
Dong--Zhao except when $N=1$.
\end{definition}
\noindent File \leanpath{VOA/\allowbreak Modules/\allowbreak SuperTwistedModule.lean}.
\begin{leancode}
noncomputable def SuperAut.twistedEigenspace {W : SuperVertexOperatorAlgebra 𝒱}
    (g : SuperAut W) (N : ℕ) (j : ℕ) : Submodule ℂ 𝒱 :=
  Module.End.eigenspace (g.toLinearEquiv : End ℂ 𝒱)
    (Complex.exp (-(2 * (Real.pi : ℂ) * Complex.I * j) / N))
...
structure SuperTwistedModule (W : SuperVertexOperatorAlgebra 𝒱) (g : SuperAut W)
    (N : ℕ) (M : Type v) [AddCommGroup M] [Module ℂ M] where
  order_pos : 0 < N
  order_pow : g ^ N = 1
  YM : 𝒱 →ₗ[ℂ] (ℤ → End ℂ M)
  parity : ZMod 2 → Submodule ℂ M
  parity_isInternal : DirectSum.IsInternal parity
  parity_mode : ∀ {p q : ZMod 2} (u : 𝒱), u ∈ W.parity p →
    ∀ (x : M), x ∈ parity q → ∀ (k : ℤ), YM u k x ∈ parity (p + q)
  vacuum_YM : ∀ k : ℤ, YM W.vacuum k = if k = -(N : ℤ) then 1 else 0
  truncation : ∀ (u : 𝒱) (x : M), ∃ K : ℤ, ∀ k : ℤ, K ≤ k → YM u k x = 0
  support : ∀ (j : ℕ) (u : 𝒱), u ∈ g.twistedEigenspace N j →
    ∀ k : ℤ, (k - j) % (N : ℤ) ≠ 0 → YM u k = 0
  jacobi : ∀ {p q : ZMod 2} (j : ℕ) (u : 𝒱),
    u ∈ g.twistedEigenspace N j → u ∈ W.parity p →
    ∀ (w : 𝒱), w ∈ W.parity q → ∀ (x : M) (ℓ k m : ℤ),
    (∑ᶠ i : ℤ, (negOnePow i * cbz ℓ i) •
        (YM u (ℓ * N + k - i * N) (YM w (m + i * N) x)
          - (koszul p q * negOnePow ℓ) •
            YM w (ℓ * N + m - i * N) (YM u (k + i * N) x))) =
      ∑ᶠ i : ℕ, fractionalBinom ((k : ℂ) / N) i •
        YM (W.smul u (ℓ + i) w) (k + m - i * N) x
\end{leancode}
\begin{dictionary}
The remaining fields (\lean{translation}, \lean{weight}, \lean{weight_isInternal}, \lean{L0_weight},
\lean{weight_finite}, \lean{boundedBelow}, \lean{weight_mode}) are the graded data and are not needed
below. \lean{fractionalBinom x i} $=\binom xi=\prod_{k<i}(x-k)/i!$ for complex $x$. In the
\lean{jacobi} field the label \lean{ℓ * N + k - i * N} is the fractional mode $\ell+r-i$ with
$r=k/N$, and \lean{m + i * N} is $s+i$ with $s=m/N$: the identity reads, for $r=k/N$, $s=m/N$,
\begin{equation*}
\sum_i(-1)^i\binom\ell i\bigl(u_{(\ell+r-i)}w_{(s+i)}-\varepsilon(-1)^\ell w_{(\ell+s-i)}u_{(r+i)}\bigr)
=\sum_{i\ge0}\binom ri(u_{(\ell+i)}w)_{(r+s-i)} ,
\end{equation*}
which is Dong--Li--Mason's mode identity with the Koszul sign (Chapter~\ref{ch:06},
\eqref{ch06:eq:jacobi-mode}). \textbf{But it is required for every $k\in\Z$}, with no hypothesis
$k\equiv j\pmod N$: the restriction that Chapter~\ref{ch:06} added to the bosonic
\lean{TwistedModule.jacobi} is absent. The \lean{support} field uses $k\equiv j$, the bosonic
\lean{TwistedModule} uses $k\equiv-j$ for the same eigenvalue $e^{-2\pi ij/N}$; so a super twisted
module for $g$ corresponds to a bosonic one for $g^{-1}$ (finding (F7), harmless, undocumented).
\end{dictionary}
\begin{physics}
The \emph{intended} meaning is a $g$-twisted sector: fields acquire the monodromy $g$ when transported
around the origin, so their modes are fractional. For $g=(-1)^F$ and $N=2$ odd fields get
half-integral vertex-algebra modes; for a weight-$\frac12$ fermion that means integral physics modes
$\psi_n$, $n\in\Z$, i.e.\ the Ramond sector with its zero mode $\psi_0$. \textbf{The Lean structure
does not model this at the documented commit}: for $N\ge2$, and in particular for the Ramond sector,
its only inhabitant is the zero space (F1), so no Ramond sector, and no twisted sector of any kind,
exists in the library.
\end{physics}
\begin{status}
A definition; standard axioms. \textbf{Uninhabited except by $M=0$ for every $N\ge2$}
(Theorem~\ref{l17:thm:twistvacuous});
meaningful only for $N=1$.
\end{status}

\begin{theorem}[Twisted super modules with $N\ge2$ are zero]\label{l17:thm:twistvacuous}
Let $W$ be any super VOA, $g$ any even automorphism, $N\ge2$, and $T$ any
\lean{SuperTwistedModule W g N M}. Then every $x\in M$ is $0$. In particular $M$ is a
subsingleton, every weight space of $T$ is zero, and every parity-twisted (Ramond) carrier
\lean{W.ParityTwistedModule M} is zero.
\end{theorem}
\begin{dictionary}
\lean{W.ParityTwistedModule M} is the library's abbreviation
\lean{SuperTwistedModule W W.parityAut 2 M}. No hypothesis beyond $N\ge2$ is used; in particular the
statement holds for $g=1$ as well (an identity automorphism with a chosen $N\ge 2$).
\end{dictionary}
\begin{physics}
This is a statement about a mis-specified axiom system, not about conformal field theory. Physically,
the twisted Jacobi identity is the mode expansion of a generating identity in which a field of
monodromy $e^{-2\pi ij/N}$ only has modes on the coset $k\equiv j\pmod N$; the instance used in the
proof below (the vacuum field, which has trivial monodromy, at the label $k=1$) is not part of any
genuine twisted sector when $N\ge2$, because that coefficient of the generating identity is simply
absent. The Lean axiom nevertheless demands it, and the demand kills every vector. In plain words:
the library's ``Ramond sector'' and every other twisted sector of order $N\ge2$ is empty.
\end{physics}
\begin{proof}[Proof on paper]
The identical computation as Chapter~\ref{ch:06}, Theorem~\ref{ch06:thm:defect}. Take
$u=w=\lseventeenvac$ (even, in $W^0$ since $g\lseventeenvac=\lseventeenvac$), $p=q=0$ (Koszul sign
$1$), $\ell=-1$, $k=1$, $m=-N-1$. Left side: each term contains $\lseventeenvac$ at a label
$-N+1-iN$ or $1+iN$; $\lseventeenvac_{(k'/N)}\neq0$ only for $k'=-N$, which would force $N\mid1$.
So the left side is $0$. Right side: $\lseventeenvac_{(-1+i)}\lseventeenvac=\delta_{i,0}\lseventeenvac$,
$\binom{1/N}0=1$, and the label is $1-N-1=-N$, so the right side is $x$. Hence $x=0$.
\end{proof}
\begin{status}
Checked in Lean against the library (the check is not part of this release); not a library theorem. \texttt{\#print axioms} of
the three statements: \lean{[propext, Classical.choice, Quot.sound]}.
\end{status}

\begin{warning}[title={Finding (F1): every twisted super theorem with $N\ge2$ is vacuous}]
This confirms with a machine check a defect recorded earlier in the project: the \lean{jacobi} field of
\lean{SuperTwistedModule} quantifies over all mode indices instead of the coset allowed by \lean{support}
(Section~\ref{ch15:sec:super}). Consequences
inside this chapter's files: the Ramond support theorems \lean{parityTwist_even_support},
\lean{parityTwist_odd_support} and the order-three support checks are true of the zero module only;
\leanpath{Modularity/\allowbreak SuperTwistedTrace.lean} (twisted supertraces for $N\ge2$),
\leanpath{SuperHalfTwistedTrace.lean} (every \lean{SuperTwistedModule W g 2 M}, i.e.\ the whole
``genuine half-integral trace carrier''), \leanpath{SuperHalfSpinRecursionBridge.lean},
\leanpath{SuperSectorRationality.lean} (the \lean{twisted\-Module} field of its realization data) and
\leanpath{Modules/\allowbreak SuperRelativeEtingofFinite.lean} (twisted-module carriers), and the
twisted zero-mode action of \leanpath{Modules/\allowbreak SuperTwistedZhu.lean}
(\lean{SuperTwistedModule.twistedZeroModeBlock}, \lean{twistedZeroModeAction}: the bridge from
twisted modules to the twisted Zhu quotient of Definition~\ref{l17:def:twistedzhu}) are about the
zero space whenever $N\ge2$. In particular every half-integral trace and supertrace is identically
zero, and \lean{SuperTwistedModule.IsIrreducible} for $N\ge2$ is satisfied only by zero modules
(§\ref{l17:sec:status}). The identity-twist case $N=1$ (\lean{SuperOrdVModule.toIdentityTwist}, used
for instance by the torus recursion of \leanpath{Modularity/\allowbreak SuperTorusRecursionClose.lean})
is unaffected. Severity: serious. Remedy: add the hypothesis \texttt{(k - j) \% (N : ℤ) = 0 →} to
\lean{jacobi}. Note the sign: the super \lean{support} field uses $k\equiv j$, so the restriction is
$(k-j)\bmod N=0$; the bosonic \lean{TwistedModule.jacobi} of Chapter~\ref{ch:06} uses
\texttt{(k + j) \% (N : ℤ) = 0} because its support is $k\equiv-j$ (F7). Copying the bosonic
hypothesis verbatim would be wrong here.
\end{warning}

%======================================================================
\section{The four free-field examples}\label{l17:sec:examples}
%======================================================================

All four examples are built on infinite-dimensional Fock spaces, and this is forced: locality makes
the twisted commutator a finite sum of derivatives of the delta function, which has infinite support
along the diagonal, so \emph{no finite-dimensional space carries a genuinely fermionic local pair}
(the project's internal records note this negative result, found when finite-dimensional
fermionic examples were attempted). The fermionic Fock spaces are exterior algebras, realised in Mathlib as Clifford algebras of
the zero quadratic form; creation is left multiplication by a generator and annihilation is
contraction against a coordinate functional, and Mathlib's \lean{CliffordAlgebra.contractLeft_ι_mul}
\emph{is} the canonical anticommutation relation. The bosonic Fock space of $\beta\gamma$ is a
polynomial ring, with multiplication and partial derivatives.

\subsection{The free fermion}\label{l17:sec:ff}

\begin{definition}[The Neveu--Schwarz free fermion]\label{l17:def:psi}
Let $\mathcal F=\Lambda(\bigoplus_{k\in\N}\C e_k)$. Let $c_k$ be exterior multiplication by $e_k$ and
$a_k$ contraction with the $k$-th coordinate. The library defines $\psi:\Z\to\End\mathcal F$ by
$\psi(n)=a_n$ for $n\ge0$ and $\psi(n)=c_{-n-1}$ for $n<0$. \textbf{Caution (F3):} the library
reads $\psi(n)$ as the coefficient of $z^{n}$, not as the mode $\psi_{(n)}$; with that reading $\psi$ is
\emph{not} a field (Theorem~\ref{l17:thm:psinotfield}), and the library contains no super vertex
algebra of the free fermion.
\end{definition}
\noindent File \leanpath{VOA/\allowbreak Examples/\allowbreak FreeFermion.lean}.
\begin{leancode}
noncomputable abbrev Fock := CliffordAlgebra Qzero
noncomputable def cre (k : ℕ) : Module.End ℂ Fock :=
  LinearMap.mulLeft ℂ (ι Qzero (Finsupp.single k (1 : ℂ)))
noncomputable def ann (k : ℕ) : Module.End ℂ Fock := contractLeft (Finsupp.lapply k)
...
noncomputable def psi : FormalDist (Module.End ℂ Fock) :=
  fun n => if 0 ≤ n then ann n.toNat else cre (-n - 1).toNat
\end{leancode}
\begin{dictionary}
\lean{Qzero} is the zero quadratic form on \lean{ModeSpace = ℕ →₀ ℂ} (finitely supported sequences),
so \lean{Fock} is the exterior algebra. The docstring says ``$\psi(z)=\sum_{n\in\Z}\psi_{(n)}z^{-n-1}$,
creation modes for $n\le-1$'', i.e.\ it reads the argument of \lean{psi} as a \emph{mode} index. But
\lean{psi} has type \lean{FormalDist}, whose argument is the \emph{$z$-exponent}
(§\ref{l17:sec:conv}). Read as the library reads it, $\psi(z)=\sum_k\psi(k)z^k$ has creation
operators at all negative powers of $z$, so $\psi(z)x\notin\mathcal F((z))$: it is not a field
(Theorem~\ref{l17:thm:psinotfield}).
\end{dictionary}
\begin{physics}
The Majorana fermion of weight $\frac12$, $\psi(z)\psi(w)\sim\frac1{z-w}$, central charge
$c=\frac12$ (the chiral half of the critical Ising model). In the Neveu--Schwarz sector its modes are
$\psi_r$, $r\in\Z+\frac12$; here $c_k=\psi_{-k-1/2}$ and $a_k=\psi_{k+1/2}$.
\end{physics}
\begin{status}
A definition; standard axioms.
\end{status}

\begin{theorem}[Canonical anticommutation relations and locality]\label{l17:thm:psicar}
$\psi(m)\psi(n)+\psi(n)\psi(m)=\delta_{m+n+1,0}$ for all $m,n\in\Z$, and $\psi$ is $(-1)$-local with
itself.
\end{theorem}
\noindent File \leanpath{VOA/\allowbreak Examples/\allowbreak FreeFermion.lean}.
\begin{leancode}
theorem psi_anticomm (m n : ℤ) :
    psi m * psi n + psi n * psi m = if m + n + 1 = 0 then 1 else 0 := by
...
theorem psi_localε : Localε (-1) psi psi :=
  localε_neg_one_of_anticomm 1 psi_anticomm
\end{leancode}
\begin{dictionary}
Both statements are invariant under $n\mapsto-n-1$, so they are true whichever indexing one has in
mind. \lean{psi_localε} is the first inhabitant of \lean{Localε} at a sign other than $1$; it checks
that the sign conventions are right. It does \emph{not} make $\psi$ a field.
\end{dictionary}
\begin{physics}
$\{\psi_r,\psi_s\}=\delta_{r+s,0}$: the Clifford algebra of the free Majorana fermion, whose
anticommutator is $\delta(z-w)$, a simple pole in the OPE.
\end{physics}
\begin{proof}[Proof on paper]
Case analysis on the signs of $m,n$, using $c_jc_k+c_kc_j=0$, $a_ja_k+a_ka_j=0$, $a_jc_k+c_ka_j=\delta_{jk}$;
locality follows from the free-field criterion with $S=\mathrm{id}$.
\end{proof}
\begin{status}
\texttt{\#print axioms} of \lean{VOA.FreeFermion.psi\_anticomm}, \lean{VOA.FreeFermion.psi\_localε}:
\lean{[propext, Classical.choice, Quot.sound]}.
\end{status}

\begin{theorem}[$\psi$ is not a field]\label{l17:thm:psinotfield}
$\neg\,\mathrm{IsField}(\psi)$: there is no $B$ with $\psi(n)\lseventeenvac=0$ for all $n<B$.
\end{theorem}
\begin{dictionary}
\lean{IsField a} $:\Leftrightarrow\forall x\,\exists B\,\forall n<B,\ a(n)x=0$ (\leanpath{Core/\allowbreak
Dong.lean}). The witness is $x=1\in\mathcal F$.
\end{dictionary}
\begin{physics}
A field must annihilate any given state after finitely many negative powers of $z$ (finitely many
singular terms in an OPE with a fixed state). An operator family that keeps creating new particles at
ever more singular orders is not a field.
\end{physics}
\begin{proof}[Proof on paper]
For $n<0$, $\psi(n)1=c_{-n-1}1=e_{-n-1}\neq0$, because $a_k(c_k1)=1-c_k(a_k1)=1$.
\end{proof}
\begin{status}
Checked in Lean against the library (the check is not part of this release); not a library theorem. \texttt{\#print axioms}: \lean{[propext, Classical.choice, Quot.sound]}. Consequently \textbf{the free
fermion is not announceable}: the library has only anticommuting operator families that are not fields, and no vertex
superalgebra of the free fermion.
\end{status}

\begin{warning}[title={Findings (F3): the free fermion is not a vertex superalgebra in Lean}]
The same inverted orientation occurs in \lean{BCGhosts.bfield}, \lean{BCGhosts.cfield}
(\leanpath{Examples/\allowbreak BCGhosts.lean}) and in \lean{BetaGamma.beta}, \lean{BetaGamma.gamma}
(\leanpath{Examples/\allowbreak BetaGamma.lean}); the later files \leanpath{GhostSystem.lean} and
\leanpath{BetaGammaField.lean} say so in their docstrings and rebuild $bc$ and $\beta\gamma$ in the
correct orientation. The free fermion has no such rebuild: \textbf{the library contains no
\lean{SuperVertexAlgebra} or \lean{SuperVertexOperatorAlgebra} for the free fermion}, only the
operators of Definition~\ref{l17:def:psi}. The docstrings of \leanpath{FreeFermion.lean}
(``creation modes for $n\le-1$''), \leanpath{BetaGamma.lean} and \leanpath{BetaGammaVirasoro.lean}
(which says \lean{BetaGamma} ``indexes fields by the $z$-exponent'', contradicting
\leanpath{BetaGammaField.lean}) are misleading. Severity: moderate (definitions do not mean what their
docstrings say; no downstream theorem is affected because the only consumers are locality
statements). Remedy: reindex $n\mapsto -n-1$ or delete the three early files.
\end{warning}

\subsection{The bc ghost system}\label{l17:sec:bc}

\begin{definition}[$bc$ fields]\label{l17:def:bc}
On the exterior algebra $\Lambda\bigl(\bigoplus_{k\in\N,\,a\in\{0,1\}}\C e_{k,a}\bigr)$ (species $0=b$,
$1=c$) define, for $a\in\{0,1\}$, the distribution $\mathrm{gh}_a$ with $z^j$-coefficient equal to
creation of $e_{j,a}$ for $j\ge0$ and to annihilation against the \emph{conjugate} species $e_{-j-1,a+1}$
for $j<0$; $b=\mathrm{gh}_0$, $c=\mathrm{gh}_1$.
\end{definition}
\noindent File \leanpath{VOA/\allowbreak Examples/\allowbreak GhostSystem.lean}.
\begin{leancode}
noncomputable def ghField (a : Fin 2) : FormalDist (End ℂ FockSF) :=
  fun j => if 0 ≤ j then creGh a j.toNat else annGh a (-j - 1).toNat
noncomputable def bfield : FormalDist (End ℂ FockSF) := ghField 0
noncomputable def cfield : FormalDist (End ℂ FockSF) := ghField 1
\end{leancode}
\begin{dictionary}
Here the $z$-exponent indexing is respected: creation at $j\ge0$. The carrier \lean{FockSF} is reused
from the symplectic-fermion files (it is the same exterior algebra). \lean{creGh a k} $=b$- or
$c$-creation $a_{(-k-1)}$; \lean{annGh a k} contracts against species $a+1$, which is what makes $b$
and $c$ conjugate to each other.
\end{dictionary}
\begin{physics}
The reparametrisation ghosts of the bosonic string are a $bc$ system of weights $(\lambda,1-\lambda)$
with $\lambda=2$; the modes $c_n$ of the weight $-1$ ghost and $b_n$ of the weight-$2$ antighost
generate the ghost Fock space of BRST quantisation. At $\lambda=1$ ($b$ of weight $1$, $c$ of
weight $0$) the system appears in the Drinfeld--Sokolov reduction of $\widehat{\mathfrak{sl}}_2$.
\end{physics}
\begin{status}
Definitions; standard axioms.
\end{status}

\begin{theorem}[$bc$ anticommutators, locality, and the super vertex algebra]\label{l17:thm:bc}
$\{\mathrm{gh}_a(j),\mathrm{gh}_b(k)\}=\delta_{j+k+1,0}\,\omega^{ab}$ with $\omega^{ab}=1$ for $a\neq b$
and $0$ for $a=b$; all ghost fields are mutually $(-1)$-local fields; and there is a super vertex algebra
\lean{bcGhostSVA} on the Fock space, graded by $\operatorname{wt}(b_{(-k-1)})=k+1$,
$\operatorname{wt}(c_{(-k-1)})=k$ (i.e.\ $\lambda=1$), with Clifford parity, whose vertex operator is
produced by super reconstruction.
\end{theorem}
\noindent File \leanpath{VOA/\allowbreak Examples/\allowbreak GhostSystem.lean}.
\begin{leancode}
theorem ghField_anticomm (a b : Fin 2) (j k : ℤ) :
    ghField a j * ghField b k + ghField b k * ghField a j
      = (if j + k + 1 = 0 then ghPair a b else 0) • (1 : End ℂ FockSF) := by
...
theorem ghField_localε (a b : Fin 2) : Localε (-1) (ghField a) (ghField b) := by
...
theorem ghField_isField (a : Fin 2) : IsField (ghField a) := by
...
noncomputable def bcGhostSVA : SuperVertexAlgebra FockSF where
  grading := gradingGh
  isInternal := gradingGh_isInternal
  finrank := gradingGh_finrank
  parity := evenOdd QSF
\end{leancode}
\begin{dictionary}
\lean{ghPair a b} is $1$ when $b=a+1$ in \lean{Fin 2} (species differ) and $0$ otherwise.
\lean{evenOdd QSF} is Mathlib's $\Z/2$-grading of the Clifford algebra by word length mod $2$.
Finite-dimensionality of the weight spaces at $\lambda=1$ is not automatic (the weight-$0$ mode $c_{(-1)}$
can be repeated) and uses that a creation monomial with a repeated mode vanishes. Spin--statistics
fails ($b_{(-1)}\lseventeenvac$ is odd of weight $1$). There is no conformal vector: \lean{bcGhostSVA}
is not a \lean{SuperNearVOA}.
\end{dictionary}
\begin{physics}
$b(z)c(w)\sim\frac1{z-w}$: a first-order pole between two fermionic fields of conjugate weights. The
additional $\Z$-grading by ghost number (\lean{ghostGrading}, $b$ has $-1$, $c$ has $+1$) reduces mod $2$
to the parity.
\end{physics}
\begin{proof}[Proof on paper]
The anticommutators are the Clifford relations; locality is the free-field criterion; lower
truncation holds because at negative exponent the operators are annihilations of increasing level,
which kill any fixed monomial. The super vertex algebra comes from super reconstruction applied to
the word fields (iterated twisted normal-ordered products of derivatives of $b$ and $c$); the two
compatibility axioms VA1, VA1$'$ are checked word by word.
\end{proof}
\begin{status}
\texttt{\#print axioms} of \lean{VOA.Ghost.ghField\_anticomm}, \lean{VOA.Ghost.bcGhostSVA}:
\lean{[propext, Classical.choice, Quot.sound]}. Unconditional, \textbf{in the library's weak sense}: \lean{SuperVertexAlgebra}
has no translation-covariance axiom (finding (F2), Theorem~\ref{l17:thm:nocov}), so \lean{bcGhostSVA} is a weaker statement
than ``the $bc$ system is a vertex superalgebra'' in the standard sense; no super Jacobi identity is proved for it, and its
Virasoro relation is open.
\end{status}

\begin{definition}[Ghost central charge and the ghost Virasoro statement]\label{l17:def:ghostvir}
For $\lambda\in\C$ let $c(\lambda)=-2(6\lambda^2-6\lambda+1)$ and
$L(z)=-\lambda{:}b\,\partial c{:}+(1-\lambda){:}(\partial b)c{:}$. \lean{IsGhostVirasoro λ} is the
proposition that the modes of $L(z)$ satisfy the Virasoro relations with central charge $c(\lambda)$.
\end{definition}
\noindent File \leanpath{VOA/\allowbreak Examples/\allowbreak GhostSystem.lean}.
\begin{leancode}
def ghostCentralCharge (lam : ℂ) : ℂ := -2 * (6 * lam ^ 2 - 6 * lam + 1)
...
theorem ghostCentralCharge_two : ghostCentralCharge 2 = -26 := by
...
def IsGhostVirasoro (lam : ℂ) : Prop :=
  ∀ m n : ℤ, Lgh lam m * Lgh lam n - Lgh lam n * Lgh lam m
    = ((m : ℂ) - (n : ℂ)) • Lgh lam (m + n)
      + (if m + n = 0 then ((m : ℂ) ^ 3 - (m : ℂ)) / 12 * ghostCentralCharge lam else 0)
          • (1 : End ℂ FockSF)
\end{leancode}
\begin{dictionary}
\lean{Lgh lam n} is the $z^{-n-2}$-coefficient of \lean{LghField lam}, the stated normal-ordered
bilinear at Koszul sign $-1$ (so $L_n$ for a weight-$2$ field). \lean{IsGhostVirasoro} is a
\emph{definition of a proposition}; it is \textbf{not proved for any $\lambda$}.
\lean{ghostCentralCharge_two} is the arithmetic fact $-2(24-12+1)=-26$; it says nothing about the
ghost system.
\end{dictionary}
\begin{physics}
$c(\lambda)=-3(2\lambda-1)^2+1$. At $\lambda=2$, $c=-26$: the ghost anomaly cancels the $c=26$ of the
matter in the critical bosonic string (Friedan--Martinec--Shenker). At $\lambda=1$, $c=-2$; at
$\lambda=\frac12$ the $bc$ system is a complex (Dirac) fermion, $c=+1$.
\end{physics}
\begin{status}
\texttt{\#print axioms VOA.Ghost.ghostCentralCharge\_two}: \lean{[propext, Classical.choice,
Quot.sound]}. The Virasoro relation for the $bc$ system is \textbf{open} in the library.
\end{status}

\begin{warning}[title={Finding (F6): what the bc files do not contain}]
The $bc$ files contain no stress tensor theorem, no \lean{SuperNearVOA}, no grading for $\lambda\neq1$,
and in particular nothing about the $\lambda=2$ string ghosts beyond the polynomial identity
$c(2)=-26$. The docstring of \leanpath{GhostSystem.lean} is honest about this. Severity: open gap (an
advertised physical value is only arithmetic), in the same family as the other places in the library where a physically named
value is checked only as arithmetic.
\end{warning}

\subsection{The beta-gamma system}\label{l17:sec:bg}

\begin{definition}[$\beta\gamma$ fields]\label{l17:def:bg}
On $\C[x_k,y_k:k\in\N]$ let $\beta(z)=\sum_j\beta_jz^j$ with $\beta_j$ = multiplication by $x_j$ for
$j\ge0$ and $\partial/\partial y_{-j-1}$ for $j<0$, and $\gamma_j$ = multiplication by $y_j$ for $j\ge0$
and $-\partial/\partial x_{-j-1}$ for $j<0$.
\end{definition}
\noindent File \leanpath{VOA/\allowbreak Examples/\allowbreak BetaGammaField.lean}.
\begin{leancode}
noncomputable def betaZ : FormalDist (Module.End ℂ Fock) :=
  fun j => if 0 ≤ j then mulX (Sum.inl j.toNat) else der (Sum.inr (-j - 1).toNat)
...
noncomputable def gammaZ : FormalDist (Module.End ℂ Fock) :=
  fun j => if 0 ≤ j then mulX (Sum.inr j.toNat) else -der (Sum.inl (-j - 1).toNat)
\end{leancode}
\begin{dictionary}
\lean{Fock = MvPolynomial (ℕ ⊕ ℕ) ℂ}; \lean{Sum.inl k} is $x_k$ ($\beta$-creation $\beta_{(-k-1)}$),
\lean{Sum.inr k} is $y_k$. These are the correctly oriented fields; \lean{betaZ j = BetaGamma.beta (-j-1)}
(\lean{betaZ_eq}). The relation is $[\beta_{(m)},\gamma_{(n)}]=\delta_{m+n+1,0}$ (bosonic commutator,
$\varepsilon=1$).
\end{dictionary}
\begin{physics}
Symplectic bosons: commuting fields of weights $(\lambda,1-\lambda)$, $\beta(z)\gamma(w)\sim\frac{1}{z-w}$.
At $\lambda=\frac32$ they are the superconformal ghosts of the RNS superstring ($c=+11$); at
$\lambda=1$, $c=+2$; at $\lambda=\frac12$ (the library's choice), $c=-1$. In general
$c_{\beta\gamma}(\lambda)=-c_{bc}(\lambda)$.
\end{physics}
\begin{status}
Definitions; standard axioms.
\end{status}

\begin{theorem}[The $\beta\gamma$ super VOA]\label{l17:thm:bgsvoa}
There is a purely even super VOA \lean{bgSVOA} on $\C[x_k,y_k]$ with $x_k,y_k$ of weight $k+\frac12$,
$Y(x_0,z)=\beta(z)$, $Y(y_0,z)=\gamma(z)$, conformal vector
$\omega=\frac12(x_0y_1-x_1y_0)$ and central charge $-1$. It has a nonzero weight space at the
non-integer weight $\frac12$ and violates spin--statistics.
\end{theorem}
\noindent File \leanpath{VOA/\allowbreak Examples/\allowbreak BetaGammaInstance.lean}.
\begin{leancode}
noncomputable def bgSVA : SuperVertexAlgebra Fock where
  grading := gradingBG
  isInternal := gradingBG_isInternal
  finrank := gradingBG_finiteDimensional
  parity := (ParityGrading.even Fock).parity
...
theorem bgSVA_Y_beta : bgSVA.Y (X (Sum.inl 0) : Fock) = betaZ := by
...
theorem bgSVA_grading_notMem_intCast :
    ∃ c : ℂ, bgSVA.grading c ≠ ⊥ ∧ ∀ n : ℤ, (n : ℂ) ≠ c :=
...
theorem bgSVA_not_spinStatistics : ¬ SpinStatistics bgSVA := by
...
noncomputable def bgOmega : Fock :=
  (1 / 2 : ℂ) • stWordBG [Sum.inl 0, Sum.inr 1]
    - (1 / 2 : ℂ) • stWordBG [Sum.inl 1, Sum.inr 0]
...
noncomputable def bgSVOA : SuperVertexOperatorAlgebra Fock :=
  bgSNVOA.toSVOA (by simpa only using gradingBG_zero) (fun h hh => gradingBG_eq_bot_of_re_neg hh)
\end{leancode}
\noindent File \leanpath{VOA/\allowbreak Examples/\allowbreak BetaGammaVirasoro.lean}.
\begin{leancode}
theorem Lbg_virasoro_id (m n : ℤ) :
    Lbg m * Lbg n - Lbg n * Lbg m
      = ((m : ℂ) - (n : ℂ)) • Lbg (m + n)
        + (if m + n = 0 then ((m : ℂ) ^ 3 - (m : ℂ)) / 12 * (-1) else 0)
            • (LinearMap.id : Module.End ℂ Fock)
\end{leancode}
\begin{dictionary}
\lean{gradingBG c} is the space of weighted-homogeneous polynomials of weight $c$ for the weight
function $x_k,y_k\mapsto k+\frac12$ (\lean{bgWt}). \lean{stWordBG [i, j]} is the monomial $X_iX_j$, so
\lean{bgOmega} $=\frac12x_0y_1-\frac12x_1y_0$ (\lean{bgOmega_eq}). \lean{bgSNVOA} is the super near-VOA
whose \lean{virasoro} field is \lean{Lbg_virasoro} transported (\lean{bgOmega_eq_BGVir}: this $\omega$
is the state whose modes are the \lean{Lbg n}), and \lean{toSVOA} adds $\lseventeenV_0=\C1$ and
positivity. \lean{bgSVA_grading_notMem_intCast} is the non-vacuity of Axis A: the grading cannot be
reindexed by $\Z$. The weights $h_\beta=h_\gamma=\frac12$ are derived in \leanpath{BetaGammaVirasoro.lean}
from primarity (\lean{weights_sum_one}: $h_\beta+h_\gamma=1$ is forced, symmetry gives $\frac12$).
\end{dictionary}
\begin{physics}
$T(z)=\frac12{:}\beta\partial\gamma{:}-\frac12{:}\partial\beta\gamma{:}$, the stress tensor at
$\lambda=\frac12$, $c=-1$. The state $\beta_{-1/2}|0\rangle$ is a boson of spin $\frac12$: the
theory is non-unitary and violates spin--statistics in the opposite direction to symplectic fermions.
\end{physics}
\begin{proof}[Proof on paper]
Reconstruction from the monomial words (the word fields are normal-ordered products of derivatives
of $\beta,\gamma$, all bosonic). The Virasoro relation is a direct computation with normally ordered
bilinears $L_n=\sum_jA_j{:}\beta\gamma{:}$; the coefficient family $A_j$ is forced by requiring
$\beta,\gamma$ to be primary of weight $\frac12$, and the central term is the usual double
contraction, here $-1$ because the fields are bosonic.
\end{proof}
\begin{status}
\texttt{\#print axioms} of \lean{VOA.BGInst.bgSVOA}, \lean{VOA.BGInst.bgSVA\_not\_spinStatistics}:
\lean{[propext, Classical.choice, Quot.sound]}. Unconditional. \textbf{Purely even}: the parity is \lean{ParityGrading.even},
so this instance never exercises the Koszul sign rule; it tests the half-integral weight axis only.
\end{status}

\begin{definition}[Super $C_2$, Poisson relation span, Etingof finiteness]\label{l17:def:c2}
For a super VOA $W$: $C_2(W)=\operatorname{span}\{u_{(-2)}v\}$; $W$ is \emph{super $C_2$-cofinite}
if $W/C_2(W)$ is finite-dimensional; $R=C_2(W)+\operatorname{span}\{a_{(0)}b\}$, and $W$ is
\emph{super Etingof-finite} if $W/R$ is finite-dimensional.
\end{definition}
\noindent File \leanpath{VOA/\allowbreak Modularity/\allowbreak GammaThetaStatement.lean}.
\begin{leancode}
def superO₂ (W : SuperVertexOperatorAlgebra 𝒱) : Submodule ℂ 𝒱 :=
  Submodule.span ℂ {x | ∃ u v : 𝒱, x = W.toSuperVertexAlgebra.smul u (-2) v}
def SuperC2Cofinite (W : SuperVertexOperatorAlgebra 𝒱) : Prop :=
  FiniteDimensional ℂ (𝒱 ⧸ superO₂ W)
\end{leancode}
\noindent File \leanpath{VOA/\allowbreak Modularity/\allowbreak SuperPoissonTrace.lean}.
\begin{leancode}
def superZeroModeSpan (W : SuperVertexOperatorAlgebra 𝒱) : Submodule ℂ 𝒱 :=
  Submodule.span ℂ {x : 𝒱 | ∃ a b : 𝒱, x = W.toSuperNearVOA.smul a 0 b}
...
def superPoissonRelSpan (W : SuperVertexOperatorAlgebra 𝒱) : Submodule ℂ 𝒱 :=
  superO₂ W ⊔ superZeroModeSpan W
...
def SuperEtingofFinite (W : SuperVertexOperatorAlgebra 𝒱) : Prop :=
  FiniteDimensional ℂ (SuperPoissonHP₀ W)
\end{leancode}
\begin{dictionary}
\lean{SuperPoissonHP₀ W} is the quotient $\lseventeenV/R$ (\lean{SuperEtingofFinite_iff}). On a bosonic
VOA these are literally the bosonic notions of Chapter~\ref{ch:L12}, Definitions~\ref{l12:def:c2}
and~\ref{l12:def:etingof} (\lean{superC2Cofinite_toSuper_iff} is \lean{Iff.rfl}). $W/C_2(W)$ is
Zhu's Poisson superalgebra $R_W$ and $W/R$ is its zeroth Poisson homology $R_W/\{R_W,R_W\}$; the
library does not construct the Poisson structure on the quotient, only the vector space. Super
Etingof-finiteness is the hypothesis of Arakawa--Kawasetsu (finitely many symplectic leaves implies it,
by Etingof--Schedler).
\end{dictionary}
\begin{physics}
$C_2$-cofinite (``lisse'') theories have finitely many sectors and modular characters; quasi-lisse
theories (finitely many symplectic leaves in the associated variety) still have characters satisfying
modular differential equations. Higgs-branch chiral algebras of 4d $\mathcal N=2$ theories are
quasi-lisse; free $\beta\gamma$ is the simplest (associated variety $\C^2$, one leaf).
\end{physics}
\begin{status}
Definitions; standard axioms.
\end{status}

\begin{theorem}[$\beta\gamma$ is quasi-lisse but not lisse]\label{l17:thm:bgquasi}
\lean{bgSVOA} is not super $C_2$-cofinite, and it is super Etingof-finite (indeed $R$ is everything).
\end{theorem}
\noindent File \leanpath{VOA/\allowbreak Examples/\allowbreak BetaGammaC2LowerBound.lean}.
\begin{leancode}
theorem not_SuperC2Cofinite_bgSVOA : ¬ SuperC2Cofinite bgSVOA := by
\end{leancode}
\noindent File \leanpath{VOA/\allowbreak Examples/\allowbreak BetaGammaQuasiLisse.lean}.
\begin{leancode}
theorem superPoissonRelSpan_eq_top : superPoissonRelSpan bgSVOA = ⊤ :=
...
theorem bgSVOA_strict_halfIntegral_quasiLisse :
    ¬ SuperC2Cofinite bgSVOA ∧ SuperEtingofFinite bgSVOA :=
\end{leancode}
\begin{dictionary}
Both halves are unconditional. ``$R=\top$'' means $W/R=0$, which is finite-dimensional; this is not
a degenerate trick: for the Weyl-type Poisson algebra $\C[\beta,\gamma]$ with $\{\beta,\gamma\}=1$ every
polynomial is a bracket, $f=\{\beta,\int f\,d\gamma\}$, so $HP_0=0$.
\end{dictionary}
\begin{physics}
The associated variety of $\beta\gamma$ is the symplectic plane $\C^2$: positive-dimensional
(not lisse) but a single symplectic leaf (quasi-lisse). Its vacuum character
$q^{1/24}\prod_{n}(1-q^{n-1/2})^{-2}$ is modular only for $\Gamma_\theta$-type subgroups and only up
to the subtleties of non-$C_2$ theories.
\end{physics}
\begin{proof}[Proof on paper]
Not $C_2$: the powers $x_0^n$ survive in $W/C_2$; a linear detector (the coefficient of $x_0^n$) kills
every $u_{(-2)}v$ by a weight and $\beta$-minus-$\gamma$ charge count, and the detectors separate the
classes $[x_0^n]$, which are therefore linearly independent. Etingof-finite: the zero mode of $x_0$
acts as $\partial/\partial y_0$ and that of $y_0$ as $-\partial/\partial x_0$, so every derivative
image lies in $R$; partial differentiation is surjective on polynomials over $\C$.
\end{proof}
\begin{status}
\texttt{\#print axioms} of \lean{VOA.BGNonC2.not\_SuperC2Cofinite\_bgSVOA} and
\lean{VOA.BGInst.bgSVOA\_strict\_halfIntegral\_quasiLisse}: \lean{[propext, Classical.choice,
Quot.sound]}. Unconditional.
\end{status}

\subsection{Symplectic fermions}\label{l17:sec:sf}

\begin{definition}[Symplectic fermion fields]\label{l17:def:chi}
On the exterior algebra $\mathrm{Fock}_{SF}$ on generators $e_{k,a}$ ($k\in\N$, $a\in\{0,1\}$), the
fields $\chi^a(z)$ have $z^j$-coefficient creation of $e_{j,a}$ for $j\ge0$ and a contraction weighted by
the symplectic form for $j<0$, with
$\{\chi^a(j),\chi^b(k)\}=(-j-1)\,\delta_{j+k+2,0}\,\omega^{ab}$.
\end{definition}
\noindent File \leanpath{VOA/\allowbreak Examples/\allowbreak SymplecticFermion.lean}.
\begin{leancode}
noncomputable def chi (a : Fin 2) : FormalDist (Module.End ℂ FockSF) :=
  fun j => if 0 ≤ j then creSF a j.toNat else annSF a (-j - 1).toNat
...
theorem chi_anticomm (a b : Fin 2) (j k : ℤ) :
    chi a j * chi b k + chi b k * chi a j
      = (if j + k + 2 = 0 then (((-j - 1 : ℤ) : ℂ) * sfOmega a b) else 0)
        • (1 : Module.End ℂ FockSF) := by
...
theorem chi_localε (a b : Fin 2) : Localε (-1) (chi a) (chi b) := by
\end{leancode}
\begin{dictionary}
In modes $\chi^a_m:=\chi^a(-m-1)$ (\lean{chiM}) the relation reads $\{\chi^a_m,\chi^b_n\}=m\,\delta_{m+n,0}\,\omega^{ab}$,
with $\omega^{ab}$ the antisymmetric form \lean{sfOmega} ($\omega^{01}=-\omega^{10}$). The factor $m$
makes the anticommutator a derivative of the delta function, so the locality order is $2$, not $1$
(proved through \lean{deltaDeriv 1}, not the first-order free-field criterion). Weight of $\chi^a$: $1$.
\end{dictionary}
\begin{physics}
Kausch's symplectic fermions: a pair of weight-$1$ fermions with $\chi^a(z)\chi^b(w)\sim\frac{\omega^{ab}}{(z-w)^2}$,
central charge $c=-2$; equivalently the $\eta\xi$ system with $\xi$ replaced by $\partial\xi$. The
prototypical logarithmic CFT: its even part is the triplet algebra $W(2)$, and its zero modes
$\chi^a_0$ generate a Grassmann algebra on non-vacuum modules, responsible for Jordan blocks of $L_0$.
\end{physics}
\begin{status}
\texttt{\#print axioms VOA.SympFerm.chi\_anticomm}: \lean{[propext, Classical.choice, Quot.sound]}.
\end{status}

\begin{theorem}[The symplectic-fermion super VOA]\label{l17:thm:sfsvoa}
There is a super VOA \lean{sympFermSVOA} on $\mathrm{Fock}_{SF}$, with $Y(\chi^a_{-1}\lseventeenvac,z)=\chi^a(z)$,
integer weights in both parities, conformal vector $\omega_{SF}=-\chi^0_{-1}\chi^1_{-1}\lseventeenvac$, and
central charge $-2$.
\end{theorem}
\noindent File \leanpath{VOA/\allowbreak Examples/\allowbreak SymplecticFermionVOA.lean}.
\begin{leancode}
noncomputable def sympFermSNVOA : SuperNearVOA FockSF where
  __ := sympFermSVA
  conformalVector := ωSF
...
noncomputable def sympFermSVOA : SuperVertexOperatorAlgebra FockSF :=
  sympFermSNVOA.toSVOA sympFermSNVOA_vacuum_grade sympFermSNVOA_positivity
...
@[simp] theorem sympFermSVOA_centralCharge : sympFermSVOA.centralCharge = -2 := rfl
\end{leancode}
\begin{dictionary}
\lean{sympFermSVA} is the \lean{SuperVertexAlgebra} of \leanpath{SymplecticFermionInstance.lean}
(grading by word weight, parity \lean{evenOdd QSF}, $Y$ from super reconstruction). The Virasoro field
is \lean{YSF_omega_virasoro}, proved in \leanpath{SymplecticFermionVirasoroA.lean} by an explicit mode
computation; VOA3 is \lean{YSF_translation}. The central charge being $-2$ by \lean{rfl} means it is the
\lean{centralCharge} field, set to $-2$ and \emph{justified} by the Virasoro field, not merely asserted.
\end{dictionary}
\begin{physics}
$T(z)=\omega_{ab}{:}\chi^a\chi^b{:}/2$ up to normalisation; $c=-2$. Spin--statistics fails (odd
states of integer weight), as it must in a non-unitary logarithmic theory.
\end{physics}
\begin{proof}[Proof on paper]
Reconstruction from the creation words; the Virasoro relation is the standard bilinear computation
with a double-contraction anomaly $\frac{m^3-m}{12}\cdot(-2)$, the sign coming from the fermionic
loop.
\end{proof}
\begin{status}
\texttt{\#print axioms VOA.SympFerm.sympFermSVOA}: \lean{[propext, Classical.choice, Quot.sound]}.
Unconditional. It has genuinely odd states (the release audit elaborated \lean{sfCre 0 0 ∈ sympFermSVOA.parity 1}), so the Koszul
sign is exercised; it is the library's only super VOA instance of which this is true.
\end{status}

\begin{theorem}[Symplectic fermions are super $C_2$-cofinite]\label{l17:thm:sfc2}
$\dim\mathrm{Fock}_{SF}/C_2\le4$, spanned by the classes of $\lseventeenvac,\chi^0_{-1}\lseventeenvac,
\chi^1_{-1}\lseventeenvac,\chi^0_{-1}\chi^1_{-1}\lseventeenvac$.
\end{theorem}
\noindent File \leanpath{VOA/\allowbreak Examples/\allowbreak SympFermSuperC2.lean}.
\begin{leancode}
theorem sympFerm_superC2Cofinite : SuperC2Cofinite sympFermSVOA := by
...
theorem finrank_superC2Quotient_le_four :
    Module.finrank ℂ (FockSF ⧸ superO₂ sympFermSVOA) ≤ 4 := by
\end{leancode}
\begin{dictionary}
\lean{superO₂} is Definition~\ref{l17:def:c2}; the hand-built \lean{O₂SF} of
\leanpath{SympFermC2.lean} is \emph{definitionally} the same span (\lean{O₂SF_eq_superO₂ : ... := rfl}).
This is ambient $C_2$-cofiniteness; it does \emph{not} give $C_2$-cofiniteness of the even part
$W(2)$ (whose $C_2$ uses only even vertex operators), as the file warns.
\end{dictionary}
\begin{physics}
$R_{SF}=\Lambda(\chi^0,\chi^1)$: the associated scheme is a fat point with odd directions, so the
theory is lisse, as expected of a $C_2$-cofinite logarithmic CFT.
\end{physics}
\begin{proof}[Proof on paper]
$(\chi^a_{-1}\lseventeenvac)_{(-2)}v=\chi^a_{-2}v$ for every $v$, so every state with a level-$1$ creation
operator in front is in $C_2$; levels $\ge2$ follow by translation, $T$ being an even derivation and
$Tv=v_{(-2)}\lseventeenvac\in C_2$. A word whose leading operator is at level $0$ is handled by
anticommuting it past; level-$0$ generators square to zero, leaving four survivors. (Using
$T\lseventeenV\subseteq C_2$ alone would not terminate: $\lseventeenV/T\lseventeenV$ is
infinite-dimensional here.)
\end{proof}
\begin{status}
\texttt{\#print axioms VOA.SympFerm.sympFerm\_superC2Cofinite}: \lean{[propext, Classical.choice,
Quot.sound]}. Unconditional.
\end{status}

\begin{theorem}[The super Zhu ring of symplectic fermions is not semisimple]\label{l17:thm:sfzhu}
The super Zhu algebra $A(SF)$ (Definition~\ref{l17:def:superzhu}) is an associative ring of dimension
at most $4$, $[\chi^a]\neq0$, and the two-sided ideal generated by $[\chi^0]$ is nonzero and
square-zero; hence $A(SF)$ is not a semisimple ring.
\end{theorem}
\noindent File \leanpath{VOA/\allowbreak Examples/\allowbreak SympFermSuperZhu.lean}.
\begin{leancode}
theorem finrank_ASF_le_four : Module.finrank ℂ ASF ≤ 4 := by
...
theorem chiClass_ne_zero (a : Fin 2) : sfN.superMk hSF (chiStateSF a) ≠ 0 :=
...
theorem ASF_not_isSemisimpleRing :
    letI := instRingASF
    ¬ IsSemisimpleRing ASF :=
\end{leancode}
\begin{dictionary}
\lean{sfN} is \lean{sympFermSVOA.toSuperNearVOA}, \lean{hSF} its integrality proof, \lean{ASF} the
quotient \lean{sfN.superA hSF}, \lean{superMk} the class map. The ring structure \lean{instRingASF}
comes from the discharged \lean{SuperZhuRingData} (Theorem~\ref{l17:thm:superzhuring}).
\lean{IsSemisimpleRing} is Mathlib's notion (every module semisimple). The obstruction used is the
correct one: a nonzero square-zero two-sided ideal; a nonzero nilpotent element would not suffice
($M_2(\C)$ is semisimple). The bundled form is \lean{sympFerm_integralSuperC2FiniteZhuBenchmark}
(\leanpath{SympFermSuperFiniteBenchmark.lean}): super $C_2$, ring data, and the ideal; it is
deliberately not called ``strongly finite'', since simplicity and self-contragredience are not
formalised on the super side.
\end{dictionary}
\begin{physics}
For a rational theory Zhu's algebra is semisimple (its simple modules are the ground states of the
sectors). Non-semisimplicity is the algebraic shadow of the logarithmic, non-split sectors of
symplectic fermions; here $A(SF)\twoheadleftarrow\Lambda(\chi^0,\chi^1)$, the zero-mode Grassmann
algebra.
\end{physics}
\begin{proof}[Proof on paper]
$\dim\le4$: modulo $O(SF)$ one has $\chi^a_{-k-2}v\equiv-\chi^a_{-k-1}v$ (from $(L_{-1}+L_0)\chi^a\ast v\in O$),
which reduces every word to the four bottom words. $[\chi^a]\neq0$: an evaluation homomorphism
$\mathrm{Fock}_{SF}\to\C[\varepsilon]/(\varepsilon^2)$, $\chi^a_{-k-1}\mapsto(-1)^k\delta_{a,c}\varepsilon$,
kills every generator of $O(SF)$ but not $\chi^c$. In the ring, $[\chi^a]^2=0$ and
$[\chi^0][\chi^1]=-[\omega]=-[\chi^1][\chi^0]$, so the ideal generated by $[\chi^0]$ is spanned by
$[\chi^0]$ and $[\chi^0][\chi^1]$ and squares to zero. In a semisimple ring a square-zero two-sided
ideal is zero.
\end{proof}
\begin{status}
\texttt{\#print axioms VOA.SympFerm.ASF\_not\_isSemisimpleRing} and
\lean{VOA.SympFerm.sympFerm\_integralSuperC2FiniteZhuBenchmark}: \lean{[propext, Classical.choice,
Quot.sound]}. Unconditional. The dimension is proved only to be \emph{at most} $4$ (\lean{finrank_ASF_le_four}; also
$\ge1$, \lean{one_le_finrank_ASF}); that $A(SF)$ is four-dimensional is not proved.
\end{status}

\begin{definition}[Super-rationality, library sense]\label{l17:def:superrat}
$W$ is \lean{IsSuperRational} if every \lean{SympFermSuperOrdVModule} of $W$ (a super weak module with
a $\C$-grading by finite-dimensional $L_0$-eigenspaces bounded below) is completely reducible: every
submodule closed under all modes has a complement closed under all modes.
\end{definition}
\noindent File \leanpath{VOA/\allowbreak Examples/\allowbreak SympFermIrrational.lean}.
\begin{leancode}
structure SympFermSuperOrdVModule (W : SuperVertexOperatorAlgebra 𝒱) (M : Type v)
    [AddCommGroup M] [Module ℂ M] extends SuperWeakModule W.toSuperVertexAlgebra M where
  weight : ℂ → Submodule ℂ M
  isInternal : DirectSum.IsInternal weight
  L0_weight : ∀ (α : ℂ) (x : M), x ∈ weight α → mode (YM W.conformalVector) 1 x = α • x
  finrank : ∀ α, FiniteDimensional ℂ (weight α)
  boundedBelow : ∃ B : ℝ, ∀ α : ℂ, weight α ≠ ⊥ → B ≤ α.re
...
def IsSuperRational (W : SuperVertexOperatorAlgebra 𝒱) : Prop :=
  ∀ {M : Type u} [AddCommGroup M] [Module ℂ M] (A : SympFermSuperOrdVModule W M),
    SuperWeakModule.IsCompletelyReducible A.toSuperWeakModule
\end{leancode}
\begin{dictionary}
\lean{IsCompletelyReducible A} $:=\forall P\,\exists Q,\ \mathrm{IsCompl}\,P\,Q$ over mode-stable
submodules. Compared with Dong--Li--Mason rationality (complete reducibility of \emph{admissible},
$\N$-gradable modules) the class is different: graded by $L_0$-eigenvalues, so logarithmic modules
are excluded from the start. The structure is a second copy of \lean{SuperOrdVModule} that
\textbf{omits} the mode-weight axiom \lean{weight_mode}, so it is a \emph{larger} class than
\lean{SuperOrdVModule}; the modules are quantified in universe \lean{u}, the universe of $\lseventeenV$.
The library's own docstring flags that the name should carry the category (\lean{GradedSuperRational}).
\end{dictionary}
\begin{physics}
Rationality = finitely many sectors, all completely reducible = no logarithmic partners. Symplectic
fermions are the textbook irrational $C_2$-cofinite theory: their module with free zero modes is an
indecomposable, non-simple extension.
\end{physics}
\begin{status}
Definitions; standard axioms.
\end{status}

\begin{theorem}[``Symplectic fermions are not super-rational'' --- conditional]\label{l17:thm:sfirrat}
If $N$ is a graded super module of $W$ with a proper nonzero submodule contained in every nonzero
submodule, then $W$ is not super-rational. In particular, given such an $N$ for
\lean{sympFermSVOA}, symplectic fermions are not super-rational. \textbf{This is a socle lemma, not
a proof that symplectic fermions are irrational}: no such $N$ for \lean{sympFermSVOA} is constructed
in the library (F4), and the hypothesis is precisely the missing mathematical content.
\end{theorem}
\noindent File \leanpath{VOA/\allowbreak Examples/\allowbreak SympFermIrrational.lean}.
\begin{leancode}
structure NonSplitSuperModule (W : SuperVertexOperatorAlgebra 𝒱) where
  carrier : Type u
  [addCommGroup : AddCommGroup carrier]
  [module : Module ℂ carrier]
  isMod : SympFermSuperOrdVModule W carrier
  minimal : SuperWeakModule.Submodule' isMod.toSuperWeakModule
  minimal_ne_bot : minimal.carrier ≠ ⊥
  minimal_ne_top : minimal.carrier ≠ ⊤
  minimal_le : ∀ Q : SuperWeakModule.Submodule' isMod.toSuperWeakModule,
    Q.carrier = ⊥ ∨ minimal.carrier ≤ Q.carrier
...
theorem SympFerm.sympFermSVOA_not_isSuperRational
    (N : NonSplitSuperModule SympFerm.sympFermSVOA) : ¬ IsSuperRational SympFerm.sympFermSVOA :=
  not_isSuperRational_of_nonSplit N
\end{leancode}
\begin{dictionary}
The hypothesis \lean{N} is the entire mathematical content: a non-split module. No term of type
\lean{NonSplitSuperModule sympFermSVOA} exists in the library. The library proves, conversely, that the
adjoint module cannot serve (\lean{not_hasFreeZeroModes_of_adjoint}: on $\mathrm{Fock}_{SF}$ the zero
modes $\chi^a_0$ are the zero operator, \lean{zeroMode_adjoin_eq_bot}). So the theorem is a correct
socle lemma (if a module has a minimal proper submodule, it is not completely reducible) and
\emph{not} an irrationality theorem.
\end{dictionary}
\begin{physics}
The missing object is the ``free zero-mode'' module $\Lambda(\chi^+_0,\chi^-_0)\otimes(\text{oscillators})$
of Kausch and Gaberdiel--Kausch, whose socle $\C\chi^+_0\chi^-_0x$ has no complement; on it $L_0$
is diagonalisable but the representation is non-split.
\end{physics}
\begin{proof}[Proof on paper]
If $Q$ complemented the minimal submodule $P$, then either $Q=0$, forcing $P=\top$, or $P\le Q$,
forcing $P=P\cap Q=0$.
\end{proof}
\begin{status}
\texttt{\#print axioms VOA.SympFerm.sympFermSVOA\_not\_isSuperRational}: \lean{[propext,
Classical.choice, Quot.sound]}. \textbf{Conditional} on \lean{N : NonSplitSuperModule sympFermSVOA},
never constructed (and essentially the conclusion; moreover \lean{IsSuperRational} quantifies over
\lean{SympFermSuperOrdVModule}, which lacks \lean{weight_mode}). Irrationality of symplectic fermions is \textbf{open} in Lean and
\textbf{not announceable}. What \emph{is}
unconditional is the algebraic obstruction of Theorem~\ref{l17:thm:sfzhu}; no theorem in the library
converts it into non-rationality.
\end{status}

\begin{warning}[title={Finding (F4): a conditional that assumes its conclusion; stale status block}]
The name \lean{sympFermSVOA_not_isSuperRational} suggests a theorem about symplectic fermions; its
only hypothesis is the non-split module, which is the whole content. In addition the ``Status'' block
of \leanpath{SympFermIrrational.lean} says ``no super module of any kind exists, because
\lean{SuperWeakModule} requires the super Jacobi identity and the library has no super
Borcherds/Jacobi theorem''; this is false at the documented commit (Theorem~\ref{l17:thm:superjacobi},
Proposition~\ref{l17:prop:adjoint}). And \lean{SympFermSuperOrdVModule} duplicates
\lean{SuperOrdVModule} without \lean{weight_mode} (F8). Severities: open gap (conditional presented near an
unconditional name), documentation only (docstring), moderate (duplicate, weaker structure).
\end{warning}

\subsection{The lattice super ambient}\label{l17:sec:lattice}

The doublet and triplet constructions of Chapter~\ref{ch:L18} need an ambient lattice algebra whose
conformal vector is deformed by a background charge, which produces rational (in particular
half-integer) and \emph{negative} weights. \leanpath{LatticeSuperAmbientA.lean} and
\leanpath{LatticeSuperAmbientB.lean} build it as a purely even \lean{SuperNearVOA}.

\begin{theorem}[Lattice super near-VOA with background charge]\label{l17:thm:lattice}
For an even positive-definite lattice $L$, a $\Q$-valued quadratic refinement $q$ with finite level
sets, and a background charge $\lambda$ with $\langle\lambda,\beta\rangle=\frac12\langle\beta,\beta\rangle-q(\beta)$
for all $\beta\in L$, the lattice vertex algebra with grading by $q$-refined weight and conformal vector
$\omega_\lambda=\omega_L+\lambda_{(-2)}\lseventeenvac$ is a super near-VOA, bounded below, in general not
positive.
\end{theorem}
\noindent File \leanpath{VOA/\allowbreak Examples/\allowbreak LatticeSuperAmbientB.lean}.
\begin{leancode}
def BgCompat (qr : QuadRefine L) (lam : Fin (rank L) → ℂ) : Prop :=
  ∀ β : L.carrier, BgB.bgPair L lam β = ((halfNorm L β : ℤ) : ℂ) - ((qr.q β : ℚ) : ℂ)
...
noncomputable def latSNVOA (hEven : L.Even) (hPos : L.PosDef) (qr : QuadRefine L)
    (hqf : QFinite L qr) (lam : Fin (rank L) → ℂ) (hcompat : BgCompat L qr lam) :
    SuperNearVOA (LatModuleGen L) where
\end{leancode}
\begin{dictionary}
The vertex operators, vacuum, creation, locality and regularity are the bosonic lattice ones of
Chapter~\ref{ch:03}; only the grading (by $q$) and the conformal vector change. Parity is trivial.
The new obligation VA1 for the refined grading holds because the refinement differs from the
undeformed weight by $-\langle\lambda,\text{charge}\rangle$, additive in the charge.
\lean{deformedAmbient_not_positivity} shows positivity genuinely fails in the doublet case.
\end{dictionary}
\begin{physics}
A Coulomb gas: a free boson with background charge, $c=\operatorname{rank}L-12\langle\lambda,\lambda\rangle$;
vertex operators $e^{\beta}$ acquire weights $\frac12\beta^2-\langle\lambda,\beta\rangle$, which can be
fractional or negative. Screening operators inside it cut out the triplet and doublet algebras.
\end{physics}
\begin{status}
\texttt{\#print axioms VOA.LatGen.SAmbB.latSNVOA}: \lean{[propext, Classical.choice, Quot.sound]}. Conditional on the displayed
lattice data, which are instantiated in the files (\lean{tripletSNVOA}, \lean{Half.halfSuperNearVOA}).
\end{status}

%======================================================================
\section{Super Zhu algebras, untwisted and twisted}\label{l17:sec:zhu}
%======================================================================

\begin{definition}[Super Zhu algebra]\label{l17:def:superzhu}
Let $V$ be a super near-VOA with integer weights. For $u\in V_k$ put
$u\ast v=\sum_{m=0}^{k}\binom km u_{(m-1)}v$, extended linearly; $O(V)=\operatorname{span}\{(L_{-1}u+L_0u)\ast v\}$;
$A(V)=V/O(V)$.
\end{definition}
\noindent File \leanpath{VOA/\allowbreak Modules/\allowbreak SuperZhuRing.lean}.
\begin{leancode}
noncomputable def superZhuFamily (v : 𝒱) (k : ℤ) : (V.gradingZ k) →ₗ[ℂ] 𝒱 where
  toFun uk := ∑ m ∈ Finset.range (k.toNat + 1), cbz k m • V.smul (uk : 𝒱) ((m : ℤ) - 1) v
...
def superO (hV : V.IntegralWeights) : Submodule ℂ 𝒱 :=
  Submodule.span ℂ {x | ∃ u v : 𝒱, x = V.superZhuMul hV (V.L (-1) u + V.L 0 u) v}
...
abbrev superA (hV : V.IntegralWeights) : Type u := 𝒱 ⧸ V.superO hV
\end{leancode}
\begin{dictionary}
\lean{V.IntegralWeights} says $V_h=0$ for $h\notin\Z$; \lean{gradingZ k} is $V_k$. \lean{superZhuMul}
applies \lean{superZhuFamily} to the homogeneous components. The generator
$(L_{-1}u+L_0u)\ast v$ for $u\in V_k$ equals Zhu's $u\circ v=\Res_z Y(u,z)v\frac{(1+z)^{k}}{z^2}$ (the
Pascal identity is checked by hand in a docstring, not formalised). \textbf{No parity case split}: for
an integrally graded super VOA this is the Zhu algebra relevant to untwisted (``Neveu--Schwarz'' in the
integral sense) modules, as in Abe's treatment of symplectic fermions; Kac--Wang and Dong--Zhao's
untwisted $A(V)$ for a $\frac12\Z$-graded VOSA (with $u\circ v=\Res_z Y(u,z)v\,(1+z)^{\wt u-1/2}/z$ for odd $u$)
is \emph{not} covered: the library requires integer weights.
\end{dictionary}
\begin{physics}
Ground states of a sector form a module over $A(V)$; the zero mode $o(u)=u_{(\wt u-1)}$ of each field
acts on them, and $O(V)$ is the set of combinations whose zero modes vanish on every ground state.
\end{physics}
\begin{status}
Definitions; standard axioms.
\end{status}

\begin{theorem}[$A(V)$ is an associative ring; Dong--Zhao Thm.~3.1]\label{l17:thm:superzhuring}
For every integrally graded super VOA $W$, $O(W)$ is a two-sided ideal for $\ast$ and $\ast$ is
associative modulo $O(W)$; so $A(W)$ is an associative unital $\C$-algebra. If $W$ is super
$C_2$-cofinite, $A(W)$ is finite-dimensional.
\end{theorem}
\noindent File \leanpath{VOA/\allowbreak Modules/\allowbreak SuperZhuRingDischarge.lean}.
\begin{leancode}
theorem superZhuRingData (W : SuperVertexOperatorAlgebra 𝒱)
    (hV : W.toSuperNearVOA.IntegralWeights) :
    SuperZhuRingData W.toSuperNearVOA hV :=
\end{leancode}
\noindent File \leanpath{VOA/\allowbreak Modularity/\allowbreak SuperZhuAlgebraFinite.lean}.
\begin{leancode}
theorem finiteDimensional_superA_of_superC2Cofinite (W : SuperVertexOperatorAlgebra 𝒱)
    (hV : W.toSuperNearVOA.IntegralWeights) (hC2 : SuperC2Cofinite W) :
    FiniteDimensional ℂ (W.toSuperNearVOA.superA hV) := by
\end{leancode}
\begin{dictionary}
\lean{SuperZhuRingData V hV} is the proposition with fields \lean{left_ideal}, \lean{right_ideal},
\lean{assoc} (\leanpath{Modules/\allowbreak SuperZhuRing.lean}); from it the ring and algebra instances
are built (\lean{instRingSuperA}). The general near-VOA version \lean{superZhuRingData_holds} needs
only integrality and vanishing of negative weight spaces. No spin--statistics, no skew gate.
\end{dictionary}
\begin{physics}
Zhu's theorem: $A(V)$ is the algebra of zero modes on ground states; it is associative because the
zero modes compose. The finite dimensionality reflects finitely many ground-state types in a lisse
theory.
\end{physics}
\begin{proof}[Proof on paper]
Zhu's residue calculus, run with the super Jacobi identity: the right-ideal and associator identities
are residue manipulations in which every exchange of two vectors happens through the Jacobi identity,
so the Koszul sign appears and cancels consistently; the left-ideal property follows from
associativity. Finite-dimensionality: by triangularity in the weight, every $u_{(-2)}v$ equals
$(L_{-1}u+L_0u)\ast v$ plus lower-weight terms, so a finite spanning set of $W/C_2$ spans $A(W)$.
\end{proof}
\begin{status}
\texttt{\#print axioms} of \lean{VOA.SuperVertexOperatorAlgebra.superZhuRingData} and
\lean{VOA.SuperNearVOA.finiteDimensional\_superA\_of\_superC2Cofinite}: \lean{[propext,
Classical.choice, Quot.sound]}. Unconditional beyond \lean{hV} (and \lean{hC2}).
\end{status}

\begin{definition}[Twisted super Zhu quotient]\label{l17:def:twistedzhu}
Let $g$ be an even automorphism with $g^N=1$, $\delta_r=1$ if $r=0$ and $0$ otherwise. For $u\in W_k$
in the eigenspace $W^r$ ($gu=e^{-2\pi ir/N}u$) put
$u\circ_gv=\sum_{i\ge0}\binom{k-1+\delta_r+r/N}{i}u_{(i-1-\delta_r)}v$. Let $\bar u=\frac1N\sum_{a<N}g^au$
(Reynolds average). $O_g(W)$ is spanned by the fixed-sector generators $(L_{-1}\bar u+L_0\bar u)\ast v$
and the nonfixed residues $u\circ_gv$ ($u\in W_k\cap W^r$, $r\neq0$); $A_g(W)=W/O_g(W)$.
\end{definition}
\noindent File \leanpath{VOA/\allowbreak Modules/\allowbreak SuperTwistedZhu.lean}.
\begin{leancode}
noncomputable def twistedZhuCircleHom (W : SuperVertexOperatorAlgebra 𝒱)
    (N r : ℕ) (k : ℤ) (u v : 𝒱) : 𝒱 :=
  ∑ᶠ i : ℕ, fractionalBinom (twistedZhuExponent N r k) i •
    W.smul u ((i : ℤ) - 1 - (twistedZhuDelta r : ℤ)) v
...
def SuperVertexOperatorAlgebra.twistedO
    (W : SuperVertexOperatorAlgebra 𝒱) (hW : W.toSuperNearVOA.IntegralWeights)
    (g : SuperAut W) (N : ℕ) : Submodule ℂ 𝒱 :=
  Submodule.span ℂ (W.twistedZhuFixedGeneratorSet hW g N ∪
    W.twistedZhuNonfixedGeneratorSet g N)
...
abbrev SuperVertexOperatorAlgebra.twistedA
    (W : SuperVertexOperatorAlgebra 𝒱) (hW : W.toSuperNearVOA.IntegralWeights)
    (g : SuperAut W) (N : ℕ) : Type u := 𝒱 ⧸ W.twistedO hW g N
\end{leancode}
\begin{dictionary}
\lean{twistedZhuExponent N r k} $=k-1+\delta_r+r/N$; in the fixed sector the circle has exponent $k$
and pole order $2$ (\lean{twistedZhuCircleHom_fixed}). At the identity twist $O_1=O$ exactly
(\lean{twistedO_identity}). This matches Dong--Zhao (math/0411524), §3, in the integral-weight case;
their sign-adjusted convention ($\sigma g$ in place of $g$) is kept separate in the library
(\lean{SuperAut.twistedZhuSignAdjusted}). The ring axioms for $g\neq1$ are the unproved proposition
\lean{SuperTwistedZhuRingData} (including the right-unit statement $u-\bar u\in O_g$); so for $g\ne1$
$A_g$ is only a vector space in Lean.
\end{dictionary}
\begin{physics}
Ground states of $g$-twisted sectors are modules for $A_g(V)$; the nonfixed fields have no zero modes
in a $g$-twisted sector (their modes are fractional), which is why $W^r$, $r\neq0$, is killed.
In the library the quotient $A_g(W)$ is a genuine object, but the map from twisted modules to it
(\lean{SuperTwistedModule.twistedZeroModeAction}) takes a \lean{SuperTwistedModule}, which for
$N\ge2$ is the zero space (F1); so no nonzero ground-state module of a twisted sector is ever
produced.
\end{physics}
\begin{status}
Definitions; standard axioms.
\end{status}

\begin{theorem}[Twisted Zhu finiteness]\label{l17:thm:twistedzhufinite}
If $W$ is integrally graded and super $C_2$-cofinite, $g^N=1$ with $N>0$, and the eigenspace
decomposition of $g$ is internal and compatible with weight and parity, then $A_g(W)$ is
finite-dimensional.
\end{theorem}
\noindent File \leanpath{VOA/\allowbreak Modules/\allowbreak SuperTwistedZhuFinite.lean}.
\begin{leancode}
theorem finiteDimensional_twistedA_of_superC2Cofinite
    (W : SuperVertexOperatorAlgebra 𝒱) (hW : W.toSuperNearVOA.IntegralWeights)
    (g : SuperAut W) {N : ℕ} (hN : 0 < N) (_hpow : g ^ N = 1)
    (hE : TwistEigenspaceDecomposition W g N) (hC2 : SuperC2Cofinite W) :
    FiniteDimensional ℂ (W.twistedA hW g N) := by
\end{leancode}
\begin{dictionary}
\lean{TwistEigenspaceDecomposition W g N} has three fields: \lean{isInternal} (the eigenspaces
$W^r$, $r\in\Z/N$, form an internal direct sum, i.e.\ $g$ is diagonalisable) and the two statements that
the projections preserve each $W_k$ and each parity. Over $\C$ a finite-order operator is
diagonalisable, and on each finite-dimensional $W_k$ the projections are polynomials in $g$; so this
hypothesis is in principle derivable from $g^N=1$, but the library supplies it as data. \lean{_hpow}
is recorded but not used. The conclusion concerns a vector space (no ring structure needed). This
theorem does \emph{not} involve \lean{SuperTwistedModule}, so finding (F1) does not touch it: it is a
true statement about the vector space $A_g(W)$. What (F1) does touch is any intended application to
twisted modules: there is no nonzero twisted module with $N\ge2$ in the library for $A_g(W)$ to act on.
\end{dictionary}
\begin{physics}
A lisse theory has finitely many ground-state types in each twisted sector too.
\end{physics}
\begin{proof}[Proof on paper]
Triangular reduction by weight within each eigen-sector: for $u\in W^r_k$, the leading term of the
circle residue is $u_{(-2)}v$ (fixed sector) or $u_{(-1)}v$ plus a pole-one tail (nonfixed), which
reduces $C_2$ modulo $O_g$ and strictly lower weights; a finite spanning set of $W/C_2$ then spans
$A_g$.
\end{proof}
\begin{status}
\texttt{\#print axioms VOA.finiteDimensional\_twistedA\_of\_superC2Cofinite}: \lean{[propext,
Classical.choice, Quot.sound]}. Conditional on \lean{hE} (diagonalisability data), \lean{hW},
\lean{hC2}.
\end{status}

%======================================================================
\section{The modular toolkit for super and half-integral theories}\label{l17:sec:modularity}
%======================================================================

For a bosonic rational $C_2$-cofinite VOA, Zhu's theorem (Chapter~\ref{ch:L15}) says that the torus
one-point functions $\tr_Mo(v)q^{L_0-c/24}$ span an $SL_2(\Z)$-invariant space. With fermions or
half-integer weights two things change. First, $T:\tau\mapsto\tau+1$ multiplies a $q^{h}$-term by
$e^{2\pi ih}=-1$ when $h\in\frac12+\Z$, so only $T^2$ is a symmetry of a single character and the
natural group is the theta group $\Gamma_\theta=\langle S,T^2\rangle$. Second, a fermion on a torus
needs a spin structure: a choice of periodic or antiperiodic boundary condition along each of the two
cycles. The four spin structures $(\mu,\lambda)\in\{\pm\}^2$ correspond to the four traces
$\tr_{NS}$, $\tr_{NS}(-1)^F$, $\tr_R$, $\tr_R(-1)^F$ and are permuted by $SL_2(\Z)$ through its
action on $(\Z/2)^2$; $\Gamma_\theta$ is the stabiliser of the NS--NS structure. This section presents
the group theory, the automorphy factor, the four ``spin'' Eisenstein-type coefficient families of
van Ekeren, and the proved transformation statements.

\begin{definition}[The theta group]\label{l17:def:gammatheta}
$\Gamma_\theta=\{g\in SL_2(\Z): g\equiv\begin{psmallmatrix}1&0\\0&1\end{psmallmatrix}\text{ or }
g\equiv\begin{psmallmatrix}0&1\\1&0\end{psmallmatrix}\pmod 2\}$.
\end{definition}
\noindent File \leanpath{VOA/\allowbreak Modularity/\allowbreak ThetaGroup.lean}.
\begin{leancode}
def redEntry (g : SL(2, ℤ)) (i j : Fin 2) : ZMod 2 :=
  (((g : Matrix (Fin 2) (Fin 2) ℤ) i j : ℤ) : ZMod 2)
...
def IsThetaCongr (g : SL(2, ℤ)) : Prop :=
  (redEntry g 0 0 = 1 ∧ redEntry g 0 1 = 0 ∧ redEntry g 1 0 = 0 ∧ redEntry g 1 1 = 1) ∨
    (redEntry g 0 0 = 0 ∧ redEntry g 0 1 = 1 ∧ redEntry g 1 0 = 1 ∧ redEntry g 1 1 = 0)
def GammaTheta : Subgroup SL(2, ℤ) where
  carrier := {g | IsThetaCongr g}
\end{leancode}
\begin{dictionary}
Mod $2$, $S=\begin{psmallmatrix}0&-1\\1&0\end{psmallmatrix}\equiv\begin{psmallmatrix}0&1\\1&0\end{psmallmatrix}$,
so this is ``$g\equiv I$ or $g\equiv S$''. Closure under products and inverses is proved by computation
in $\Z/2$. Standard (e.g.\ Knopp, \emph{Modular functions in analytic number theory}, Ch.~2).
\end{dictionary}
\begin{physics}
$\Gamma_\theta$ is the subgroup of the mapping class group of the torus preserving the NS--NS
(antiperiodic--antiperiodic) spin structure, i.e.\ the modular group of a single fermionic partition
function $\tr_{NS}q^{L_0-c/24}$.
\end{physics}
\begin{status}
A definition; standard axioms.
\end{status}

\begin{theorem}[Generation and index]\label{l17:thm:gammatheta}
$\Gamma_\theta=\langle S,T^2\rangle$, $T\notin\Gamma_\theta$, $\Gamma(2)\le\Gamma_\theta$, and
$[SL_2(\Z):\Gamma_\theta]=3$.
\end{theorem}
\noindent File \leanpath{VOA/\allowbreak Modularity/\allowbreak ThetaGroup.lean}.
\begin{leancode}
theorem Gamma_two_le : CongruenceSubgroup.Gamma 2 ≤ GammaTheta := by
...
theorem T_not_mem : ModularGroup.T ∉ GammaTheta := by
...
def genSTsq : Subgroup SL(2, ℤ) :=
  Subgroup.closure ({ModularGroup.S, ModularGroup.T ^ 2} : Set SL(2, ℤ))
...
theorem GammaTheta_eq_genSTsq : GammaTheta = genSTsq := by
...
theorem index_eq_three : GammaTheta.index = 3 := by
\end{leancode}
\begin{dictionary}
\lean{Subgroup.closure} is the subgroup generated. \lean{index} is Mathlib's \lean{Subgroup.index}
(the cardinality of $SL_2(\Z)/\Gamma_\theta$). \lean{genSTsq} is definitionally the \lean{thetaGroup}
of \leanpath{Modularity/\allowbreak SlashSubgroup.lean} (Chapter~\ref{ch:L14}), and
\lean{GammaThetaStmt.thetaGroup_eq} records \lean{SlashSub.thetaGroup = ThetaGroup.GammaTheta}.
\end{dictionary}
\begin{physics}
The three cosets $\Gamma_\theta$, $T^{-1}\Gamma_\theta$, $(TS)^{-1}\Gamma_\theta$ correspond to the
three even spin structures (NS--NS, NS--R, R--NS), permuted by $SL_2(\Z)$; the odd R--R structure is
fixed.
\end{physics}
\begin{proof}[Proof on paper]
$\langle S,T^2\rangle\le\Gamma_\theta$ since both generators are in it. Conversely, $SL_2(\Z)$ acts on
the three right cosets of $\langle S,T^2\rangle$ represented by $1,T,TS$ (checked on the generators
$S,T$); an element of $\Gamma_\theta$ maps the coset of $1$ to that of $t\in\{1,T,TS\}$ with
$t\in\Gamma_\theta$, and $T,TS\notin\Gamma_\theta$ force $t=1$. Index: every $g$ is in exactly one of
three explicit cosets, by its reduction mod $2$ ($SL_2(\Z/2)\cong S_3$ has order $6$ and the image of
$\Gamma_\theta$ has order $2$).
\end{proof}
\begin{status}
\texttt{\#print axioms} of \lean{VOA.ThetaGroup.GammaTheta\_eq\_genSTsq},
\lean{VOA.ThetaGroup.index\_eq\_three}: \lean{[propext, Classical.choice, Quot.sound]}. Unconditional.
\end{status}

\begin{theorem}[A half-integral automorphy factor on $\Gamma_\theta$]\label{l17:thm:halfaut}
The principal square root $s(A,\tau)=\sqrt{c\tau+d}$ satisfies, for all $A,B\in SL_2(\Z)$ and $\tau$ in
the upper half-plane, $s(AB,\tau)=m(A,B)\,s(A,B\tau)\,s(B,\tau)$ with a constant $m(A,B)$ independent
of $\tau$; hence a \lean{HalfAutFactor} exists on every subgroup, in particular on $\Gamma_\theta$.
\end{theorem}
\noindent File \leanpath{VOA/\allowbreak Modularity/\allowbreak SlashSubgroup.lean}.
\begin{leancode}
structure HalfAutFactor (Γ : Subgroup SL(2, ℤ)) where
  sq : SL(2, ℤ) → UpperHalfPlane → ℂ
  sq_sq : ∀ A ∈ Γ, ∀ τ : UpperHalfPlane, sq A τ ^ (2 : ℤ) = autFactor A τ
  sq_ne_zero : ∀ A ∈ Γ, ∀ τ : UpperHalfPlane, sq A τ ≠ 0
  mult : SL(2, ℤ) → SL(2, ℤ) → ℂ
  cocycle : ∀ A ∈ Γ, ∀ B ∈ Γ, ∀ τ : UpperHalfPlane,
    sq (A * B) τ = mult A B * (sq A (B • τ) * sq B τ)
  sq_one : ∀ τ : UpperHalfPlane, sq 1 τ = 1
\end{leancode}
\noindent File \leanpath{VOA/\allowbreak Modularity/\allowbreak ThetaHalfAutFactor.lean}.
\begin{leancode}
noncomputable def principalHalfRoot (A : SL(2, ℤ)) (τ : UpperHalfPlane) : ℂ :=
  Complex.sqrt (autFactor A τ)
...
noncomputable def thetaHalfAutFactor : HalfAutFactor thetaGroup :=
  principalHalfAutFactor thetaGroup
...
theorem thetaHalfAutFactorExists : ThetaHalfAutFactorExists :=
  ⟨SlashSub.thetaHalfAutFactor⟩
\end{leancode}
\begin{dictionary}
\lean{autFactor A τ} $=c\tau+d$ (Chapter~\ref{ch:L14}). The crucial point is that \lean{mult} does
\emph{not} depend on $\tau$; that is a real constraint, proved by continuity of the ratio, the fact that
its square is $1$, and connectedness of the upper half-plane (evaluate at $\tau=i$). The resulting
multiplier is $\pm1$-valued; the docstring's ``intended'' theta multiplier (an eighth root of unity) is
a different, equally valid choice, the difference being absorbable into the representation $\rho$ of a
modularity statement. \lean{ThetaHalfAutFactorExists} is
\lean{Nonempty (SlashSub.HalfAutFactor SlashSub.thetaGroup)}.
\end{dictionary}
\begin{physics}
Characters of weight $\frac12$ objects ($\eta$, $\theta$-functions, fermionic partition functions with
insertions of half-integral weight) transform with $(c\tau+d)^{1/2}$ times a multiplier system; this is
the analytic datum needed to even state $\Gamma_\theta$-modularity.
\end{physics}
\begin{proof}[Proof on paper]
$s(AB,\tau)^2=j(AB,\tau)=j(A,B\tau)j(B,\tau)=(s(A,B\tau)s(B,\tau))^2$ ($j$ the automorphy factor, a
cocycle), so the ratio $m$ is a continuous function on the upper half-plane with values in $\{\pm1\}$
(the denominator never vanishes), hence constant.
\end{proof}
\begin{status}
\texttt{\#print axioms VOA.GammaThetaStmt.thetaHalfAutFactorExists}: \lean{[propext,
Classical.choice, Quot.sound]} (umbrella check). Unconditional.
\end{status}

\begin{definition}[Spin pairs and their $SL_2(\Z)$-action]\label{l17:def:spinpair}
A spin pair is $(\mu,\lambda)\in(\Z/2)^2$ (bit $0$ for $+1$, $1$ for $-1$); $A=\begin{psmallmatrix}a&b\\c&d\end{psmallmatrix}$
acts on the right by $(\mu,\lambda)\cdot A=(a\mu+c\lambda,\,b\mu+d\lambda)$ mod $2$.
\end{definition}
\noindent File \leanpath{VOA/\allowbreak Modularity/\allowbreak SuperSpinCoefficients.lean}.
\begin{leancode}
structure SpinPair where
  mu : SpinPhase
  lambda : SpinPhase
deriving DecidableEq, Fintype
...
def act (A : SL(2, ℤ)) (p : SpinPair) : SpinPair :=
  ⟨p.mu * ThetaGroup.redEntry A 0 0 + p.lambda * ThetaGroup.redEntry A 1 0,
   p.mu * ThetaGroup.redEntry A 0 1 + p.lambda * ThetaGroup.redEntry A 1 1⟩
\end{leancode}
\begin{dictionary}
\lean{SpinPhase} is \lean{ZMod 2}. Proved: a right action (\lean{act_mul}), $S$ swaps the labels,
$T$ sends $(\mu,\lambda)\mapsto(\mu,\mu+\lambda)$, $T^2$ fixes everything, and every element of
$\Gamma_\theta$ either fixes or swaps (\lean{act_gammaTheta_self_or_swap}). The labels are \emph{not}
identified with parity: \lean{mu} and \lean{lambda} are spatial and temporal eigenvalue labels.
\end{dictionary}
\begin{physics}
$\mu$ labels the boundary condition along the spatial circle (the twist $g$ of the sector), $\lambda$
along the time circle (the insertion $h$ in the trace). The action is the orbifold rule
$(g,h)\cdot A=(g^ah^c,g^bh^d)$ of Dong--Li--Mason read on $\pm1$ eigenvalues. For fermions the physical
spin structure combines these labels with the conformal sign $\kappa=e^{2\pi iL_0}$ (\leanpath{SuperConformalSign.lean}),
which is why $T$ fixes the label $(+,+)$ and yet $T\notin\Gamma_\theta$: on half-integral weights $T$
itself produces a sign.
\end{physics}
\begin{status}
Definitions; standard axioms.
\end{status}

\begin{definition}[Characteristic theta functions]\label{l17:def:thetachar}
$\Theta[a,b](z,\tau)=e^{\pi ia^2\tau+2\pi ia(z+b)}\,\vartheta(z+a\tau+b,\tau)$ with
$\vartheta(z,\tau)=\sum_ne^{\pi in^2\tau+2\pi inz}$ (Mathlib's \lean{jacobiTheta₂}); for a spin pair
the characteristics are $a,b\in\{0,\frac12\}$.
\end{definition}
\noindent File \leanpath{VOA/\allowbreak Modularity/\allowbreak SuperSpinThetaBridge.lean}.
\begin{leancode}
noncomputable def spinThetaChar (a b z tau : Complex) : Complex :=
  Complex.exp (Real.pi * Complex.I * a ^ 2 * tau +
    2 * Real.pi * Complex.I * a * (z + b)) *
      jacobiTheta₂ (z + a * tau + b) tau
...
noncomputable def spinPairTheta (p : SpinPair) (z tau : Complex) : Complex :=
  spinThetaChar (spinThetaHalf p.mu) (spinThetaHalf p.lambda) z tau
\end{leancode}
\begin{dictionary}
\lean{spinThetaHalf} sends $+\mapsto0$, $-\mapsto\frac12$. This is the standard theta function with
characteristics (Mumford, \emph{Tata lectures on theta I}, §II.1).
\end{dictionary}
\begin{physics}
$\vartheta[a,b](0,\tau)/\eta(\tau)$ are the partition functions of a complex fermion in the four
spin structures; their ratios and logarithmic derivatives are the propagators of a fermion on the
torus, which is how they enter the genus-one recursion.
\end{physics}
\begin{status}
Definitions; standard axioms.
\end{status}

\begin{theorem}[Spin theta $S$-law]\label{l17:thm:thetaS}
For all $a,b,z\in\C$ and $\tau\neq0$,
\begin{equation*}
\Theta[a,b]\Bigl(\frac z\tau,-\frac1\tau\Bigr)=(-i\tau)^{1/2}\,e^{\pi iz^2/\tau+2\pi iab}\,\Theta[b,-a](z,\tau),
\end{equation*}
and, for $\operatorname{Im}\tau>0$, the cross-multiplied difference of logarithmic derivatives of two
such theta functions transforms homogeneously with the factor $((-i\tau)^{1/2})^2$ and the product of
the two Gaussian-and-phase exponentials.
\end{theorem}
\noindent File \leanpath{VOA/\allowbreak Modularity/\allowbreak SuperSpinThetaBridge.lean}.
\begin{leancode}
theorem spinThetaChar_S (a b z tau : Complex) (htau : tau ≠ 0) :
    spinThetaChar a b (z / tau) (-1 / tau) =
      (-Complex.I * tau) ^ (1 / 2 : Complex) *
        Complex.exp (Real.pi * Complex.I * z ^ 2 / tau +
          2 * Real.pi * Complex.I * a * b) * spinThetaChar b (-a) z tau := by
...
theorem spinThetaCross_S
    (a b c d z tau : Complex) (htauIm : 0 < tau.im) :
    thetaCross
        (fun w => spinThetaChar a b (w / tau) (-1 / tau))
        (fun w => spinThetaChar c d (w / tau) (-1 / tau)) z =
      ((-Complex.I * tau) ^ (1 / 2 : Complex)) ^ 2 *
        Complex.exp (2 * (Real.pi * Complex.I * z ^ 2 / tau) +
          2 * Real.pi * Complex.I * (a * b + c * d)) *
        thetaCross (fun w => spinThetaChar b (-a) w tau)
          (fun w => spinThetaChar d (-c) w tau) z := by
\end{leancode}
\begin{dictionary}
\lean{(-Complex.I * tau) ^ (1 / 2 : Complex)} is the principal power $(-i\tau)^{1/2}$.
\lean{thetaCross f g z} $=f'(z)g(z)-g'(z)f(z)$, so $\mathrm{thetaCross}=0$ says the logarithmic
derivatives agree; the second statement is the denominator-free form of the law for
$\partial_z\log\Theta[a,b]-\partial_z\log\Theta[c,d]$. The characteristic $[b,-a]$ is unreduced
($-a$ rather than $-a$ mod $1$); reduction only changes a constant sign. \textbf{This is the analytic
kernel only}: the identification of the coefficient families $P^{\mu,\lambda}_k$ below with Taylor
coefficients of these logarithmic derivatives (a Jacobi triple product / Lambert expansion) is
\emph{not} in the library, so no $S$-law for the non-$(+,+)$ spin coefficients follows.
\end{dictionary}
\begin{physics}
Poisson summation: $S$ exchanges space and time, hence the two characteristics, i.e.\ the NS--R and
R--NS spin structures, with the modular anomaly $(-i\tau)^{1/2}$ of a weight-$\frac12$ object.
\end{physics}
\begin{proof}[Proof on paper]
Mathlib's \lean{jacobiTheta₂_functional_equation} is the case $a=b=0$; substitute
$w=z-a+b\tau$ and collect the Gaussian exponentials (a polynomial identity in the exponents). For the
logarithmic-derivative form, the $z$-independent phase and the common Gaussian cancel in
$f'g-g'f$.
\end{proof}
\begin{status}
\texttt{\#print axioms} of \lean{VOA.spinThetaChar\_S}, \lean{VOA.spinThetaCross\_S}: \lean{[propext,
Classical.choice, Quot.sound]}. Unconditional.
\end{status}

\begin{definition}[Van Ekeren's spin coefficient series]\label{l17:def:spinP}
For a spin pair $p=(\mu,\lambda)$ and $k\in\N$, $P^{\mu,\lambda}_k$ is the formal series on the half-lattice
$\frac12\N$ of exponents with constant term $\delta_{k,0}\delta_{p,(+,-)}\pi i-\frac{(2\pi i)^{k+1}}{(k+1)!}B_{k+1}(1+\varepsilon_\mu)$
($\varepsilon_+=0$, $\varepsilon_-=-\frac12$) and higher coefficients the Lambert sums of van Ekeren's
Lemma 3.1; $Q^{\mu,\lambda}_k=(2\pi i)^{-k-1}P^{\mu,\lambda}_k$, and $\mathcal Q^{\mu,\lambda}_k(\tau)=\sum_NQ_k(N/2)e^{\pi iN\tau}$.
\end{definition}
\noindent File \leanpath{VOA/\allowbreak Modularity/\allowbreak SuperSpinCoefficients.lean}.
\begin{leancode}
noncomputable def spinP (p : SpinPair) (k : ℕ) : FormalQSeries :=
  halfLatticeQSeries (spinPNatCoeff p k)
\end{leancode}
\noindent File \leanpath{VOA/\allowbreak Modularity/\allowbreak SuperSpinRelations.lean}.
\begin{leancode}
noncomputable def spinQ (p : SpinPair) (k : Nat) : FormalQSeries :=
  (spinTwoPiI ^ (k + 1))⁻¹ • spinP p k
\end{leancode}
\noindent File \leanpath{VOA/\allowbreak Modularity/\allowbreak SuperSpinAnalytic.lean}.
\begin{leancode}
noncomputable def spinQAnalyticTerm (p : SpinPair) (k N : Nat) (z : Complex) : Complex :=
  spinQ p k (halfQExponent N) *
    Complex.exp (spinTwoPiI * z * halfQExponent N)
noncomputable def spinQAnalytic (p : SpinPair) (k : Nat) (z : Complex) : Complex :=
  ∑' N : Nat, spinQAnalyticTerm p k N z
\end{leancode}
\begin{dictionary}
\lean{FormalQSeries} is $\C\to\C$ (coefficient at a complex exponent, Chapter~\ref{ch:L14});
\lean{halfLatticeQSeries a} places $a(N)$ at exponent $N/2$ and $0$ elsewhere; \lean{spinPNatCoeff}
is the explicit cusp-plus-Lambert coefficient (quoted in full in the source; it uses
\lean{Polynomial.bernoulli} and divisor sums over \lean{Nat.divisorsAntidiagonal}). \lean{∑'} is
Mathlib's \lean{tsum} (value $0$ if not summable; summability is proved,
\lean{summable_spinQAnalyticTerm}, for $\operatorname{Im}z>0$). Proved fidelity check: at $(+,+)$ and
$k=2j-1$, $Q$ is \emph{exactly} Zhu's normalised Eisenstein series \lean{zhuG j}
(Chapter~\ref{ch:L14}, Definition~\ref{l14:def:zhuG}; \lean{spinP_pp_eq_spinTwoPiI_pow_zhuG}); for
$\mu=+$ the odd half-exponents vanish.
\end{dictionary}
\begin{physics}
These are the coefficient functions of the twisted Weierstrass functions $P_k\bigl[{\mu\atop\lambda}\bigr](z,\tau)$
that appear in genus-one recursion for traces over twisted sectors with twisted insertions (Mason--Tuite--Zuevsky,
van Ekeren). In the untwisted sector they are the Eisenstein series $G_{2k}$. (The coefficient
functions are genuine; but the twisted-sector traces they are meant to enter are, in the library,
traces over \lean{SuperTwistedModule}s with $N\ge2$, which are zero spaces (F1).)
\end{physics}
\begin{status}
Definitions; standard axioms. Holomorphy (\lean{differentiableOn_spinQAnalytic}) and uniform convergence
on closed sub-half-planes are proved.
\end{status}

\begin{theorem}[Proved transformation cells]\label{l17:thm:spinS}
For every spin pair, $\mathcal Q^{\mu,\lambda}_k(\tau+2)=\mathcal Q^{\mu,\lambda}_k(\tau)$. In the invariant
sector $(+,+)$: for $j\ge2$, $\mathcal Q_{2j-1}(-1/\tau)=\tau^{2j}\mathcal Q_{2j-1}(\tau)$, and
$\mathcal Q_1(-1/\tau)=\tau^2\mathcal Q_1(\tau)-\tau/(2\pi i)$.
\end{theorem}
\noindent File \leanpath{VOA/\allowbreak Modularity/\allowbreak SuperSpinAnalytic.lean}.
\begin{leancode}
theorem spinQAnalytic_add_two (p : SpinPair) (k : Nat) (z : Complex) :
    spinQAnalytic p k (z + 2) = spinQAnalytic p k z := by
\end{leancode}
\noindent File \leanpath{VOA/\allowbreak Modularity/\allowbreak SuperSpinS.lean}.
\begin{leancode}
theorem spinQAnalytic_pp_higher_S {j : Nat} (hj : 2 ≤ j) {z : Complex}
    (hz : z ∈ imGT 0) :
    spinQAnalytic SpinPair.pp (2 * j - 1) (-1 / z) =
      z ^ (2 * j) * spinQAnalytic SpinPair.pp (2 * j - 1) z := by
...
theorem spinQAnalytic_pp_one_S {z : Complex} (hz : z ∈ imGT 0) :
    spinQAnalytic SpinPair.pp 1 (-1 / z) =
      z ^ 2 * spinQAnalytic SpinPair.pp 1 z - z / spinTwoPiI := by
\end{leancode}
\begin{dictionary}
\lean{imGT 0} is $\{z:\operatorname{Im}z>0\}$. The $(+,+)$ cells are the Eisenstein $S$-laws of
Chapter~\ref{ch:L14} transported through the identification with \lean{zhuGC}; the weight-$2$ cell
carries the $G_2$ anomaly. \textbf{The three other spin sectors have no proved $S$-law}; the
$T^2$-law holds for all four. The weight-$1$ case $k=0$ (a constant) is deliberately not given a
transformation law.
\end{dictionary}
\begin{physics}
Only the untwisted Eisenstein series are covered; the twisted ones, which mix sectors under $S$
($(+,-)\leftrightarrow(-,+)$), are exactly what a $\Gamma_\theta$-modularity proof needs.
\end{physics}
\begin{proof}[Proof on paper]
$T^2$: every exponent $N/2$ satisfies $e^{2\pi i(N/2)\cdot2}=1$. $S$ at $(+,+)$: the odd half-exponents
vanish, the even ones reindex to $\sum_nG_{2j}(n)q^n$, and the Eisenstein law applies.
\end{proof}
\begin{status}
\texttt{\#print axioms} of the three theorems: \lean{[propext, Classical.choice, Quot.sound]}.
Unconditional, as statements about coefficient functions. Their intended use, traces over twisted (spin) sectors, is
vacuous in the library: those sectors are \lean{SuperTwistedModule}s with $N\ge2$, hence zero (F1).
\end{status}

\subsection{Genus-one blocks for integrally graded super VOAs}

\textbf{Notation: square-bracket modes.} Genus-one (torus) quantities are naturally written in the
cylinder coordinate $w$, with $z=e^{w}-1$. Zhu's change of coordinates
(Chapter~\ref{ch:L15}, Definition~\ref{l15:def:sqmode}; Chapter~\ref{ch:L16}, §\ref{l16:sec:formal})
defines, for $a$ of integral weight $k$, the field $Y[a,w]=Y(a,e^{w}-1)e^{kw}=\sum_{m\in\Z}a[m]w^{-m-1}$;
its coefficients $a[m]$ are the \emph{square-bracket modes}. Each $a[m]$ is a finite combination
$a[m]=a_{(m)}+(\text{terms }a_{(m+i)},\ i\ge1)$ of ordinary modes, so $a[-2]b$ differs from $a_{(-2)}b$
by terms of lower weight. The square-bracket Virasoro operators are
$L[n]=\omega[n+1]-\frac c{24}\delta_{n,-2}\,\mathrm{id}$, and $a$ is \emph{square-bracket primary} of weight $\Delta$
if $L[0]a=\Delta a$ and $L[n]a=0$ for $n\ge1$. The functions $G_{2m+2}(z)$ below are Zhu's normalised
Eisenstein series (Chapter~\ref{ch:L14}, Definition~\ref{l14:def:zhuGC}).

\begin{definition}[Analytic super block]\label{l17:def:superblock}
For an integrally graded super near-VOA $V$, a function $F:V\times\{\operatorname{Im}z>0\}\to\C$, linear
in the state and holomorphic in $z$, is an analytic super block if $F(a[0]v)=0$ for $a\in V_k$, $k\ge1$,
$F(a[-2]v)+\sum_{m\ge0}(2m+1)G_{2m+2}(z)F(a[2m]v)=0$, and the $\omega$-recursion
$F(L[-2]v)=\frac1{2\pi i}\partial_zF(v)+\sum_mG_{2m+2}(z)F(L[2m]v)$ holds.
\end{definition}
\noindent File \leanpath{VOA/\allowbreak Modularity/\allowbreak SuperAnalyticBlock.lean}.
\begin{leancode}
structure IsSuperAnalyticBlock (F : 𝒱 → ℂ → ℂ) : Prop where
  holo : ∀ v : 𝒱, DifferentiableOn ℂ (F v) (imGT 0)
  addv : ∀ (v w : 𝒱) (z : ℂ), F (v + w) z = F v z + F w z
  smulv : ∀ (c : ℂ) (v : 𝒱) (z : ℂ), F (c • v) z = c * F v z
  vanish : ∀ (k : ℕ), 1 ≤ k → ∀ a ∈ V.gradingZ (k : ℤ), ∀ v : 𝒱,
    F (V.superVSqZ a k 0 v) = 0
  negTwo : ∀ (k : ℕ), 1 ≤ k → ∀ a ∈ V.gradingZ (k : ℤ), ∀ (v : 𝒱), ∀ z ∈ imGT 0,
    F (V.superVSqZ a k (-2) v) z
      + ∑ᶠ m : ℕ, ((2 * (m + 1) - 1 : ℕ) : ℂ)
          * (zhuGC (m + 1) z
            * F (V.superVSqZ a k (2 * (m + 1) - 2 : ℤ) v) z) = 0
  omega : ∀ (v : 𝒱), ∀ z ∈ imGT 0,
    F (V.superLsq (-2) v) z
      = (2 * (Real.pi : ℂ) * Complex.I)⁻¹ * _root_.deriv (F v) z
        + ∑ᶠ m : ℕ, zhuGC (m + 1) z
            * F (V.superLsq (2 * (m + 1) - 2 : ℤ) v) z
\end{leancode}
\begin{dictionary}
\lean{V.superVSqZ a k m} is the square-bracket mode $a[m]$ for $a$ of weight $k$ (Zhu's change of
coordinates, Chapter~\ref{ch:L15}, Definition~\ref{l15:def:sqmode}); \lean{superLsq n} $=L[n]$;
\lean{zhuGC j} $=G_{2j}$ in Zhu's normalisation (Chapter~\ref{ch:L14}, Definition~\ref{l14:def:zhuGC}).
The $m=0$ term of \lean{negTwo} is $G_2F(a[0]v)$, killed by \lean{vanish}, so this is Zhu's recursion
$\sum_{k\ge2}$. The Koszul sign does not appear in the definition: it appears only when $a[0]v$ is
identified, modulo $O(V)$, with a supercommutator $[a]\ast[v]-(-1)^{pq}[v]\ast[a]$
(\leanpath{SuperZhuCommutator.lean}); a supertrace kills supercommutators.
\end{dictionary}
\begin{physics}
The torus one-point functions $\lseventeenstr_Mo(v)q^{L_0-c/24}$ (with the $(-1)^F$ insertion, i.e.\ the
$(+,-)$-type temporal structure for integral weights) satisfy these Ward identities; the definition
axiomatises them.
\end{physics}
\begin{status}
A definition; standard axioms.
\end{status}

\begin{theorem}[Super Arike--Nagatomo upper bound]\label{l17:thm:AN}
If $W$ is an integrally graded, super $C_2$-cofinite super VOA, the space of analytic super blocks
(restricted to the upper half-plane) is finite-dimensional and its dimension is at most the dimension
of the space of supersymmetric linear functionals on $A(W)$ (those $\phi$ with
$\phi(ab)=(-1)^{pq}\phi(ba)$ on homogeneous $a,b$).
\end{theorem}
\noindent File \leanpath{VOA/\allowbreak Modularity/\allowbreak SuperBlockBoundC2.lean}.
\begin{leancode}
theorem finrank_superAnalyticBlocks_le_finrank_superZhuSymmetricFunctionals
    (W : SuperVertexOperatorAlgebra 𝒱) (hV : W.toSuperNearVOA.IntegralWeights)
    (hC2 : SuperC2Cofinite W) :
    finrank ℂ (Submodule.map restrictImGT W.toSuperNearVOA.superAnalyticBlocks)
      ≤ (let hZ := W.superZhuRingData hV
        letI := hZ.AsRing
        letI := hZ.AsAlgebra
        finrank ℂ hZ.superZhuSymmetricFunctionals) := by
\end{leancode}
\begin{dictionary}
\lean{restrictImGT} restricts functions to the upper half-plane. \lean{finrank} of an
infinite-dimensional space is $0$ in Mathlib, so the companion
\lean{finiteDimensional_superAnalyticBlocks_of_superC2} (finite-dimensionality) is what prevents the
inequality from being vacuous; it is proved. No rationality or semisimplicity is assumed. The bound
is the upper-bound half of Arike--Nagatomo's theorem (super version); the realisation of the bound by
pseudo-supertraces is not proved.
\end{dictionary}
\begin{physics}
The number of independent torus amplitudes with a $(-1)^F$ insertion is at most the number of
supersymmetric functionals on the zero-mode algebra; in a logarithmic theory such as symplectic
fermions, pseudo-supertraces over non-semisimple sectors are expected to saturate it.
\end{physics}
\begin{proof}[Proof on paper]
Every block kills the relation space spanned by the $a[0]v$ and the $a[-2]v+\dots$; by $C_2$-cofiniteness
and a triangular descent (the square-bracket $a[-2]b$ differs from $a_{(-2)}b$ by lower-weight terms),
finitely many states $T$ span modulo the relations at some interior point, with $|T|$ bounded by the
cusp quotient, which is identified with $A(W)/[A(W),A(W)]_{\mathrm{super}}$
(\leanpath{SuperCuspFunctionalCount.lean}); a block vanishing at $T$ at one point vanishes identically
by the $\omega$-ODE. Its dual is the space of supersymmetric functionals.
\end{proof}
\begin{status}
\texttt{\#print axioms
VOA.SuperNearVOA.\allowbreak finrank\_superAnalyticBlocks\_\allowbreak le\_finrank\_\allowbreak superZhuSymmetricFunctionals}:
\lean{[propext, Classical.choice, Quot.sound]}. Unconditional beyond \lean{hV}, \lean{hC2}. Instance:
symplectic fermions (\leanpath{SympFermSuperBlockBound.lean}, dimension $\le4$).
\end{status}

\begin{theorem}[Supertraces are formal super blocks]\label{l17:thm:formalblock}
For every graded super module $N$ of an integrally graded super VOA, the formal one-point supertrace
$v\mapsto\sum_e\lseventeenstr_{N_{e+c/24}}o(v)\,q^e$ is a formal super block: it kills the relation
module and satisfies the $\omega$-recursion with the formal Eisenstein coefficients.
\end{theorem}
\noindent File \leanpath{VOA/\allowbreak Modularity/\allowbreak SuperTorusRecursionClose.lean}.
\begin{leancode}
theorem superTorusOnePoint_isSuperFormalBlock
    (N : SuperOrdVModule W M) (hW : W.toSuperNearVOA.IntegralWeights) :
    SuperNearVOA.IsSuperFormalBlock W (N.superTorusOnePoint hW) := by
\end{leancode}
\begin{dictionary}
\lean{IsSuperFormalBlock W S} (\leanpath{SuperFormalBlock.lean}) is the conjunction of
``\lean{superBlockRelSpanG} is in the kernel of every supported lift'' and the formal $\omega$-recursion
$S(L[-2]v)=q\frac{d}{dq}S(v)+\sum_m G_{2m+2}\ast S(L[2m]v)$ (\lean{qConv} is convolution with the
coefficient sequence). \lean{superTorusOnePoint} takes at exponent $e$ the supertrace (\lean{supertrace},
$\tr((-1)^F\cdot)$, basis-free) of the zero mode on $N_{e+c/24}$. Since the only \lean{SuperOrdVModule}s
constructed are bosonic (\lean{OrdVModule.toSuper}), this theorem has no instance with a nonzero odd
part in the library.
\end{dictionary}
\begin{physics}
Zhu's recursion for torus amplitudes with $(-1)^F$: the two signs (Koszul in the commutator, Koszul in
supertrace cyclicity) cancel, so the bosonic Lambert--Eisenstein resummation goes through.
\end{physics}
\begin{proof}[Proof on paper]
The super module commutator formula from the $\ell=0$ slice of the module Jacobi identity gives
$o(a[0]b)=o(a)o(b)-(-1)^{pq}o(b)o(a)$, which has zero supertrace by Koszul cyclicity; the $[-2]$ and
$\omega$ relations follow from Zhu's geometric-series argument, the compensated pair carrying one
Koszul scalar from the iterate and the same from cyclicity, whose product is $1$.
\end{proof}
\begin{status}
\texttt{\#print axioms VOA.SuperOrdVModule.superTorusOnePoint\_isSuperFormalBlock}: \lean{[propext,
Classical.choice, Quot.sound]}. Unconditional; untwisted sector only. Odd (type-Q) traces: only the
square-bracket-zero slice is proved.
\end{status}

\begin{theorem}[Super Arakawa--Kawasetsu formal ODE]\label{l17:thm:AK}
If $W$ is integrally graded and super Etingof-finite, $S$ is a formal super block supported on
$q^{h+\N}$, and $a$ is square-bracket primary, then $S(a)$ satisfies a monic formal $q$-differential
equation $(q\frac d{dq})^sS(a)+\sum_{i<s}g_i\ast(q\frac d{dq})^iS(a)=0$ with subexponential coefficient
sequences $g_i$.
\end{theorem}
\noindent File \leanpath{VOA/\allowbreak Modularity/\allowbreak SuperAKFormalODE.lean}.
\begin{leancode}
theorem formalODE_of_etingofFinite_of_sqPrimary {S : 𝒱 →ₗ[ℂ] FormalQSeries} {hbase : ℂ}
    (hsupp : ∀ v, S v ∈ QSeriesAt hbase) (hS : SuperNearVOA.IsSuperFormalBlock W S)
    (hV : W.toSuperNearVOA.IntegralWeights) (hE : SuperEtingofFinite W) {a : 𝒱} {Δ : ℂ}
    (ha : W.toSuperNearVOA.IsSuperSqPrimary a Δ) :
    ∃ (s : ℕ) (g : ℕ → ℕ → ℂ), 1 ≤ s ∧ (∀ i, Subexponential (fun k => ‖g i k‖)) ∧
      ∀ e : ℂ, (qDeriv^[s] (S a)) e
        + ∑ i ∈ Finset.range s, qConv (g i) (qDeriv^[i] (S a)) e = 0 :=
\end{leancode}
\begin{dictionary}
\lean{QSeriesAt hbase}: series supported on $hbase+\N$. \lean{IsSuperSqPrimary a Δ}: $L[0]a=\Delta a$,
$L[n]a=0$ for $n\ge1$. \lean{qDeriv} is $q\frac d{dq}$ on formal series; \lean{Subexponential}: growth
$e^{o(k)}$ (Chapter~\ref{ch:L14}, Definition~\ref{l14:def:subexp}). The coefficients are formal series,
not elements of $\C[G_4,G_6]$; so this is \textbf{not} a modular linear differential equation, and
nothing about modularity or convergence follows (the docstring says so).
\end{dictionary}
\begin{physics}
Arakawa--Kawasetsu: characters of quasi-lisse theories satisfy modular linear differential equations.
The formal shadow proved here is the Noetherian argument; the modular coefficients and the analytic
continuation are missing.
\end{physics}
\begin{proof}[Proof on paper]
By Etingof-finiteness the block span is a finitely generated module over the Noetherian ring of Zhu's
quasimodular coefficients, so the tower $S(L[-2]^ia)$ satisfies a monic relation; for a primary $a$ the
$\omega$-recursion makes $S(L[-2]^ia)$ unitriangular in $(q\frac d{dq})^jS(a)$, which converts the relation.
\end{proof}
\begin{status}
\texttt{\#print axioms VOA.SuperAK.formalODE\_of\_etingofFinite\_of\_sqPrimary}: \lean{[propext,
Classical.choice, Quot.sound]}. Conditional on \lean{hS} (supplied by Theorem~\ref{l17:thm:formalblock}
for supertraces), \lean{hV}, \lean{hE}, \lean{ha}. No instance in the library combines integral weights,
super Etingof-finiteness and a genuinely odd part.
\end{status}

\subsection{The half-integral side: finite cusp descent}

\begin{theorem}[Finite square-spin cusp quotient]\label{l17:thm:cusp}
Let $W$ be a half-integrally graded super VOA ($W_h=0$ unless $h\in\frac12\Z$), super $C_2$-cofinite, and
$(g,h)$ a commuting pair of involutive even automorphisms. Then the quotient of $W$ by the
square-spin cusp relation span of the $(+,+)$ inner eigensector of $(g,h)$ is finite-dimensional.
\end{theorem}
\noindent File \leanpath{VOA/\allowbreak Modularity/\allowbreak SuperSpinDescentSqFinite.lean}.
\begin{leancode}
theorem BoundaryPair.finiteDimensional_superSpinSqCuspQuotient_of_superC2
    (b : BoundaryPair W) (hhalf : IsHalfIntegerGraded W)
    (hspatial : b.spatial ^ 2 = 1) (htemporal : b.temporal ^ 2 = 1)
    (hC2 : SuperC2Cofinite W) :
    FiniteDimensional Complex
      (VType ⧸ superSpinSqCuspRelSpan W hhalf b.spinEigenspace) :=
\end{leancode}
\noindent File \leanpath{VOA/\allowbreak Modularity/\allowbreak SuperSectorSystem.lean}.
\begin{leancode}
structure BoundaryPair (W : SuperVertexOperatorAlgebra VType) where
  spatial : SuperAut W
  temporal : SuperAut W
  commute : Commute spatial temporal
\end{leancode}
\begin{dictionary}
\lean{IsHalfIntegerGraded W}: $W_k=0$ unless $k=m/2$ for some $m\in\Z$ (\leanpath{GammaThetaStatement.lean}).
\lean{b.spinEigenspace} assigns to each \lean{SpinPair} the simultaneous $\pm1$ eigenspace of the two
involutions; \lean{superSpinSqCuspRelSpan} is spanned by the cusp ($q\to0$) limits of the square-bracket
spin relations together with the three non-$(+,+)$ eigenspaces, the square $C_2$ being formed with the
square-bracket $[-2]$ products (a separate span \lean{superSqO₂}, because ordinary $C_2$ is not
square-graded). No spin--statistics and no identification of $g,h$ with parity. This does not involve
twisted modules, so (F1) does not touch it.
\end{dictionary}
\begin{physics}
The finiteness of the space of genus-one conformal blocks in one spin structure, at the level of the
cusp: the algebraic precondition of any $\Gamma_\theta$-modularity theorem.
\end{physics}
\begin{proof}[Proof on paper]
Ordinary $C_2$-cofiniteness implies square $C_2$-cofiniteness by unitriangular induction on doubled
weight (the square $[-2]$ mode has leading coefficient one); projecting to the $(+,+)$ eigenspace is
compatible with the square grading.
\end{proof}
\begin{status}
\texttt{\#print axioms
VOA.SuperNearVOA.\allowbreak BoundaryPair.\allowbreak finiteDimensional\_\allowbreak superSpinSqCuspQuotient\_\allowbreak of\_superC2}:
\lean{[propext, Classical.choice, Quot.sound]}. Unconditional beyond its hypotheses. No half-integrally
graded super $C_2$-cofinite instance exists in the library ($\beta\gamma$ is not $C_2$-cofinite; the free
fermion is not built).
\end{status}

%======================================================================
\section{The four-by-three grid: what each cell contains}\label{l17:sec:grid}
%======================================================================

The project organises super modularity as a grid. The \emph{columns} are four independent kinds of
object: integral non-super (ordinary VOAs), integral super (integer weights, nontrivial parity),
half-integral non-super (weights in $\frac12\Z$, purely even; represented as purely even
\lean{SuperVertexOperatorAlgebra}s, since \lean{VertexOperatorAlgebra} is $\Z$-graded), and
half-integral super. The \emph{rows} are three routes to modularity: (R1) strong rationality and Zhu's
theorem (Chapter~\ref{ch:L15}); (R2) strong finiteness and the Arike--Nagatomo / pseudo-trace theorem
(Chapter~\ref{ch:L18}); (R3) the quasi-lisse (Etingof-finite) Arakawa--Kawasetsu route
(Chapter~\ref{ch:L16}). A cell is \emph{complete} only when the corresponding modular-invariance
theorem is proved for the abstract class \emph{and} instantiated. The table below says exactly what is
in Lean. ``(F1)'' marks a component that is vacuous by Theorem~\ref{l17:thm:twistvacuous}: it takes
a \lean{SuperTwistedModule} with $N\ge2$ as input, and (when this chapter was first written, and still in the released source) every such module is
the zero space. Such components are listed
for completeness only; they have no mathematical content.

\begin{center}
\scriptsize
\renewcommand{\arraystretch}{1.3}
\begin{longtable}{@{}>{\raggedright\arraybackslash}p{1.4cm}>{\raggedright\arraybackslash}p{3.1cm}>{\raggedright\arraybackslash}p{3.1cm}>{\raggedright\arraybackslash}p{3.1cm}>{\raggedright\arraybackslash}p{3.1cm}@{}}
\toprule
& \textbf{Integral non-super} & \textbf{Integral super} & \textbf{Half-integral non-super} &
\textbf{Half-integral super}\\
\midrule
\endhead
\textbf{R1 Zhu} &
Bosonic $SL_2(\Z)$ modularity proved \emph{conditionally} on \lean{StronglyRational} plus named
Dong--Li--Mason inputs (Chapter~\ref{ch:L15}). No model instantiated (Lee--Yang lacks
\lean{StronglyRational}). &
Proved: super Zhu ring (Thm.~\ref{l17:thm:superzhuring}); identity-sector supertrace is a formal
block (Thm.~\ref{l17:thm:formalblock}); finite type-M/type-Q Wedderburn functional calculus for a
\emph{supplied} decomposition (\leanpath{SuperWedderburn}\allowbreak\leanpath{Trace.lean}, \leanpath{QueerOddTrace.lean}).
Supplied-data sector realization (\lean{FiniteCompleteSector}\allowbreak\lean{TraceRealization}): vacuous for $N\ge2$
(F1). No modularity theorem. No model. &
Proved: $\Gamma_\theta$, half-automorphy factor, spin coefficients with $T^2$-law and $(+,+)$ $S$-cells
(Thms.~\ref{l17:thm:gammatheta}--\ref{l17:thm:spinS}); finite cusp descent
(Thm.~\ref{l17:thm:cusp}). Denominator-two trace carrier: vacuous (F1). No $\Gamma_\theta$ theorem.
No model. &
Same foundation as the previous column (parity independent of weight). No $\Gamma_\theta$ theorem.
Intended model, the free fermion, exists only as operators (Thm.~\ref{l17:thm:psinotfield}).\\
\addlinespace
\textbf{R2 AN} &
Analytic-block upper bound and partial pseudo-trace blocks (Chapter~\ref{ch:L18}); realization,
convergence, spanning, modular invariance open. &
Proved: super AN upper bound (Thm.~\ref{l17:thm:AN}). Model: symplectic fermions are super
$C_2$-cofinite with non-semisimple $A(SF)$, $\dim\le4$ (Thms.~\ref{l17:thm:sfc2},
\ref{l17:thm:sfzhu}); bound $\le4$ for them. No pseudo-supertraces; irrationality only conditional
(Thm.~\ref{l17:thm:sfirrat}). &
Finite cusp descent (Thm.~\ref{l17:thm:cusp}); the recursion bridge
(\lean{SuperSpinSq}\allowbreak\lean{RecursionData}) has three unproved fields and acts on a vacuous carrier (F1). &
Same, parity independent. No model.\\
\addlinespace
\textbf{R3 AK} &
A common monic MLDE for analytic blocks from \lean{EtingofFinite} and supplied grading/Ramanujan
inputs (Chapter~\ref{ch:L16}); character modularity open. &
Proved: formal super AK ODE (Thm.~\ref{l17:thm:AK}), conditional on a formal block and
\lean{SuperEtingofFinite}; relative quotient carriers for twisted modules vacuous for $N\ge2$ (F1).
No convergence, no modularity. No model. &
Proved: Noetherian bounded coefficient algebra and a $pp$ cusp recurrence
(\leanpath{SuperSpinQuasiLisse}\allowbreak\leanpath{Algebra.lean}), conditional on \lean{SuperSpinCusp}\allowbreak\lean{EulerData} and
\lean{SuperSpinRelative}\allowbreak\lean{EtingofData}. Model: $\beta\gamma$ is super Etingof-finite and not super $C_2$
(Thm.~\ref{l17:thm:bgquasi}); no trace or ODE instantiated. &
Same conditional infrastructure, parity independent. No model.\\
\bottomrule
\end{longtable}
\end{center}

\begin{warning}[title={The grid in one sentence}]
\textbf{No cell of the three new columns contains a modular-invariance theorem}, conditional or not;
every proved item is a foundation, a finiteness or upper-bound result, a formal (non-analytic) ODE, or
a transformation law of the Eisenstein-type coefficient functions. In addition, finding (F1) empties
every component that uses a \lean{SuperTwistedModule} with $N\ge2$: the half-integral trace
carriers, the Ramond support statements, the twisted sector realization, and the relative quotients
of twisted modules. The integral non-super column is documented in Chapters~\ref{ch:L15},
\ref{ch:L16}, \ref{ch:L18}; its status was not re-verified for this chapter.
\end{warning}

%======================================================================
\section{Machine status of the headlines}\label{l17:sec:status}
%======================================================================

All axiom checks below were run for this chapter on a scratch file importing the relevant modules
(the umbrellas \leanpath{Modularity/\allowbreak SuperIntegral.lean} and
\leanpath{Modularity/\allowbreak HalfIntegralGridFoundations.lean} among them) with
\texttt{lake env lean -DmaxSynthPendingDepth=3}; no full \texttt{lake build} was run. Every
declaration listed reports exactly the three standard axioms; \textbf{none uses \lean{VOA.pbw} (the PBW theorem of Chapter~\ref{ch:L02}; formerly a project axiom, now a theorem)}.

\begin{plaincode}
'VOA.SuperVertexAlgebra' depends on axioms: [propext, Classical.choice, Quot.sound]
'VOA.SuperVertexAlgebra.dong_fields' depends on axioms: [propext, Classical.choice, Quot.sound]
'VOA.SuperVertexAlgebra.goddard_uniquenessε' depends on axioms: [propext, Classical.choice, Quot.sound]
'VOA.SuperNearVOA.superBorcherds' depends on axioms: [propext, Classical.choice, Quot.sound]
'VOA.SuperNearVOA.superJacobi' depends on axioms: [propext, Classical.choice, Quot.sound]
'VOA.SuperNearVOA.superSmul_eq_skew_Lpow' depends on axioms: [propext, Classical.choice, Quot.sound]
'VOA.SuperVertexOperatorAlgebra.skew_modeε' depends on axioms: [propext, Classical.choice, Quot.sound]
'VOA.SuperVertexOperatorAlgebra.orderOf_parityAut_eq_two' depends on axioms: [propext, Classical.choice, Quot.sound]
'VOA.SuperVertexOperatorAlgebra.evenPartVOA' depends on axioms: [propext, Classical.choice, Quot.sound]
'VOA.SuperNearVOA.adjointSuperWeakModule' depends on axioms: [propext, Classical.choice, Quot.sound]
'VOA.FreeFermion.psi_localε' depends on axioms: [propext, Classical.choice, Quot.sound]
'VOA.Ghost.bcGhostSVA' depends on axioms: [propext, Classical.choice, Quot.sound]
'VOA.BGInst.bgSVOA' depends on axioms: [propext, Classical.choice, Quot.sound]
'VOA.BGNonC2.not_SuperC2Cofinite_bgSVOA' depends on axioms: [propext, Classical.choice, Quot.sound]
'VOA.SympFerm.sympFermSVOA' depends on axioms: [propext, Classical.choice, Quot.sound]
'VOA.SympFerm.sympFerm_superC2Cofinite' depends on axioms: [propext, Classical.choice, Quot.sound]
'VOA.SympFerm.ASF_not_isSemisimpleRing' depends on axioms: [propext, Classical.choice, Quot.sound]
'VOA.SympFerm.sympFermSVOA_not_isSuperRational' depends on axioms: [propext, Classical.choice, Quot.sound]
'VOA.SuperVertexOperatorAlgebra.superZhuRingData' depends on axioms: [propext, Classical.choice, Quot.sound]
'VOA.finiteDimensional_twistedA_of_superC2Cofinite' depends on axioms: [propext, Classical.choice, Quot.sound]
'VOA.ThetaGroup.GammaTheta_eq_genSTsq' depends on axioms: [propext, Classical.choice, Quot.sound]
'VOA.ThetaGroup.index_eq_three' depends on axioms: [propext, Classical.choice, Quot.sound]
'VOA.spinThetaChar_S' depends on axioms: [propext, Classical.choice, Quot.sound]
'VOA.spinQAnalytic_pp_higher_S' depends on axioms: [propext, Classical.choice, Quot.sound]
'VOA.spinQAnalytic_pp_one_S' depends on axioms: [propext, Classical.choice, Quot.sound]
'VOA.SuperNearVOA.finrank_superAnalyticBlocks_le_finrank_superZhuSymmetricFunctionals' depends on axioms: [propext, Classical.choice, Quot.sound]
'VOA.SuperAK.formalODE_of_etingofFinite_of_sqPrimary' depends on axioms: [propext, Classical.choice, Quot.sound]
'VOA.SuperOrdVModule.superTorusOnePoint_isSuperFormalBlock' depends on axioms: [propext, Classical.choice, Quot.sound]
'VOA.SuperNearVOA.BoundaryPair.finiteDimensional_superSpinSqCuspQuotient_of_superC2' depends on axioms: [propext, Classical.choice, Quot.sound]
\end{plaincode}

The Lean checks behind Theorems~\ref{l17:thm:nocov}, \ref{l17:thm:twistvacuous} and~\ref{l17:thm:psinotfield}
(not part of this release) compile with the same three axioms.

\paragraph{Release audit 2026-10-27: RELEASE\_AUDIT\_L17\_L18\_L23.md.} The independent release audit
\texttt{Notes/\allowbreak Textbook/\allowbreak referee/\allowbreak RELEASE\_AUDIT\_L17\_L18\_L23.md} re-elaborated, at a later
pre-release commit and from the existing build products (\texttt{lake env lean -DmaxSynthPendingDepth=3}; \texttt{lake build}
not run), \lean{SuperNearVOA.superJacobi}, \lean{evenPartVOA}, \lean{sympFermSVOA}, \lean{sympFerm_superC2Cofinite},
\lean{ASF_not_isSemisimpleRing}, \lean{finrank_ASF_le_four}, \lean{sympFermSVOA_not_isSuperRational}, \lean{bgSVOA},
\lean{not_SuperC2Cofinite_bgSVOA}, \lean{bgSVOA_superEtingofFinite} and \lean{bcGhostSVA}; each reports only the three standard
axioms. It recompiled the checks mentioned above. Verdicts: \emph{super VOAs, even part, symplectic fermions, $\beta\gamma$:
announceable with caveats} --- the super Jacobi identity for \lean{SuperNearVOA}, \lean{evenPartVOA} and \lean{sympFermSVOA}
($c=-2$, genuinely odd) are sound; super $C_2$-cofiniteness and non-semisimplicity of the super Zhu algebra are proved, its
dimension only shown $\le4$ (m11); \lean{bgSVOA} ($c=-1$) is purely even (m10); \lean{SuperVertexAlgebra} lacks translation
covariance, which weakens \lean{bcGhostSVA} (m9). \emph{Twisted/Ramond super modules, super-irrationality and the free fermion:
not announceable} (M6, m10, m12): \lean{SuperTwistedModule} is the zero module for $N\ge2$, still at the audited commit and
then unflagged in the library (the released source now flags it in docstrings) (warning box in Section~\ref{l17:warn:twistzero}); \lean{NonSplitSuperModule sympFermSVOA} is never
constructed; the free fermion is not a field. No modular-invariance theorem exists in any of the three new grid columns (m13).
The audit also records (M7) that the aggregate build is not certified for the release commit.

\begin{warning}[title={Further minor finding: irreducible twisted modules may be zero}]
\lean{SuperTwistedModule.IsIrreducible} (\leanpath{Modularity/\allowbreak SuperSectorRationality.lean})
requires only that every graded invariant subspace be $\bot$ or $\top$; it does not require $M\neq0$,
so the zero module is ``irreducible''. Combined with (F1), for $N\ge2$ the irreducible objects of
\lean{PackagedIrreducibleSuperTwistedModule} are exactly zero modules. Severity: moderate (compare \lean{IsLowestWeight}, which likewise does not require $M_h\neq0$;
Chapter~\ref{ch:L12}).
\end{warning}

%======================================================================
\section{Referee's checklist}\label{l17:sec:checklist}
%======================================================================

\begin{center}
\scriptsize
\renewcommand{\arraystretch}{1.25}
\begin{longtable}{@{}>{\raggedright\arraybackslash}p{4.1cm}>{\raggedright\arraybackslash}p{4.3cm}>{\raggedright\arraybackslash}p{3.6cm}>{\raggedright\arraybackslash}p{2.4cm}@{}}
\toprule
\textbf{Lean declaration} & \textbf{Intended meaning} & \textbf{Hypotheses; status} & \textbf{Fidelity}\\
\midrule
\endhead
\lean{SuperVertexAlgebra} & super vertex algebra with $\C$-grading & definition; lacks translation
covariance (F2) & Def.~\ref{l17:def:sva}, Thm.~\ref{l17:thm:nocov}\\
\lean{SpinStatistics} & integer weight iff even & predicate; satisfiable & Def.~\ref{l17:def:spinstat}\\
\lean{SuperNearVOA}, \lean{SuperVertexOperatorAlgebra} & conformal super vertex algebra (of CFT type) &
definitions & Defs.~\ref{l17:def:snvoa}, \ref{l17:def:svoa}\\
\lean{dong_fields} & OPE of local fields is local, with product sign & homogeneity; unconditional &
Thm.~\ref{l17:thm:dong}\\
\lean{goddard_uniquenessε} & a field is determined by the state it creates & locality, vacuum tower;
unconditional & Thm.~\ref{l17:thm:goddard}\\
\lean{superBorcherds}, \lean{superJacobi} & full super Borcherds identity, all $\ell,m,n$ &
\lean{SuperNearVOA}, homogeneity; unconditional & Thm.~\ref{l17:thm:superjacobi}\\
\lean{superSmul_eq_skew_Lpow} & super skew-symmetry & \lean{SuperNearVOA}; unconditional &
Thm.~\ref{l17:thm:skew}\\
\lean{skew_modeε} & super skew-symmetry & gate \lean{WeakAssocLocε} never proved; superseded (F5) &
Prop.~\ref{l17:prop:skewgated}\\
\lean{parityAut}, \lean{orderOf_parityAut_eq_two} & $(-1)^F$ is an automorphism of order $2$ &
odd part nonzero; unconditional & Thm.~\ref{l17:thm:parityaut}\\
\lean{evenPartVOA} & even part is an ordinary VOA & integrality of even weights &
Thm.~\ref{l17:thm:evenpart}\\
\lean{SuperWeakModule}, \lean{adjointSuperWeakModule} & super module; adjoint module & definition;
unconditional & Def.~\ref{l17:def:superweak}, Prop.~\ref{l17:prop:adjoint}\\
\lean{SuperTwistedModule} & $g$-twisted super module & \textbf{uninhabited except by $M=0$ for $N\ge2$}
(F1; still so in the released source); content only at $N=1$; every twisted/Ramond claim built on it vacuous &
Def.~\ref{l17:def:supertwisted}, Thm.~\ref{l17:thm:twistvacuous}\\
\lean{FreeFermion.psi}, \lean{psi_localε} & NS free fermion & CAR, locality true; not a field, no SVA
(F3) & Def.~\ref{l17:def:psi}, Thm.~\ref{l17:thm:psinotfield}\\
\lean{Ghost.bcGhostSVA} & $bc$ system at $\lambda=1$ & unconditional in the weak sense (no translation covariance, F2); no Virasoro (F6) &
Thm.~\ref{l17:thm:bc}, Def.~\ref{l17:def:ghostvir}\\
\lean{BGInst.bgSVOA} & $\beta\gamma$ at $\lambda=\frac12$, $c=-1$, purely even & unconditional; never exercises the Koszul sign &
Thm.~\ref{l17:thm:bgsvoa}\\
\lean{not_SuperC2Cofinite_bgSVOA}, \lean{bgSVOA_strict_}\allowbreak\lean{halfIntegral_quasiLisse} & $\beta\gamma$
quasi-lisse, not lisse & unconditional & Thm.~\ref{l17:thm:bgquasi}\\
\lean{SympFerm.sympFermSVOA} & symplectic fermions, $c=-2$, genuinely odd states & unconditional & Thm.~\ref{l17:thm:sfsvoa}\\
\lean{sympFerm_superC2Cofinite} & SF are $C_2$-cofinite & unconditional & Thm.~\ref{l17:thm:sfc2}\\
\lean{ASF_not_isSemisimpleRing}, \lean{finrank_ASF_le_four} & $A(SF)$ not semisimple, $\dim\le4$ (not ``$=4$'') & unconditional & Thm.~\ref{l17:thm:sfzhu}\\
\lean{IsSuperRational} & complete reducibility of graded super modules & definition; quantifies over
\lean{SympFermSuperOrdVModule}, which omits \lean{weight_mode} (a larger class than
\lean{SuperOrdVModule}; F8) & Def.~\ref{l17:def:superrat}\\
\lean{sympFermSVOA_not_}\allowbreak\lean{isSuperRational} & SF not super-rational & \textbf{conditional on the missing
module} (F4); a socle lemma, not an irrationality theorem & Thm.~\ref{l17:thm:sfirrat}\\
\lean{superZhuRingData} & $A(W)$ associative (Dong--Zhao 3.1) & integral weights; unconditional &
Thm.~\ref{l17:thm:superzhuring}\\
\lean{finiteDimensional_}\allowbreak\lean{twistedA_of_}\allowbreak\lean{superC2Cofinite} & twisted Zhu finiteness & integral, $C_2$,
diagonalisability data & Thm.~\ref{l17:thm:twistedzhufinite}\\
\lean{GammaTheta_eq_genSTsq}, \lean{index_eq_three} & $\Gamma_\theta=\langle S,T^2\rangle$, index $3$
& unconditional & Thm.~\ref{l17:thm:gammatheta}\\
\lean{thetaHalfAutFactorExists} & half-automorphy factor on $\Gamma_\theta$ & unconditional &
Thm.~\ref{l17:thm:halfaut}\\
\lean{spinThetaChar_S} & theta-characteristic $S$-law & $\tau\neq0$; unconditional &
Thm.~\ref{l17:thm:thetaS}\\
\lean{spinQAnalytic_pp_higher_S}, \lean{_pp_one_S}, \lean{_add_two} & $(+,+)$ $S$-cells, $T^2$ for all &
unconditional; other three $S$-cells open & Thm.~\ref{l17:thm:spinS}\\
\lean{finrank_superAnalytic}\allowbreak\lean{Blocks_le_...} & super AN upper bound & integral, $C_2$ &
Thm.~\ref{l17:thm:AN}\\
\lean{superTorusOnePoint_}\allowbreak\lean{isSuperFormalBlock} & supertraces are formal blocks & integral; untwisted &
Thm.~\ref{l17:thm:formalblock}\\
\lean{formalODE_of_}\allowbreak\lean{etingofFinite_of_sqPrimary} & formal AK ODE & formal block, Etingof, primary; not
modular & Thm.~\ref{l17:thm:AK}\\
\lean{finiteDimensional_}\allowbreak\lean{superSpinSqCusp}\allowbreak\lean{Quotient_of_superC2} & finite cusp quotient, $(+,+)$ sector &
half-integral, $C_2$, involutive pair & Thm.~\ref{l17:thm:cusp}\\
\bottomrule
\end{longtable}
\end{center}

%======================================================================
\section{Notes and references}\label{l17:sec:notes}
%======================================================================

\textbf{Vertex superalgebras.} V.~G.~Kac, \emph{Vertex Algebras for Beginners}, University Lecture
Series 10, American Mathematical Society, second edition 1998: Definition~1.3 and §4.1 (vertex
superalgebras, the Koszul sign), Theorem~4.4 (Goddard uniqueness), §4.8 (the Borcherds identity
(4.8.3) and skew-symmetry), §3.6 (free fermions, $bc$ and $\beta\gamma$ systems). E.~Frenkel and
D.~Ben-Zvi, \emph{Vertex Algebras and Algebraic Curves}, AMS, 2nd ed.\ 2004. H.~Li, ``Local systems of
vertex operators, vertex superalgebras and modules'', J.\ Pure Appl.\ Algebra 109 (1996) 143--195.

\textbf{Twisted modules and Zhu algebras for vertex operator superalgebras.} C.~Dong and Z.~Zhao,
``Modularity in orbifold theory for vertex operator superalgebras'', Comm.\ Math.\ Phys.\ 260 (2005)
227--256, arXiv:math/0411523; and ``Twisted representations of vertex operator superalgebras'', Commun.\
Contemp.\ Math.\ 8 (2006) 101--121, arXiv:math/0411524 (the twisted Zhu algebra $A_g(V)$ and
Theorem~3.1). V.~G.~Kac and W.~Wang, ``Vertex operator superalgebras and their representations'',
Contemp.\ Math.\ 175 (1994). C.~Dong, H.~Li and G.~Mason, ``Modular-invariance of trace functions in
orbifold theory and generalized moonshine'', Comm.\ Math.\ Phys.\ 214 (2000) 1--56.

\textbf{Super modularity.} J.~van Ekeren, ``Modular invariance for twisted modules over a vertex operator
superalgebra'', Comm.\ Math.\ Phys.\ 322 (2013) 333--371, arXiv:1111.0682 (Section~3, equation (3.1)
and Lemma~3.1 for the coefficients $P^{\mu,\lambda}_k$; type-M and type-Q supertraces). G.~Höhn,
\emph{Selbstduale Vertexoperatorsuperalgebren und das Babymonster}, PhD thesis, Bonn 1995. Y.~Arike and
K.~Nagatomo, ``Some remarks on pseudo-trace functions for orbifold models associated with symplectic
fermions'', Int.\ J.\ Math.\ 24 (2013), arXiv:1104.0068. T.~Arakawa and K.~Kawasetsu, ``Quasi-lisse
vertex algebras and modular linear differential equations'', in \emph{Lie Groups, Geometry, and
Representation Theory}, Progr.\ Math.\ 326 (2018), arXiv:1610.05865. P.~Etingof and T.~Schedler,
``Poisson traces and D-modules on Poisson varieties'', GAFA 20 (2010).

\textbf{Symplectic fermions and logarithmic CFT.} H.~G.~Kausch, ``Symplectic fermions'', Nucl.\ Phys.\
B583 (2000) 513--541, arXiv:hep-th/0003029. M.~R.~Gaberdiel and H.~G.~Kausch, ``A rational logarithmic
conformal field theory'', Phys.\ Lett.\ B386 (1996) 131--137, arXiv:hep-th/9606050, and ``A local
logarithmic conformal field theory'', Nucl.\ Phys.\ B538 (1999) 631--658, arXiv:hep-th/9807091.
T.~Abe, ``A $\Z_2$-orbifold model of the symplectic fermionic vertex operator superalgebra'', Math.\
Z.\ 255 (2007) 755--792, arXiv:math/0503472.

\textbf{Ghosts and strings.} D.~Friedan, E.~Martinec and S.~Shenker, ``Conformal invariance,
supersymmetry and string theory'', Nucl.\ Phys.\ B271 (1986) 93--165 ($bc$ and $\beta\gamma$ systems,
$c=\mp(3(2\lambda-1)^2-1)$, the ghost anomalies $-26$ and $+11$). J.~Polchinski, \emph{String Theory}
vol.~1, §2.5.

\textbf{Theta functions.} D.~Mumford, \emph{Tata Lectures on Theta I}, Birkhäuser 1983; M.~Knopp,
\emph{Modular Functions in Analytic Number Theory}, Markham 1970 (the theta group).

\textbf{Within the library.} Project notes written before \leanpath{BetaGammaInstance.lean} and
\leanpath{GhostSystem.lean} describe the $bc$ and $\beta\gamma$ systems as ``not yet'' super VOA instances; that is out of
date for them (it remains true of the free fermion). The bosonic counterparts are Chapters~\ref{ch:L03} (axioms, Borcherds), \ref{ch:L12} (modules,
Zhu), \ref{ch:06} (twisted modules and the repaired \lean{TwistedModule.jacobi}), \ref{ch:L14} and
\ref{ch:L15} (modular forms and Zhu's theorem).

\endgroup
